% !TEX program = lualatex
% ---------------------------------------------------------------------------------------------
%  Readback note 25: Columns on a definite unitary group: divisibility and size of the chain of theta lifts
%
%  Build:  python3 tools/rb.py note <slug>
%
%  One unit of a Lean artifact, rewritten as ordinary mathematics. The prose is hand-written. Lean names,
%  line ranges and counts come from data.tex; quoted Lean statements are copied from the artifact by the build.
%
%  Rules of the text: standard terminology; every symbol is explained where it first appears and listed again
%  in the notation at the end; no letter has two meanings; a proof is complete or says what it takes for granted;
%  Lean names appear only in the LEAN lines and in the concordance.
%
%  The signature rule (the build enforces it): a function, relation or predicate that the artifact defines is
%  written only through its notation macro, which the build makes from the spec (`notation`) with one slot for
%  each explicit argument of the Lean definition. No argument is fixed and dropped, also not "for brevity" and
%  not for one section. Every argument and every hypothesis of a stated declaration stands inside its
%  statement, or in a Setting that the statement cites by number; only implicit and instance arguments (the
%  ambient objects) may stand in the conventions.
% ---------------------------------------------------------------------------------------------
\documentclass[11pt,a4paper]{article}
\ifdefined\pdfvariable\pdfvariable suppressoptionalinfo \numexpr 1+2+4+8\relax\fi
\usepackage[a4paper,hmargin=35mm,top=22mm,bottom=27mm,footskip=12mm,nohead]{geometry}

\usepackage{fontspec}
\directlua{
  luaotfload.add_fallback("gfallback", {"STIXTwoMath-Regular.otf:mode=node;", "DejaVuSans.ttf:mode=node;"})
  luaotfload.add_fallback("gmonofallback", {"DejaVuSansMono.ttf:mode=node;", "STIXTwoMath-Regular.otf:mode=node;"})
}
\setmainfont{NewCM10}[Extension=.otf,UprightFont=*-Book,ItalicFont=*-BookItalic,
  BoldFont=*-Bold,BoldItalicFont=*-BoldItalic,SmallCapsFont=*-Book,SmallCapsFeatures={Letters=SmallCaps},
  RawFeature={fallback=gfallback}]
\setsansfont{NewCMSans10}[Extension=.otf,UprightFont=*-Book,ItalicFont=*-BookOblique,
  BoldFont=*-Bold,BoldItalicFont=*-BoldOblique,RawFeature={fallback=gfallback}]
\setmonofont{JetBrainsMono}[Extension=.otf,UprightFont=*-Regular,ItalicFont=*-Italic,
  BoldFont=*-SemiBold,BoldItalicFont=*-SemiBoldItalic,Scale=MatchLowercase,
  RawFeature={-calt;-liga;fallback=gmonofallback}]
\usepackage{amsmath}
\usepackage{amsthm}
\usepackage[math-style=TeX]{unicode-math}
\setmathfont{NewCMMath-Book.otf}
\emergencystretch=2.5em\relax

\usepackage{xcolor}
\usepackage{tikz}
\usetikzlibrary{calc,arrows.meta,positioning,backgrounds}
\usepackage[protrusion=true,expansion=true]{microtype}
\usepackage{enumitem}
\usepackage{needspace}
\usepackage{booktabs}
\usepackage{array}
\usepackage{longtable}
\usepackage{fancyhdr}
\usepackage{lastpage}
\usepackage{multicol}
\usepackage[hidelinks,pdfusetitle]{hyperref}

\definecolor{ink}{HTML}{0F172A}
\definecolor{inksoft}{HTML}{475569}
\definecolor{inkfaint}{HTML}{94A3B8}
\definecolor{rulegrey}{HTML}{CBD5E1}
\definecolor{arith}{HTML}{C2410C}
\definecolor{est}{HTML}{7C3AED}
\definecolor{ana}{HTML}{0891B2}
\definecolor{ok}{HTML}{047857}
\definecolor{leancol}{HTML}{3730A3}

\title{Columns on a definite unitary group: divisibility and size of the chain of theta lifts}
\author{grh-compact}
\date{}

% ------------------------------------------------------------------------------------- macros
\newcommand{\leanfont}{\ttfamily\fontsize{8.6}{10.5}\selectfont\color{leancol}}
\newcommand{\lean}[1]{\csname LN@#1\endcsname}
\newcommand{\fillin}[1]{{\color{arith}\itshape[#1]}}  % a place of the template that is not written yet; the build reports every one that is left
\newcommand{\leanquote}[1]{\csname LQ@#1\endcsname}
% generated by rb.py note; do not edit
\newcommand{\ConcDecls}{105}
\newcommand{\Pin}{72fead10}
\newcommand{\PinText}{that of commit {\leanfont 72fead10}}
\newcommand{\ArtifactName}{grh-compact}
\expandafter\newcommand\csname LN@detKernel\endcsname{\href{https://github.com/lsernau/grh-lean/blob/72fead102f663d5c7a407964a13dc0cf9d977207/grh-compact/src/New/Columns/ChainValuation.lean\#L24-L25}{{\leanfont detKernel}}}
\expandafter\newcommand\csname LN@detPeriod\endcsname{\href{https://github.com/lsernau/grh-lean/blob/72fead102f663d5c7a407964a13dc0cf9d977207/grh-compact/src/New/Columns/ChainValuation.lean\#L29-L29}{{\leanfont detPeriod}}}
\expandafter\newcommand\csname LN@detKernelembFinmul\endcsname{\href{https://github.com/lsernau/grh-lean/blob/72fead102f663d5c7a407964a13dc0cf9d977207/grh-compact/src/New/Columns/ChainValuation.lean\#L35-L39}{{\leanfont detKernel\_embFin\_mul}}}
\expandafter\newcommand\csname LN@detPeriodeqstdClassSumsum\endcsname{\href{https://github.com/lsernau/grh-lean/blob/72fead102f663d5c7a407964a13dc0cf9d977207/grh-compact/src/New/Columns/ChainValuation.lean\#L44-L49}{{\leanfont detPeriod\_eq\_stdClassSum\_sum}}}
\expandafter\newcommand\csname LN@detPeriodinsert\endcsname{\href{https://github.com/lsernau/grh-lean/blob/72fead102f663d5c7a407964a13dc0cf9d977207/grh-compact/src/New/Columns/ChainValuation.lean\#L55-L83}{{\leanfont detPeriod\_insert}}}
\expandafter\newcommand\csname LN@IsShort\endcsname{\href{https://github.com/lsernau/grh-lean/blob/72fead102f663d5c7a407964a13dc0cf9d977207/grh-compact/src/New/Columns/ChainValuation.lean\#L91-L92}{{\leanfont IsShort}}}
\expandafter\newcommand\csname LN@subtranslate\endcsname{\href{https://github.com/lsernau/grh-lean/blob/72fead102f663d5c7a407964a13dc0cf9d977207/grh-compact/src/New/Columns/ChainValuation.lean\#L94-L99}{{\leanfont IsShort.\allowbreak{}sub\_translate}}}
\expandafter\newcommand\csname LN@spanisShorteqtop\endcsname{\href{https://github.com/lsernau/grh-lean/blob/72fead102f663d5c7a407964a13dc0cf9d977207/grh-compact/src/New/Columns/ChainValuation.lean\#L105-L147}{{\leanfont span\_isShort\_eq\_top}}}
\expandafter\newcommand\csname LN@algFormCentralRep\endcsname{\href{https://github.com/lsernau/grh-lean/blob/72fead102f663d5c7a407964a13dc0cf9d977207/grh-compact/src/New/Columns/ChainValuation.lean\#L151-L153}{{\leanfont algFormCentralRep}}}
\expandafter\newcommand\csname LN@notisChosenofNmem\endcsname{\href{https://github.com/lsernau/grh-lean/blob/72fead102f663d5c7a407964a13dc0cf9d977207/grh-compact/src/New/Columns/ChainValuation.lean\#L157-L159}{{\leanfont Data.\allowbreak{}not\_isChosen\_of\_N\_mem}}}
\expandafter\newcommand\csname LN@memcolumnLevel\endcsname{\href{https://github.com/lsernau/grh-lean/blob/72fead102f663d5c7a407964a13dc0cf9d977207/grh-compact/src/New/Columns/ChainValuation.lean\#L161-L164}{{\leanfont Data.\allowbreak{}mem\_columnLevel}}}
\expandafter\newcommand\csname LN@memchosenLevelofsupportedOn\endcsname{\href{https://github.com/lsernau/grh-lean/blob/72fead102f663d5c7a407964a13dc0cf9d977207/grh-compact/src/New/Columns/ChainValuation.lean\#L166-L168}{{\leanfont Data.\allowbreak{}mem\_chosenLevel\_of\_supportedOn}}}
\expandafter\newcommand\csname LN@existscentralGen\endcsname{\href{https://github.com/lsernau/grh-lean/blob/72fead102f663d5c7a407964a13dc0cf9d977207/grh-compact/src/New/Columns/ChainValuation.lean\#L173-L181}{{\leanfont Data.\allowbreak{}exists\_centralGen}}}
\expandafter\newcommand\csname LN@centralGen\endcsname{\href{https://github.com/lsernau/grh-lean/blob/72fead102f663d5c7a407964a13dc0cf9d977207/grh-compact/src/New/Columns/ChainValuation.lean\#L185-L185}{{\leanfont Data.\allowbreak{}centralGen}}}
\expandafter\newcommand\csname LN@scalarFincentralGenmem\endcsname{\href{https://github.com/lsernau/grh-lean/blob/72fead102f663d5c7a407964a13dc0cf9d977207/grh-compact/src/New/Columns/ChainValuation.lean\#L187-L189}{{\leanfont Data.\allowbreak{}scalarFin\_centralGen\_mem}}}
\expandafter\newcommand\csname LN@centralRep\endcsname{\href{https://github.com/lsernau/grh-lean/blob/72fead102f663d5c7a407964a13dc0cf9d977207/grh-compact/src/New/Columns/ChainValuation.lean\#L193-L193}{{\leanfont Data.\allowbreak{}centralRep}}}
\expandafter\newcommand\csname LN@ramifyingOp\endcsname{\href{https://github.com/lsernau/grh-lean/blob/72fead102f663d5c7a407964a13dc0cf9d977207/grh-compact/src/New/Columns/ChainValuation.lean\#L197-L198}{{\leanfont Data.\allowbreak{}ramifyingOp}}}
\expandafter\newcommand\csname LN@ramifiedColumns\endcsname{\href{https://github.com/lsernau/grh-lean/blob/72fead102f663d5c7a407964a13dc0cf9d977207/grh-compact/src/New/Columns/ChainValuation.lean\#L202-L202}{{\leanfont Data.\allowbreak{}ramifiedColumns}}}
\expandafter\newcommand\csname LN@existseqsubcentralGen\endcsname{\href{https://github.com/lsernau/grh-lean/blob/72fead102f663d5c7a407964a13dc0cf9d977207/grh-compact/src/New/Columns/ChainValuation.lean\#L207-L209}{{\leanfont Data.\allowbreak{}exists\_eq\_sub\_centralGen}}}
\expandafter\newcommand\csname LN@sumheightmulkappale\endcsname{\href{https://github.com/lsernau/grh-lean/blob/72fead102f663d5c7a407964a13dc0cf9d977207/grh-compact/src/New/Columns/ChainValuation.lean\#L211-L212}{{\leanfont Data.\allowbreak{}sum\_height\_mul\_kappa\_le}}}
\expandafter\newcommand\csname LN@finrankkeronesubcentralGenle\endcsname{\href{https://github.com/lsernau/grh-lean/blob/72fead102f663d5c7a407964a13dc0cf9d977207/grh-compact/src/New/Columns/ChainValuation.lean\#L214-L232}{{\leanfont Data.\allowbreak{}finrank\_ker\_one\_sub\_centralGen\_le}}}
\expandafter\newcommand\csname LN@ramifyingOpfinrankcomapandshort\endcsname{\href{https://github.com/lsernau/grh-lean/blob/72fead102f663d5c7a407964a13dc0cf9d977207/grh-compact/src/New/Columns/ChainValuation.lean\#L236-L246}{{\leanfont Data.\allowbreak{}ramifyingOp\_finrank\_comap\_and\_short}}}
\expandafter\newcommand\csname LN@spanshortramifiedColumns\endcsname{\href{https://github.com/lsernau/grh-lean/blob/72fead102f663d5c7a407964a13dc0cf9d977207/grh-compact/src/New/Columns/ChainValuation.lean\#L251-L257}{{\leanfont Data.\allowbreak{}span\_short\_ramifiedColumns}}}
\expandafter\newcommand\csname LQ@spanshortramifiedColumns\endcsname{\leanquoteline{0}{lemma span\_short\_ramifiedColumns : span ℂ \{φ : D.ramifiedColumns | IsShort D.budget.exp φ.1.1\} =}\leanquoteline{2}{⊤}}
\expandafter\newcommand\csname LN@pairTranslates\endcsname{\href{https://github.com/lsernau/grh-lean/blob/72fead102f663d5c7a407964a13dc0cf9d977207/grh-compact/src/New/Columns/ChainValuation.lean\#L264-L265}{{\leanfont Data.\allowbreak{}pairTranslates}}}
\expandafter\newcommand\csname LN@admissibleTranslations\endcsname{\href{https://github.com/lsernau/grh-lean/blob/72fead102f663d5c7a407964a13dc0cf9d977207/grh-compact/src/New/Columns/ChainValuation.lean\#L269-L269}{{\leanfont Data.\allowbreak{}admissibleTranslations}}}
\expandafter\newcommand\csname LN@IsGoodDatum\endcsname{\href{https://github.com/lsernau/grh-lean/blob/72fead102f663d5c7a407964a13dc0cf9d977207/grh-compact/src/New/Columns/ChainValuation.lean\#L275-L279}{{\leanfont Data.\allowbreak{}IsGoodDatum}}}
\expandafter\newcommand\csname LN@notMempiPlacesofne\endcsname{\href{https://github.com/lsernau/grh-lean/blob/72fead102f663d5c7a407964a13dc0cf9d977207/grh-compact/src/New/Columns/ChainValuation.lean\#L281-L283}{{\leanfont Data.\allowbreak{}notMem\_piPlaces\_of\_ne}}}
\expandafter\newcommand\csname LN@memintegralAfter\endcsname{\href{https://github.com/lsernau/grh-lean/blob/72fead102f663d5c7a407964a13dc0cf9d977207/grh-compact/src/New/Columns/ChainValuation.lean\#L288-L292}{{\leanfont Data.\allowbreak{}mem\_integralAfter}}}
\expandafter\newcommand\csname LN@projOntranslateDatum\endcsname{\href{https://github.com/lsernau/grh-lean/blob/72fead102f663d5c7a407964a13dc0cf9d977207/grh-compact/src/New/Columns/ChainValuation.lean\#L294-L297}{{\leanfont Data.\allowbreak{}projOn\_translateDatum}}}
\expandafter\newcommand\csname LN@existschaineqsum\endcsname{\href{https://github.com/lsernau/grh-lean/blob/72fead102f663d5c7a407964a13dc0cf9d977207/grh-compact/src/New/Columns/ChainValuation.lean\#L302-L312}{{\leanfont Data.\allowbreak{}exists\_chain\_eq\_sum}}}
\expandafter\newcommand\csname LN@translateDatummemstandardData\endcsname{\href{https://github.com/lsernau/grh-lean/blob/72fead102f663d5c7a407964a13dc0cf9d977207/grh-compact/src/New/Columns/ChainValuation.lean\#L314-L319}{{\leanfont Data.\allowbreak{}IsGoodDatum.\allowbreak{}translateDatum\_mem\_standardData}}}
\expandafter\newcommand\csname LN@gaugeLechain\endcsname{\href{https://github.com/lsernau/grh-lean/blob/72fead102f663d5c7a407964a13dc0cf9d977207/grh-compact/src/New/Columns/ChainValuation.lean\#L324-L335}{{\leanfont Data.\allowbreak{}IsGoodDatum.\allowbreak{}gaugeLe\_chain}}}
\expandafter\newcommand\csname LN@normstarpiGleone\endcsname{\href{https://github.com/lsernau/grh-lean/blob/72fead102f663d5c7a407964a13dc0cf9d977207/grh-compact/src/New/Columns/ChainValuation.lean\#L337-L338}{{\leanfont Data.\allowbreak{}norm\_star\_piG\_le\_one}}}
\expandafter\newcommand\csname LN@nezero\endcsname{\href{https://github.com/lsernau/grh-lean/blob/72fead102f663d5c7a407964a13dc0cf9d977207/grh-compact/src/New/Columns/ColumnDefs.lean\#L22-L22}{{\leanfont Data.\allowbreak{}π\_ne\_zero}}}
\expandafter\newcommand\csname LN@piPlaces\endcsname{\href{https://github.com/lsernau/grh-lean/blob/72fead102f663d5c7a407964a13dc0cf9d977207/grh-compact/src/New/Columns/ColumnDefs.lean\#L25-L26}{{\leanfont Data.\allowbreak{}piPlaces}}}
\expandafter\newcommand\csname LN@mempiPlaces\endcsname{\href{https://github.com/lsernau/grh-lean/blob/72fead102f663d5c7a407964a13dc0cf9d977207/grh-compact/src/New/Columns/ColumnDefs.lean\#L28-L28}{{\leanfont Data.\allowbreak{}mem\_piPlaces}}}
\expandafter\newcommand\csname LN@isSplitLevelpiPlaces\endcsname{\href{https://github.com/lsernau/grh-lean/blob/72fead102f663d5c7a407964a13dc0cf9d977207/grh-compact/src/New/Columns/ColumnDefs.lean\#L30-L32}{{\leanfont Data.\allowbreak{}isSplitLevel\_piPlaces}}}
\expandafter\newcommand\csname LN@cardpiPlaces\endcsname{\href{https://github.com/lsernau/grh-lean/blob/72fead102f663d5c7a407964a13dc0cf9d977207/grh-compact/src/New/Columns/ColumnDefs.lean\#L35-L38}{{\leanfont Data.\allowbreak{}card\_piPlaces}}}
\expandafter\newcommand\csname LN@columnLevel\endcsname{\href{https://github.com/lsernau/grh-lean/blob/72fead102f663d5c7a407964a13dc0cf9d977207/grh-compact/src/New/Columns/ColumnDefs.lean\#L42-L43}{{\leanfont Data.\allowbreak{}columnLevel}}}
\expandafter\newcommand\csname LN@principalLevellecolumnLevel\endcsname{\href{https://github.com/lsernau/grh-lean/blob/72fead102f663d5c7a407964a13dc0cf9d977207/grh-compact/src/New/Columns/ColumnDefs.lean\#L45-L47}{{\leanfont Data.\allowbreak{}principalLevel\_le\_columnLevel}}}
\expandafter\newcommand\csname LN@isLevelcolumnLevel\endcsname{\href{https://github.com/lsernau/grh-lean/blob/72fead102f663d5c7a407964a13dc0cf9d977207/grh-compact/src/New/Columns/ColumnDefs.lean\#L49-L50}{{\leanfont Data.\allowbreak{}isLevel\_columnLevel}}}
\expandafter\newcommand\csname LN@lastIndexEquiv\endcsname{\href{https://github.com/lsernau/grh-lean/blob/72fead102f663d5c7a407964a13dc0cf9d977207/grh-compact/src/New/Columns/ColumnDefs.lean\#L54-L54}{{\leanfont Data.\allowbreak{}lastIndexEquiv}}}
\expandafter\newcommand\csname LN@branching\endcsname{\href{https://github.com/lsernau/grh-lean/blob/72fead102f663d5c7a407964a13dc0cf9d977207/grh-compact/src/New/Columns/ColumnDefs.lean\#L58-L62}{{\leanfont Data.\allowbreak{}branching}}}
\expandafter\newcommand\csname LN@cardX\endcsname{\href{https://github.com/lsernau/grh-lean/blob/72fead102f663d5c7a407964a13dc0cf9d977207/grh-compact/src/New/Columns/ColumnDefs.lean\#L64-L64}{{\leanfont Data.\allowbreak{}card\_X}}}
\expandafter\newcommand\csname LN@classAutcolumnLeveleqbot\endcsname{\href{https://github.com/lsernau/grh-lean/blob/72fead102f663d5c7a407964a13dc0cf9d977207/grh-compact/src/New/Columns/ColumnDefs.lean\#L66-L67}{{\leanfont Data.\allowbreak{}classAut\_columnLevel\_eq\_bot}}}
\expandafter\newcommand\csname LN@Column\endcsname{\href{https://github.com/lsernau/grh-lean/blob/72fead102f663d5c7a407964a13dc0cf9d977207/grh-compact/src/New/Columns/ColumnDefs.lean\#L71-L71}{{\leanfont Data.\allowbreak{}Column}}}
\expandafter\newcommand\csname LN@Stage\endcsname{\href{https://github.com/lsernau/grh-lean/blob/72fead102f663d5c7a407964a13dc0cf9d977207/grh-compact/src/New/Columns/ColumnDefs.lean\#L74-L74}{{\leanfont Data.\allowbreak{}Stage}}}
\expandafter\newcommand\csname LN@branchExp\endcsname{\href{https://github.com/lsernau/grh-lean/blob/72fead102f663d5c7a407964a13dc0cf9d977207/grh-compact/src/New/Columns/ColumnDefs.lean\#L78-L79}{{\leanfont Data.\allowbreak{}branchExp}}}
\expandafter\newcommand\csname LN@lebranchExp\endcsname{\href{https://github.com/lsernau/grh-lean/blob/72fead102f663d5c7a407964a13dc0cf9d977207/grh-compact/src/New/Columns/ColumnDefs.lean\#L83-L84}{{\leanfont Data.\allowbreak{}le\_branchExp}}}
\expandafter\newcommand\csname LN@borderU\endcsname{\href{https://github.com/lsernau/grh-lean/blob/72fead102f663d5c7a407964a13dc0cf9d977207/grh-compact/src/New/Columns/ColumnDefs.lean\#L88-L90}{{\leanfont Data.\allowbreak{}borderU}}}
\expandafter\newcommand\csname LN@restrictToY\endcsname{\href{https://github.com/lsernau/grh-lean/blob/72fead102f663d5c7a407964a13dc0cf9d977207/grh-compact/src/New/Columns/ColumnDefs.lean\#L93-L101}{{\leanfont Data.\allowbreak{}restrictToY}}}
\expandafter\newcommand\csname LN@instUbCastRealRrLogLx\endcsname{\href{https://github.com/lsernau/grh-lean/blob/72fead102f663d5c7a407964a13dc0cf9d977207/grh-compact/src/New/Columns/JetExponentBound.lean\#L15-L17}{{\leanfont Data.\allowbreak{}instUbCastRealRrLogLx₁}}}
\expandafter\newcommand\csname LN@mdvdorderOf\endcsname{\href{https://github.com/lsernau/grh-lean/blob/72fead102f663d5c7a407964a13dc0cf9d977207/grh-compact/src/New/Columns/JetExponentBound.lean\#L21-L28}{{\leanfont Data.\allowbreak{}IsServed.\allowbreak{}m\_dvd\_orderOf}}}
\expandafter\newcommand\csname LN@leloss\endcsname{\href{https://github.com/lsernau/grh-lean/blob/72fead102f663d5c7a407964a13dc0cf9d977207/grh-compact/src/New/Columns/JetExponentBound.lean\#L30-L30}{{\leanfont Data.\allowbreak{}le\_loss}}}
\expandafter\newcommand\csname LN@placeOfIntmulmemtopsiff\endcsname{\href{https://github.com/lsernau/grh-lean/blob/72fead102f663d5c7a407964a13dc0cf9d977207/grh-compact/src/New/Columns/JetExponentBound.lean\#L34-L36}{{\leanfont Data.\allowbreak{}placeOfInt\_mul\_mem\_tops\_iff}}}
\expandafter\newcommand\csname LN@pmeminducedPlace\endcsname{\href{https://github.com/lsernau/grh-lean/blob/72fead102f663d5c7a407964a13dc0cf9d977207/grh-compact/src/New/Columns/JetExponentBound.lean\#L40-L40}{{\leanfont Data.\allowbreak{}p\_mem\_inducedPlace}}}
\expandafter\newcommand\csname LN@inducedPlacemul\endcsname{\href{https://github.com/lsernau/grh-lean/blob/72fead102f663d5c7a407964a13dc0cf9d977207/grh-compact/src/New/Columns/JetExponentBound.lean\#L42-L43}{{\leanfont Data.\allowbreak{}inducedPlace\_mul}}}
\expandafter\newcommand\csname LN@inducedPlacenegpowmulofodd\endcsname{\href{https://github.com/lsernau/grh-lean/blob/72fead102f663d5c7a407964a13dc0cf9d977207/grh-compact/src/New/Columns/JetExponentBound.lean\#L51-L53}{{\leanfont Data.\allowbreak{}inducedPlace\_neg\_pow\_mul\_of\_odd}}}
\expandafter\newcommand\csname LN@existsisMarked\endcsname{\href{https://github.com/lsernau/grh-lean/blob/72fead102f663d5c7a407964a13dc0cf9d977207/grh-compact/src/New/Columns/JetExponentBound.lean\#L57-L62}{{\leanfont Data.\allowbreak{}exists\_isMarked}}}
\expandafter\newcommand\csname LN@inducedPlaceeqorstariff\endcsname{\href{https://github.com/lsernau/grh-lean/blob/72fead102f663d5c7a407964a13dc0cf9d977207/grh-compact/src/New/Columns/JetExponentBound.lean\#L66-L68}{{\leanfont Data.\allowbreak{}inducedPlace\_eq\_or\_star\_iff}}}
\expandafter\newcommand\csname LN@cardle\endcsname{\href{https://github.com/lsernau/grh-lean/blob/72fead102f663d5c7a407964a13dc0cf9d977207/grh-compact/src/New/Columns/JetExponentBound.lean\#L70-L71}{{\leanfont Data.\allowbreak{}IsMarked.\allowbreak{}card\_le}}}
\expandafter\newcommand\csname LN@notdvdgProdmulN\endcsname{\href{https://github.com/lsernau/grh-lean/blob/72fead102f663d5c7a407964a13dc0cf9d977207/grh-compact/src/New/Columns/JetExponentBound.lean\#L73-L75}{{\leanfont Data.\allowbreak{}IsMarked.\allowbreak{}not\_dvd\_gProd\_mul\_N}}}
\expandafter\newcommand\csname LN@notMemthick\endcsname{\href{https://github.com/lsernau/grh-lean/blob/72fead102f663d5c7a407964a13dc0cf9d977207/grh-compact/src/New/Columns/JetExponentBound.lean\#L79-L84}{{\leanfont Data.\allowbreak{}IsMarked.\allowbreak{}notMem\_thick}}}
\expandafter\newcommand\csname LN@suminducedPlacesubstareqzero\endcsname{\href{https://github.com/lsernau/grh-lean/blob/72fead102f663d5c7a407964a13dc0cf9d977207/grh-compact/src/New/Columns/JetExponentBound.lean\#L88-L92}{{\leanfont Data.\allowbreak{}IsMarked.\allowbreak{}sum\_inducedPlace\_sub\_star\_eq\_zero}}}
\expandafter\newcommand\csname LN@pmem\endcsname{\href{https://github.com/lsernau/grh-lean/blob/72fead102f663d5c7a407964a13dc0cf9d977207/grh-compact/src/New/Columns/JetExponentBound.lean\#L94-L94}{{\leanfont Data.\allowbreak{}IsMarked.\allowbreak{}p\_mem}}}
\expandafter\newcommand\csname LN@valuationp\endcsname{\href{https://github.com/lsernau/grh-lean/blob/72fead102f663d5c7a407964a13dc0cf9d977207/grh-compact/src/New/Columns/JetExponentBound.lean\#L96-L97}{{\leanfont Data.\allowbreak{}IsMarked.\allowbreak{}valuation\_p}}}
\expandafter\newcommand\csname LN@starne\endcsname{\href{https://github.com/lsernau/grh-lean/blob/72fead102f663d5c7a407964a13dc0cf9d977207/grh-compact/src/New/Columns/JetExponentBound.lean\#L99-L99}{{\leanfont Data.\allowbreak{}IsMarked.\allowbreak{}star\_ne}}}
\expandafter\newcommand\csname LN@maxCompactAtlecolumnLevel\endcsname{\href{https://github.com/lsernau/grh-lean/blob/72fead102f663d5c7a407964a13dc0cf9d977207/grh-compact/src/New/Columns/JetExponentBound.lean\#L101-L104}{{\leanfont Data.\allowbreak{}IsMarked.\allowbreak{}maxCompactAt\_le\_columnLevel}}}
\expandafter\newcommand\csname LN@markOperator\endcsname{\href{https://github.com/lsernau/grh-lean/blob/72fead102f663d5c7a407964a13dc0cf9d977207/grh-compact/src/New/Columns/JetExponentBound.lean\#L108-L111}{{\leanfont Data.\allowbreak{}IsMarked.\allowbreak{}markOperator}}}
\expandafter\newcommand\csname LN@jetExponent\endcsname{\href{https://github.com/lsernau/grh-lean/blob/72fead102f663d5c7a407964a13dc0cf9d977207/grh-compact/src/New/Columns/JetExponentBound.lean\#L115-L118}{{\leanfont Data.\allowbreak{}jetExponent}}}
\expandafter\newcommand\csname LN@neandoddanddvd\endcsname{\href{https://github.com/lsernau/grh-lean/blob/72fead102f663d5c7a407964a13dc0cf9d977207/grh-compact/src/New/Columns/JetExponentBound.lean\#L121-L122}{{\leanfont Data.\allowbreak{}IsCollision.\allowbreak{}ne\_and\_odd\_and\_dvd}}}
\expandafter\newcommand\csname LN@collisionSumle\endcsname{\href{https://github.com/lsernau/grh-lean/blob/72fead102f663d5c7a407964a13dc0cf9d977207/grh-compact/src/New/Columns/JetExponentBound.lean\#L128-L133}{{\leanfont Data.\allowbreak{}collisionSum\_le}}}
\expandafter\newcommand\csname LN@mulcollisionSumleloss\endcsname{\href{https://github.com/lsernau/grh-lean/blob/72fead102f663d5c7a407964a13dc0cf9d977207/grh-compact/src/New/Columns/JetExponentBound.lean\#L138-L144}{{\leanfont Data.\allowbreak{}mul\_collisionSum\_le\_loss}}}
\expandafter\newcommand\csname LN@collisionSumevenPlaceseqzero\endcsname{\href{https://github.com/lsernau/grh-lean/blob/72fead102f663d5c7a407964a13dc0cf9d977207/grh-compact/src/New/Columns/JetExponentBound.lean\#L148-L151}{{\leanfont Data.\allowbreak{}IsFullMark.\allowbreak{}collisionSum\_evenPlaces\_eq\_zero}}}
\expandafter\newcommand\csname LN@lowColumnslowBoxesbounds\endcsname{\href{https://github.com/lsernau/grh-lean/blob/72fead102f663d5c7a407964a13dc0cf9d977207/grh-compact/src/New/Columns/JetExponentBound.lean\#L155-L163}{{\leanfont Data.\allowbreak{}lowColumns\_lowBoxes\_bounds}}}
\expandafter\newcommand\csname LN@lejetExponent\endcsname{\href{https://github.com/lsernau/grh-lean/blob/72fead102f663d5c7a407964a13dc0cf9d977207/grh-compact/src/New/Columns/JetExponentBound.lean\#L168-L204}{{\leanfont Data.\allowbreak{}le\_jetExponent}}}
\expandafter\newcommand\csname LN@sumlowColumnsle\endcsname{\href{https://github.com/lsernau/grh-lean/blob/72fead102f663d5c7a407964a13dc0cf9d977207/grh-compact/src/New/Columns/JetExponentBound.lean\#L208-L221}{{\leanfont Data.\allowbreak{}IsMarked.\allowbreak{}sum\_lowColumns\_le}}}
\expandafter\newcommand\csname LN@liftLoss\endcsname{\href{https://github.com/lsernau/grh-lean/blob/72fead102f663d5c7a407964a13dc0cf9d977207/grh-compact/src/New/Columns/LocalStatement.lean\#L42-L43}{{\leanfont Data.\allowbreak{}liftLoss}}}
\expandafter\newcommand\csname LN@existshecketransfer\endcsname{\href{https://github.com/lsernau/grh-lean/blob/72fead102f663d5c7a407964a13dc0cf9d977207/grh-compact/src/New/Columns/LocalStatement.lean\#L50-L101}{{\leanfont Data.\allowbreak{}IsGoodDatum.\allowbreak{}exists\_hecke\_transfer}}}
\expandafter\newcommand\csname LN@memraises\endcsname{\href{https://github.com/lsernau/grh-lean/blob/72fead102f663d5c7a407964a13dc0cf9d977207/grh-compact/src/New/Columns/LocalStatement.lean\#L106-L143}{{\leanfont Data.\allowbreak{}IsGoodDatum.\allowbreak{}mem\_raises}}}
\expandafter\newcommand\csname LN@existsmovemarkOperator\endcsname{\href{https://github.com/lsernau/grh-lean/blob/72fead102f663d5c7a407964a13dc0cf9d977207/grh-compact/src/New/Columns/LocalStatement.lean\#L148-L179}{{\leanfont Data.\allowbreak{}IsGoodDatum.\allowbreak{}exists\_move\_markOperator}}}
\expandafter\newcommand\csname LN@gaugeLemarkOperatorpow\endcsname{\href{https://github.com/lsernau/grh-lean/blob/72fead102f663d5c7a407964a13dc0cf9d977207/grh-compact/src/New/Columns/LocalStatement.lean\#L184-L194}{{\leanfont Data.\allowbreak{}IsGoodDatum.\allowbreak{}gaugeLe\_markOperator\_pow}}}
\expandafter\newcommand\csname LN@weylNormSqBmule\endcsname{\href{https://github.com/lsernau/grh-lean/blob/72fead102f663d5c7a407964a13dc0cf9d977207/grh-compact/src/New/Columns/LocalStatement.lean\#L201-L228}{{\leanfont Data.\allowbreak{}weylNormSq\_Bmu\_le}}}
\expandafter\newcommand\csname LQ@weylNormSqBmule\endcsname{\leanquoteline{0}{lemma Data.weylNormSq\_Bmu\_le (H : IsGoodDatum 1 Ψ) \{σ : ℂ ≃+* ℂ\} \{φ B\} (hB : ∀ g,}\leanquoteline{2}{weylNormSq (σ ∘ φ.1.1 g) ≤ exp B) h : weylNormSq (σ ∘ (D.Bμ .one Ψ φ).1 h) ≤}\leanquoteline{2}{exp (B + D.heightPerRow + 2 * ∑ p ∈ D.nonsplitPrimes, D.gain p * log p + D.budget)}}
\expandafter\newcommand\csname LN@hasTypeprodmarkOperatorpow\endcsname{\href{https://github.com/lsernau/grh-lean/blob/72fead102f663d5c7a407964a13dc0cf9d977207/grh-compact/src/New/Columns/LocalStatement.lean\#L233-L240}{{\leanfont Data.\allowbreak{}IsMarked.\allowbreak{}hasType\_prod\_markOperator\_pow}}}
\expandafter\newcommand\csname LN@gaugeLeone\endcsname{\href{https://github.com/lsernau/grh-lean/blob/72fead102f663d5c7a407964a13dc0cf9d977207/grh-compact/src/New/Columns/LocalStatement.lean\#L246-L250}{{\leanfont Data.\allowbreak{}IsGoodDatum.\allowbreak{}gaugeLe\_one}}}
\expandafter\newcommand\csname LN@gaugeLeBmuofisMarked\endcsname{\href{https://github.com/lsernau/grh-lean/blob/72fead102f663d5c7a407964a13dc0cf9d977207/grh-compact/src/New/Columns/LocalStatement.lean\#L255-L260}{{\leanfont Data.\allowbreak{}IsGoodDatum.\allowbreak{}gaugeLe\_Bmu\_of\_isMarked}}}
\expandafter\newcommand\csname LN@gaugeLecollisionSum\endcsname{\href{https://github.com/lsernau/grh-lean/blob/72fead102f663d5c7a407964a13dc0cf9d977207/grh-compact/src/New/Columns/LocalStatement.lean\#L264-L268}{{\leanfont Data.\allowbreak{}IsGoodDatum.\allowbreak{}gaugeLe\_collisionSum}}}
\expandafter\newcommand\csname LN@gaugeLepowjetCeilsmul\endcsname{\href{https://github.com/lsernau/grh-lean/blob/72fead102f663d5c7a407964a13dc0cf9d977207/grh-compact/src/New/Columns/LocalStatement.lean\#L272-L276}{{\leanfont Data.\allowbreak{}IsGoodDatum.\allowbreak{}gaugeLe\_pow\_jetCeil\_smul}}}
\expandafter\newcommand\csname LN@existsindexDvdBmudivisible\endcsname{\href{https://github.com/lsernau/grh-lean/blob/72fead102f663d5c7a407964a13dc0cf9d977207/grh-compact/src/New/Columns/LocalStatement.lean\#L285-L307}{{\leanfont Data.\allowbreak{}exists\_indexDvd\_Bmu\_divisible}}}
\expandafter\newcommand\csname LQ@existsindexDvdBmudivisible\endcsname{\leanquoteline{0}{lemma Data.exists\_indexDvd\_Bmu\_divisible \{β : ℤ\} (hc : β < ite (p ∈ D.nonsplitPrimes) (D.gain p) 0 -}\leanquoteline{2}{D.allowedLoss p + 1) : ∃ l : ℕ, l ≤ D.allowedLoss p * finrank ℂ D.Column ∧ ∃ S,}\leanquoteline{2}{IndexDvd integralColumns S (p \textasciicircum{} l) ∧ ∀ ϑ Ψ, IsGoodDatum ϑ Ψ → p.Coprime ϑ.levelNorm → ∀ φ}\leanquoteline{2}{∈ S, φ ∈ D.ramifiedColumns → (p : ℂ) \textasciicircum{} (-β) • Bμ ϑ Ψ φ ∈ integralFormsAt ⊤ 1}}
\expandafter\newcommand\csname LN@hasTypemarkOperator\endcsname{\href{https://github.com/lsernau/grh-lean/blob/72fead102f663d5c7a407964a13dc0cf9d977207/grh-compact/src/New/Columns/MarkedPrimeType.lean\#L15-L67}{{\leanfont Data.\allowbreak{}IsMarked.\allowbreak{}hasType\_markOperator}}}
\expandafter\newcommand\csname LN@heightPerRow\endcsname{\href{https://github.com/lsernau/grh-lean/blob/72fead102f663d5c7a407964a13dc0cf9d977207/grh-compact/src/New/Columns/SizeEstimates.lean\#L20-L21}{{\leanfont Data.\allowbreak{}heightPerRow}}}
\expandafter\newcommand\csname LN@nonsplitPrimes\endcsname{\href{https://github.com/lsernau/grh-lean/blob/72fead102f663d5c7a407964a13dc0cf9d977207/grh-compact/src/New/Columns/SizeEstimates.lean\#L25-L25}{{\leanfont Data.\allowbreak{}nonsplitPrimes}}}
\expandafter\newcommand\csname LN@primeSumeqsumlocalAsk\endcsname{\href{https://github.com/lsernau/grh-lean/blob/72fead102f663d5c7a407964a13dc0cf9d977207/grh-compact/src/New/Columns/SizeEstimates.lean\#L29-L39}{{\leanfont Data.\allowbreak{}primeSum\_eq\_sum\_localAsk}}}
\expandafter\newcommand\csname LN@diagramWeightlebyindex\endcsname{\href{https://github.com/lsernau/grh-lean/blob/72fead102f663d5c7a407964a13dc0cf9d977207/grh-compact/src/New/Columns/SizeEstimates.lean\#L43-L48}{{\leanfont Data.\allowbreak{}diagramWeight\_le\_by\_index}}}
\expandafter\newcommand\csname LN@normmapstarpiGsq\endcsname{\href{https://github.com/lsernau/grh-lean/blob/72fead102f663d5c7a407964a13dc0cf9d977207/grh-compact/src/New/Columns/SizeEstimates.lean\#L50-L51}{{\leanfont Data.\allowbreak{}norm\_map\_star\_piG\_sq}}}
\expandafter\newcommand\csname LN@jetSizeleadd\endcsname{\href{https://github.com/lsernau/grh-lean/blob/72fead102f663d5c7a407964a13dc0cf9d977207/grh-compact/src/New/Columns/SizeEstimates.lean\#L55-L60}{{\leanfont Data.\allowbreak{}jetSize\_le\_add}}}
\expandafter\newcommand\csname LN@cardstdClassesleexpmullog\endcsname{\href{https://github.com/lsernau/grh-lean/blob/72fead102f663d5c7a407964a13dc0cf9d977207/grh-compact/src/New/Columns/SizeEstimates.lean\#L64-L73}{{\leanfont Data.\allowbreak{}card\_stdClasses\_le\_exp\_mul\_log}}}
\expandafter\newcommand\csname LN@cardstdClassesleexp\endcsname{\href{https://github.com/lsernau/grh-lean/blob/72fead102f663d5c7a407964a13dc0cf9d977207/grh-compact/src/New/Columns/SizeEstimates.lean\#L75-L79}{{\leanfont Data.\allowbreak{}card\_stdClasses\_le\_exp}}}
\expandafter\newcommand\csname LN@allowedLoss\endcsname{\href{https://github.com/lsernau/grh-lean/blob/72fead102f663d5c7a407964a13dc0cf9d977207/grh-compact/src/New/Columns/SizeEstimates.lean\#L83-L84}{{\leanfont Data.\allowbreak{}allowedLoss}}}
\expandafter\newcommand\csname LN@sumallowedLossmullogle\endcsname{\href{https://github.com/lsernau/grh-lean/blob/72fead102f663d5c7a407964a13dc0cf9d977207/grh-compact/src/New/Columns/SizeEstimates.lean\#L88-L103}{{\leanfont Data.\allowbreak{}sum\_allowedLoss\_mul\_log\_le}}}
\expandafter\newcommand\csname LQ@sumallowedLossmullogle\endcsname{\leanquoteline{0}{lemma sum\_allowedLoss\_mul\_log\_le U : ∑ p ∈ U, D.allowedLoss p * .log p ≤ D.budget}}
\expandafter\newcommand\csname LN@normdetCharAtdiagNle\endcsname{\href{https://github.com/lsernau/grh-lean/blob/72fead102f663d5c7a407964a13dc0cf9d977207/grh-compact/src/New/Columns/TwistedHeckeType.lean\#L31-L41}{{\leanfont Data.\allowbreak{}norm\_detCharAt\_diagN\_le}}}
\expandafter\newcommand\csname LN@namelastColDatummemraises\endcsname{\href{https://github.com/lsernau/grh-lean/blob/72fead102f663d5c7a407964a13dc0cf9d977207/grh-compact/src/New/Columns/TwistedHeckeType.lean\#L46-L76}{{\leanfont Data.\allowbreak{}name\_lastColDatum\_mem\_raises}}}
\expandafter\newcommand\csname LN@chainpairTranslateseqsmulchain\endcsname{\href{https://github.com/lsernau/grh-lean/blob/72fead102f663d5c7a407964a13dc0cf9d977207/grh-compact/src/New/Columns/TwistedHeckeType.lean\#L82-L106}{{\leanfont Data.\allowbreak{}chain₂\_pairTranslates\_eq\_smul\_chain}}}
\expandafter\newcommand\csname LN@gaugeLeBmuatN\endcsname{\href{https://github.com/lsernau/grh-lean/blob/72fead102f663d5c7a407964a13dc0cf9d977207/grh-compact/src/New/Columns/TwistedHeckeType.lean\#L113-L135}{{\leanfont Data.\allowbreak{}gaugeLe\_Bmu\_at\_N}}}
\expandafter\newcommand\csname LN@hasTypemarkOperatora\endcsname{\href{https://github.com/lsernau/grh-lean/blob/72fead102f663d5c7a407964a13dc0cf9d977207/grh-compact/src/New/Columns/TwistedHeckeType.lean\#L146-L239}{{\leanfont Columns.\allowbreak{}hasType\_markOperator}}}
\newcommand{\ConcordanceRows}{\crow{detKernel}{\href{https://github.com/lsernau/grh-lean/blob/72fead102f663d5c7a407964a13dc0cf9d977207/grh-compact/src/New/Columns/ChainValuation.lean\#L24-L25}{{\leanfont detKernel}}}{ChainValuation.lean}{24--25}
\crow{detPeriod}{\href{https://github.com/lsernau/grh-lean/blob/72fead102f663d5c7a407964a13dc0cf9d977207/grh-compact/src/New/Columns/ChainValuation.lean\#L29-L29}{{\leanfont detPeriod}}}{ChainValuation.lean}{29--29}
\crow{detKernelembFinmul}{\href{https://github.com/lsernau/grh-lean/blob/72fead102f663d5c7a407964a13dc0cf9d977207/grh-compact/src/New/Columns/ChainValuation.lean\#L35-L39}{{\leanfont detKernel\_embFin\_mul}}}{ChainValuation.lean}{35--39}
\crow{detPeriodeqstdClassSumsum}{\href{https://github.com/lsernau/grh-lean/blob/72fead102f663d5c7a407964a13dc0cf9d977207/grh-compact/src/New/Columns/ChainValuation.lean\#L44-L49}{{\leanfont detPeriod\_eq\_stdClassSum\_sum}}}{ChainValuation.lean}{44--49}
\crow{detPeriodinsert}{\href{https://github.com/lsernau/grh-lean/blob/72fead102f663d5c7a407964a13dc0cf9d977207/grh-compact/src/New/Columns/ChainValuation.lean\#L55-L83}{{\leanfont detPeriod\_insert}}}{ChainValuation.lean}{55--83}
\crow{IsShort}{\href{https://github.com/lsernau/grh-lean/blob/72fead102f663d5c7a407964a13dc0cf9d977207/grh-compact/src/New/Columns/ChainValuation.lean\#L91-L92}{{\leanfont IsShort}}}{ChainValuation.lean}{91--92}
\crow{subtranslate}{\href{https://github.com/lsernau/grh-lean/blob/72fead102f663d5c7a407964a13dc0cf9d977207/grh-compact/src/New/Columns/ChainValuation.lean\#L94-L99}{{\leanfont IsShort.\allowbreak{}sub\_translate}}}{ChainValuation.lean}{94--99}
\crow{spanisShorteqtop}{\href{https://github.com/lsernau/grh-lean/blob/72fead102f663d5c7a407964a13dc0cf9d977207/grh-compact/src/New/Columns/ChainValuation.lean\#L105-L147}{{\leanfont span\_isShort\_eq\_top}}}{ChainValuation.lean}{105--147}
\crow{algFormCentralRep}{\href{https://github.com/lsernau/grh-lean/blob/72fead102f663d5c7a407964a13dc0cf9d977207/grh-compact/src/New/Columns/ChainValuation.lean\#L151-L153}{{\leanfont algFormCentralRep}}}{ChainValuation.lean}{151--153}
\crow{notisChosenofNmem}{\href{https://github.com/lsernau/grh-lean/blob/72fead102f663d5c7a407964a13dc0cf9d977207/grh-compact/src/New/Columns/ChainValuation.lean\#L157-L159}{{\leanfont Data.\allowbreak{}not\_isChosen\_of\_N\_mem}}}{ChainValuation.lean}{157--159}
\crow{memcolumnLevel}{\href{https://github.com/lsernau/grh-lean/blob/72fead102f663d5c7a407964a13dc0cf9d977207/grh-compact/src/New/Columns/ChainValuation.lean\#L161-L164}{{\leanfont Data.\allowbreak{}mem\_columnLevel}}}{ChainValuation.lean}{161--164}
\crow{memchosenLevelofsupportedOn}{\href{https://github.com/lsernau/grh-lean/blob/72fead102f663d5c7a407964a13dc0cf9d977207/grh-compact/src/New/Columns/ChainValuation.lean\#L166-L168}{{\leanfont Data.\allowbreak{}mem\_chosenLevel\_of\_supportedOn}}}{ChainValuation.lean}{166--168}
\crow{existscentralGen}{\href{https://github.com/lsernau/grh-lean/blob/72fead102f663d5c7a407964a13dc0cf9d977207/grh-compact/src/New/Columns/ChainValuation.lean\#L173-L181}{{\leanfont Data.\allowbreak{}exists\_centralGen}}}{ChainValuation.lean}{173--181}
\crow{centralGen}{\href{https://github.com/lsernau/grh-lean/blob/72fead102f663d5c7a407964a13dc0cf9d977207/grh-compact/src/New/Columns/ChainValuation.lean\#L185-L185}{{\leanfont Data.\allowbreak{}centralGen}}}{ChainValuation.lean}{185--185}
\crow{scalarFincentralGenmem}{\href{https://github.com/lsernau/grh-lean/blob/72fead102f663d5c7a407964a13dc0cf9d977207/grh-compact/src/New/Columns/ChainValuation.lean\#L187-L189}{{\leanfont Data.\allowbreak{}scalarFin\_centralGen\_mem}}}{ChainValuation.lean}{187--189}
\crow{centralRep}{\href{https://github.com/lsernau/grh-lean/blob/72fead102f663d5c7a407964a13dc0cf9d977207/grh-compact/src/New/Columns/ChainValuation.lean\#L193-L193}{{\leanfont Data.\allowbreak{}centralRep}}}{ChainValuation.lean}{193--193}
\crow{ramifyingOp}{\href{https://github.com/lsernau/grh-lean/blob/72fead102f663d5c7a407964a13dc0cf9d977207/grh-compact/src/New/Columns/ChainValuation.lean\#L197-L198}{{\leanfont Data.\allowbreak{}ramifyingOp}}}{ChainValuation.lean}{197--198}
\crow{ramifiedColumns}{\href{https://github.com/lsernau/grh-lean/blob/72fead102f663d5c7a407964a13dc0cf9d977207/grh-compact/src/New/Columns/ChainValuation.lean\#L202-L202}{{\leanfont Data.\allowbreak{}ramifiedColumns}}}{ChainValuation.lean}{202--202}
\crow{existseqsubcentralGen}{\href{https://github.com/lsernau/grh-lean/blob/72fead102f663d5c7a407964a13dc0cf9d977207/grh-compact/src/New/Columns/ChainValuation.lean\#L207-L209}{{\leanfont Data.\allowbreak{}exists\_eq\_sub\_centralGen}}}{ChainValuation.lean}{207--209}
\crow{sumheightmulkappale}{\href{https://github.com/lsernau/grh-lean/blob/72fead102f663d5c7a407964a13dc0cf9d977207/grh-compact/src/New/Columns/ChainValuation.lean\#L211-L212}{{\leanfont Data.\allowbreak{}sum\_height\_mul\_kappa\_le}}}{ChainValuation.lean}{211--212}
\crow{finrankkeronesubcentralGenle}{\href{https://github.com/lsernau/grh-lean/blob/72fead102f663d5c7a407964a13dc0cf9d977207/grh-compact/src/New/Columns/ChainValuation.lean\#L214-L232}{{\leanfont Data.\allowbreak{}finrank\_ker\_one\_sub\_centralGen\_le}}}{ChainValuation.lean}{214--232}
\crow{ramifyingOpfinrankcomapandshort}{\href{https://github.com/lsernau/grh-lean/blob/72fead102f663d5c7a407964a13dc0cf9d977207/grh-compact/src/New/Columns/ChainValuation.lean\#L236-L246}{{\leanfont Data.\allowbreak{}ramifyingOp\_finrank\_comap\_and\_short}}}{ChainValuation.lean}{236--246}
\crow{spanshortramifiedColumns}{\href{https://github.com/lsernau/grh-lean/blob/72fead102f663d5c7a407964a13dc0cf9d977207/grh-compact/src/New/Columns/ChainValuation.lean\#L251-L257}{{\leanfont Data.\allowbreak{}span\_short\_ramifiedColumns}}}{ChainValuation.lean}{251--257}
\crow{pairTranslates}{\href{https://github.com/lsernau/grh-lean/blob/72fead102f663d5c7a407964a13dc0cf9d977207/grh-compact/src/New/Columns/ChainValuation.lean\#L264-L265}{{\leanfont Data.\allowbreak{}pairTranslates}}}{ChainValuation.lean}{264--265}
\crow{admissibleTranslations}{\href{https://github.com/lsernau/grh-lean/blob/72fead102f663d5c7a407964a13dc0cf9d977207/grh-compact/src/New/Columns/ChainValuation.lean\#L269-L269}{{\leanfont Data.\allowbreak{}admissibleTranslations}}}{ChainValuation.lean}{269--269}
\crow{IsGoodDatum}{\href{https://github.com/lsernau/grh-lean/blob/72fead102f663d5c7a407964a13dc0cf9d977207/grh-compact/src/New/Columns/ChainValuation.lean\#L275-L279}{{\leanfont Data.\allowbreak{}IsGoodDatum}}}{ChainValuation.lean}{275--279}
\crow{notMempiPlacesofne}{\href{https://github.com/lsernau/grh-lean/blob/72fead102f663d5c7a407964a13dc0cf9d977207/grh-compact/src/New/Columns/ChainValuation.lean\#L281-L283}{{\leanfont Data.\allowbreak{}notMem\_piPlaces\_of\_ne}}}{ChainValuation.lean}{281--283}
\crow{memintegralAfter}{\href{https://github.com/lsernau/grh-lean/blob/72fead102f663d5c7a407964a13dc0cf9d977207/grh-compact/src/New/Columns/ChainValuation.lean\#L288-L292}{{\leanfont Data.\allowbreak{}mem\_integralAfter}}}{ChainValuation.lean}{288--292}
\crow{projOntranslateDatum}{\href{https://github.com/lsernau/grh-lean/blob/72fead102f663d5c7a407964a13dc0cf9d977207/grh-compact/src/New/Columns/ChainValuation.lean\#L294-L297}{{\leanfont Data.\allowbreak{}projOn\_translateDatum}}}{ChainValuation.lean}{294--297}
\crow{existschaineqsum}{\href{https://github.com/lsernau/grh-lean/blob/72fead102f663d5c7a407964a13dc0cf9d977207/grh-compact/src/New/Columns/ChainValuation.lean\#L302-L312}{{\leanfont Data.\allowbreak{}exists\_chain\_eq\_sum}}}{ChainValuation.lean}{302--312}
\crow{translateDatummemstandardData}{\href{https://github.com/lsernau/grh-lean/blob/72fead102f663d5c7a407964a13dc0cf9d977207/grh-compact/src/New/Columns/ChainValuation.lean\#L314-L319}{{\leanfont Data.\allowbreak{}IsGoodDatum.\allowbreak{}translateDatum\_mem\_standardData}}}{ChainValuation.lean}{314--319}
\crow{gaugeLechain}{\href{https://github.com/lsernau/grh-lean/blob/72fead102f663d5c7a407964a13dc0cf9d977207/grh-compact/src/New/Columns/ChainValuation.lean\#L324-L335}{{\leanfont Data.\allowbreak{}IsGoodDatum.\allowbreak{}gaugeLe\_chain}}}{ChainValuation.lean}{324--335}
\crow{normstarpiGleone}{\href{https://github.com/lsernau/grh-lean/blob/72fead102f663d5c7a407964a13dc0cf9d977207/grh-compact/src/New/Columns/ChainValuation.lean\#L337-L338}{{\leanfont Data.\allowbreak{}norm\_star\_piG\_le\_one}}}{ChainValuation.lean}{337--338}
\crow{nezero}{\href{https://github.com/lsernau/grh-lean/blob/72fead102f663d5c7a407964a13dc0cf9d977207/grh-compact/src/New/Columns/ColumnDefs.lean\#L22-L22}{{\leanfont Data.\allowbreak{}π\_ne\_zero}}}{ColumnDefs.lean}{22--22}
\crow{piPlaces}{\href{https://github.com/lsernau/grh-lean/blob/72fead102f663d5c7a407964a13dc0cf9d977207/grh-compact/src/New/Columns/ColumnDefs.lean\#L25-L26}{{\leanfont Data.\allowbreak{}piPlaces}}}{ColumnDefs.lean}{25--26}
\crow{mempiPlaces}{\href{https://github.com/lsernau/grh-lean/blob/72fead102f663d5c7a407964a13dc0cf9d977207/grh-compact/src/New/Columns/ColumnDefs.lean\#L28-L28}{{\leanfont Data.\allowbreak{}mem\_piPlaces}}}{ColumnDefs.lean}{28--28}
\crow{isSplitLevelpiPlaces}{\href{https://github.com/lsernau/grh-lean/blob/72fead102f663d5c7a407964a13dc0cf9d977207/grh-compact/src/New/Columns/ColumnDefs.lean\#L30-L32}{{\leanfont Data.\allowbreak{}isSplitLevel\_piPlaces}}}{ColumnDefs.lean}{30--32}
\crow{cardpiPlaces}{\href{https://github.com/lsernau/grh-lean/blob/72fead102f663d5c7a407964a13dc0cf9d977207/grh-compact/src/New/Columns/ColumnDefs.lean\#L35-L38}{{\leanfont Data.\allowbreak{}card\_piPlaces}}}{ColumnDefs.lean}{35--38}
\crow{columnLevel}{\href{https://github.com/lsernau/grh-lean/blob/72fead102f663d5c7a407964a13dc0cf9d977207/grh-compact/src/New/Columns/ColumnDefs.lean\#L42-L43}{{\leanfont Data.\allowbreak{}columnLevel}}}{ColumnDefs.lean}{42--43}
\crow{principalLevellecolumnLevel}{\href{https://github.com/lsernau/grh-lean/blob/72fead102f663d5c7a407964a13dc0cf9d977207/grh-compact/src/New/Columns/ColumnDefs.lean\#L45-L47}{{\leanfont Data.\allowbreak{}principalLevel\_le\_columnLevel}}}{ColumnDefs.lean}{45--47}
\crow{isLevelcolumnLevel}{\href{https://github.com/lsernau/grh-lean/blob/72fead102f663d5c7a407964a13dc0cf9d977207/grh-compact/src/New/Columns/ColumnDefs.lean\#L49-L50}{{\leanfont Data.\allowbreak{}isLevel\_columnLevel}}}{ColumnDefs.lean}{49--50}
\crow{lastIndexEquiv}{\href{https://github.com/lsernau/grh-lean/blob/72fead102f663d5c7a407964a13dc0cf9d977207/grh-compact/src/New/Columns/ColumnDefs.lean\#L54-L54}{{\leanfont Data.\allowbreak{}lastIndexEquiv}}}{ColumnDefs.lean}{54--54}
\crow{branching}{\href{https://github.com/lsernau/grh-lean/blob/72fead102f663d5c7a407964a13dc0cf9d977207/grh-compact/src/New/Columns/ColumnDefs.lean\#L58-L62}{{\leanfont Data.\allowbreak{}branching}}}{ColumnDefs.lean}{58--62}
\crow{cardX}{\href{https://github.com/lsernau/grh-lean/blob/72fead102f663d5c7a407964a13dc0cf9d977207/grh-compact/src/New/Columns/ColumnDefs.lean\#L64-L64}{{\leanfont Data.\allowbreak{}card\_X}}}{ColumnDefs.lean}{64--64}
\crow{classAutcolumnLeveleqbot}{\href{https://github.com/lsernau/grh-lean/blob/72fead102f663d5c7a407964a13dc0cf9d977207/grh-compact/src/New/Columns/ColumnDefs.lean\#L66-L67}{{\leanfont Data.\allowbreak{}classAut\_columnLevel\_eq\_bot}}}{ColumnDefs.lean}{66--67}
\crow{Column}{\href{https://github.com/lsernau/grh-lean/blob/72fead102f663d5c7a407964a13dc0cf9d977207/grh-compact/src/New/Columns/ColumnDefs.lean\#L71-L71}{{\leanfont Data.\allowbreak{}Column}}}{ColumnDefs.lean}{71--71}
\crow{Stage}{\href{https://github.com/lsernau/grh-lean/blob/72fead102f663d5c7a407964a13dc0cf9d977207/grh-compact/src/New/Columns/ColumnDefs.lean\#L74-L74}{{\leanfont Data.\allowbreak{}Stage}}}{ColumnDefs.lean}{74--74}
\crow{branchExp}{\href{https://github.com/lsernau/grh-lean/blob/72fead102f663d5c7a407964a13dc0cf9d977207/grh-compact/src/New/Columns/ColumnDefs.lean\#L78-L79}{{\leanfont Data.\allowbreak{}branchExp}}}{ColumnDefs.lean}{78--79}
\crow{lebranchExp}{\href{https://github.com/lsernau/grh-lean/blob/72fead102f663d5c7a407964a13dc0cf9d977207/grh-compact/src/New/Columns/ColumnDefs.lean\#L83-L84}{{\leanfont Data.\allowbreak{}le\_branchExp}}}{ColumnDefs.lean}{83--84}
\crow{borderU}{\href{https://github.com/lsernau/grh-lean/blob/72fead102f663d5c7a407964a13dc0cf9d977207/grh-compact/src/New/Columns/ColumnDefs.lean\#L88-L90}{{\leanfont Data.\allowbreak{}borderU}}}{ColumnDefs.lean}{88--90}
\crow{restrictToY}{\href{https://github.com/lsernau/grh-lean/blob/72fead102f663d5c7a407964a13dc0cf9d977207/grh-compact/src/New/Columns/ColumnDefs.lean\#L93-L101}{{\leanfont Data.\allowbreak{}restrictToY}}}{ColumnDefs.lean}{93--101}
\crow{instUbCastRealRrLogLx}{\href{https://github.com/lsernau/grh-lean/blob/72fead102f663d5c7a407964a13dc0cf9d977207/grh-compact/src/New/Columns/JetExponentBound.lean\#L15-L17}{{\leanfont Data.\allowbreak{}instUbCastRealRrLogLx₁}}}{JetExponentBound.lean}{15--17}
\crow{mdvdorderOf}{\href{https://github.com/lsernau/grh-lean/blob/72fead102f663d5c7a407964a13dc0cf9d977207/grh-compact/src/New/Columns/JetExponentBound.lean\#L21-L28}{{\leanfont Data.\allowbreak{}IsServed.\allowbreak{}m\_dvd\_orderOf}}}{JetExponentBound.lean}{21--28}
\crow{leloss}{\href{https://github.com/lsernau/grh-lean/blob/72fead102f663d5c7a407964a13dc0cf9d977207/grh-compact/src/New/Columns/JetExponentBound.lean\#L30-L30}{{\leanfont Data.\allowbreak{}le\_loss}}}{JetExponentBound.lean}{30--30}
\crow{placeOfIntmulmemtopsiff}{\href{https://github.com/lsernau/grh-lean/blob/72fead102f663d5c7a407964a13dc0cf9d977207/grh-compact/src/New/Columns/JetExponentBound.lean\#L34-L36}{{\leanfont Data.\allowbreak{}placeOfInt\_mul\_mem\_tops\_iff}}}{JetExponentBound.lean}{34--36}
\crow{pmeminducedPlace}{\href{https://github.com/lsernau/grh-lean/blob/72fead102f663d5c7a407964a13dc0cf9d977207/grh-compact/src/New/Columns/JetExponentBound.lean\#L40-L40}{{\leanfont Data.\allowbreak{}p\_mem\_inducedPlace}}}{JetExponentBound.lean}{40--40}
\crow{inducedPlacemul}{\href{https://github.com/lsernau/grh-lean/blob/72fead102f663d5c7a407964a13dc0cf9d977207/grh-compact/src/New/Columns/JetExponentBound.lean\#L42-L43}{{\leanfont Data.\allowbreak{}inducedPlace\_mul}}}{JetExponentBound.lean}{42--43}
\crow{inducedPlacenegpowmulofodd}{\href{https://github.com/lsernau/grh-lean/blob/72fead102f663d5c7a407964a13dc0cf9d977207/grh-compact/src/New/Columns/JetExponentBound.lean\#L51-L53}{{\leanfont Data.\allowbreak{}inducedPlace\_neg\_pow\_mul\_of\_odd}}}{JetExponentBound.lean}{51--53}
\crow{existsisMarked}{\href{https://github.com/lsernau/grh-lean/blob/72fead102f663d5c7a407964a13dc0cf9d977207/grh-compact/src/New/Columns/JetExponentBound.lean\#L57-L62}{{\leanfont Data.\allowbreak{}exists\_isMarked}}}{JetExponentBound.lean}{57--62}
\crow{inducedPlaceeqorstariff}{\href{https://github.com/lsernau/grh-lean/blob/72fead102f663d5c7a407964a13dc0cf9d977207/grh-compact/src/New/Columns/JetExponentBound.lean\#L66-L68}{{\leanfont Data.\allowbreak{}inducedPlace\_eq\_or\_star\_iff}}}{JetExponentBound.lean}{66--68}
\crow{cardle}{\href{https://github.com/lsernau/grh-lean/blob/72fead102f663d5c7a407964a13dc0cf9d977207/grh-compact/src/New/Columns/JetExponentBound.lean\#L70-L71}{{\leanfont Data.\allowbreak{}IsMarked.\allowbreak{}card\_le}}}{JetExponentBound.lean}{70--71}
\crow{notdvdgProdmulN}{\href{https://github.com/lsernau/grh-lean/blob/72fead102f663d5c7a407964a13dc0cf9d977207/grh-compact/src/New/Columns/JetExponentBound.lean\#L73-L75}{{\leanfont Data.\allowbreak{}IsMarked.\allowbreak{}not\_dvd\_gProd\_mul\_N}}}{JetExponentBound.lean}{73--75}
\crow{notMemthick}{\href{https://github.com/lsernau/grh-lean/blob/72fead102f663d5c7a407964a13dc0cf9d977207/grh-compact/src/New/Columns/JetExponentBound.lean\#L79-L84}{{\leanfont Data.\allowbreak{}IsMarked.\allowbreak{}notMem\_thick}}}{JetExponentBound.lean}{79--84}
\crow{suminducedPlacesubstareqzero}{\href{https://github.com/lsernau/grh-lean/blob/72fead102f663d5c7a407964a13dc0cf9d977207/grh-compact/src/New/Columns/JetExponentBound.lean\#L88-L92}{{\leanfont Data.\allowbreak{}IsMarked.\allowbreak{}sum\_inducedPlace\_sub\_star\_eq\_zero}}}{JetExponentBound.lean}{88--92}
\crow{pmem}{\href{https://github.com/lsernau/grh-lean/blob/72fead102f663d5c7a407964a13dc0cf9d977207/grh-compact/src/New/Columns/JetExponentBound.lean\#L94-L94}{{\leanfont Data.\allowbreak{}IsMarked.\allowbreak{}p\_mem}}}{JetExponentBound.lean}{94--94}
\crow{valuationp}{\href{https://github.com/lsernau/grh-lean/blob/72fead102f663d5c7a407964a13dc0cf9d977207/grh-compact/src/New/Columns/JetExponentBound.lean\#L96-L97}{{\leanfont Data.\allowbreak{}IsMarked.\allowbreak{}valuation\_p}}}{JetExponentBound.lean}{96--97}
\crow{starne}{\href{https://github.com/lsernau/grh-lean/blob/72fead102f663d5c7a407964a13dc0cf9d977207/grh-compact/src/New/Columns/JetExponentBound.lean\#L99-L99}{{\leanfont Data.\allowbreak{}IsMarked.\allowbreak{}star\_ne}}}{JetExponentBound.lean}{99--99}
\crow{maxCompactAtlecolumnLevel}{\href{https://github.com/lsernau/grh-lean/blob/72fead102f663d5c7a407964a13dc0cf9d977207/grh-compact/src/New/Columns/JetExponentBound.lean\#L101-L104}{{\leanfont Data.\allowbreak{}IsMarked.\allowbreak{}maxCompactAt\_le\_columnLevel}}}{JetExponentBound.lean}{101--104}
\crow{markOperator}{\href{https://github.com/lsernau/grh-lean/blob/72fead102f663d5c7a407964a13dc0cf9d977207/grh-compact/src/New/Columns/JetExponentBound.lean\#L108-L111}{{\leanfont Data.\allowbreak{}IsMarked.\allowbreak{}markOperator}}}{JetExponentBound.lean}{108--111}
\crow{jetExponent}{\href{https://github.com/lsernau/grh-lean/blob/72fead102f663d5c7a407964a13dc0cf9d977207/grh-compact/src/New/Columns/JetExponentBound.lean\#L115-L118}{{\leanfont Data.\allowbreak{}jetExponent}}}{JetExponentBound.lean}{115--118}
\crow{neandoddanddvd}{\href{https://github.com/lsernau/grh-lean/blob/72fead102f663d5c7a407964a13dc0cf9d977207/grh-compact/src/New/Columns/JetExponentBound.lean\#L121-L122}{{\leanfont Data.\allowbreak{}IsCollision.\allowbreak{}ne\_and\_odd\_and\_dvd}}}{JetExponentBound.lean}{121--122}
\crow{collisionSumle}{\href{https://github.com/lsernau/grh-lean/blob/72fead102f663d5c7a407964a13dc0cf9d977207/grh-compact/src/New/Columns/JetExponentBound.lean\#L128-L133}{{\leanfont Data.\allowbreak{}collisionSum\_le}}}{JetExponentBound.lean}{128--133}
\crow{mulcollisionSumleloss}{\href{https://github.com/lsernau/grh-lean/blob/72fead102f663d5c7a407964a13dc0cf9d977207/grh-compact/src/New/Columns/JetExponentBound.lean\#L138-L144}{{\leanfont Data.\allowbreak{}mul\_collisionSum\_le\_loss}}}{JetExponentBound.lean}{138--144}
\crow{collisionSumevenPlaceseqzero}{\href{https://github.com/lsernau/grh-lean/blob/72fead102f663d5c7a407964a13dc0cf9d977207/grh-compact/src/New/Columns/JetExponentBound.lean\#L148-L151}{{\leanfont Data.\allowbreak{}IsFullMark.\allowbreak{}collisionSum\_evenPlaces\_eq\_zero}}}{JetExponentBound.lean}{148--151}
\crow{lowColumnslowBoxesbounds}{\href{https://github.com/lsernau/grh-lean/blob/72fead102f663d5c7a407964a13dc0cf9d977207/grh-compact/src/New/Columns/JetExponentBound.lean\#L155-L163}{{\leanfont Data.\allowbreak{}lowColumns\_lowBoxes\_bounds}}}{JetExponentBound.lean}{155--163}
\crow{lejetExponent}{\href{https://github.com/lsernau/grh-lean/blob/72fead102f663d5c7a407964a13dc0cf9d977207/grh-compact/src/New/Columns/JetExponentBound.lean\#L168-L204}{{\leanfont Data.\allowbreak{}le\_jetExponent}}}{JetExponentBound.lean}{168--204}
\crow{sumlowColumnsle}{\href{https://github.com/lsernau/grh-lean/blob/72fead102f663d5c7a407964a13dc0cf9d977207/grh-compact/src/New/Columns/JetExponentBound.lean\#L208-L221}{{\leanfont Data.\allowbreak{}IsMarked.\allowbreak{}sum\_lowColumns\_le}}}{JetExponentBound.lean}{208--221}
\crow{liftLoss}{\href{https://github.com/lsernau/grh-lean/blob/72fead102f663d5c7a407964a13dc0cf9d977207/grh-compact/src/New/Columns/LocalStatement.lean\#L42-L43}{{\leanfont Data.\allowbreak{}liftLoss}}}{LocalStatement.lean}{42--43}
\crow{existshecketransfer}{\href{https://github.com/lsernau/grh-lean/blob/72fead102f663d5c7a407964a13dc0cf9d977207/grh-compact/src/New/Columns/LocalStatement.lean\#L50-L101}{{\leanfont Data.\allowbreak{}IsGoodDatum.\allowbreak{}exists\_hecke\_transfer}}}{LocalStatement.lean}{50--101}
\crow{memraises}{\href{https://github.com/lsernau/grh-lean/blob/72fead102f663d5c7a407964a13dc0cf9d977207/grh-compact/src/New/Columns/LocalStatement.lean\#L106-L143}{{\leanfont Data.\allowbreak{}IsGoodDatum.\allowbreak{}mem\_raises}}}{LocalStatement.lean}{106--143}
\crow{existsmovemarkOperator}{\href{https://github.com/lsernau/grh-lean/blob/72fead102f663d5c7a407964a13dc0cf9d977207/grh-compact/src/New/Columns/LocalStatement.lean\#L148-L179}{{\leanfont Data.\allowbreak{}IsGoodDatum.\allowbreak{}exists\_move\_markOperator}}}{LocalStatement.lean}{148--179}
\crow{gaugeLemarkOperatorpow}{\href{https://github.com/lsernau/grh-lean/blob/72fead102f663d5c7a407964a13dc0cf9d977207/grh-compact/src/New/Columns/LocalStatement.lean\#L184-L194}{{\leanfont Data.\allowbreak{}IsGoodDatum.\allowbreak{}gaugeLe\_markOperator\_pow}}}{LocalStatement.lean}{184--194}
\crow{weylNormSqBmule}{\href{https://github.com/lsernau/grh-lean/blob/72fead102f663d5c7a407964a13dc0cf9d977207/grh-compact/src/New/Columns/LocalStatement.lean\#L201-L228}{{\leanfont Data.\allowbreak{}weylNormSq\_Bmu\_le}}}{LocalStatement.lean}{201--228}
\crow{hasTypeprodmarkOperatorpow}{\href{https://github.com/lsernau/grh-lean/blob/72fead102f663d5c7a407964a13dc0cf9d977207/grh-compact/src/New/Columns/LocalStatement.lean\#L233-L240}{{\leanfont Data.\allowbreak{}IsMarked.\allowbreak{}hasType\_prod\_markOperator\_pow}}}{LocalStatement.lean}{233--240}
\crow{gaugeLeone}{\href{https://github.com/lsernau/grh-lean/blob/72fead102f663d5c7a407964a13dc0cf9d977207/grh-compact/src/New/Columns/LocalStatement.lean\#L246-L250}{{\leanfont Data.\allowbreak{}IsGoodDatum.\allowbreak{}gaugeLe\_one}}}{LocalStatement.lean}{246--250}
\crow{gaugeLeBmuofisMarked}{\href{https://github.com/lsernau/grh-lean/blob/72fead102f663d5c7a407964a13dc0cf9d977207/grh-compact/src/New/Columns/LocalStatement.lean\#L255-L260}{{\leanfont Data.\allowbreak{}IsGoodDatum.\allowbreak{}gaugeLe\_Bmu\_of\_isMarked}}}{LocalStatement.lean}{255--260}
\crow{gaugeLecollisionSum}{\href{https://github.com/lsernau/grh-lean/blob/72fead102f663d5c7a407964a13dc0cf9d977207/grh-compact/src/New/Columns/LocalStatement.lean\#L264-L268}{{\leanfont Data.\allowbreak{}IsGoodDatum.\allowbreak{}gaugeLe\_collisionSum}}}{LocalStatement.lean}{264--268}
\crow{gaugeLepowjetCeilsmul}{\href{https://github.com/lsernau/grh-lean/blob/72fead102f663d5c7a407964a13dc0cf9d977207/grh-compact/src/New/Columns/LocalStatement.lean\#L272-L276}{{\leanfont Data.\allowbreak{}IsGoodDatum.\allowbreak{}gaugeLe\_pow\_jetCeil\_smul}}}{LocalStatement.lean}{272--276}
\crow{existsindexDvdBmudivisible}{\href{https://github.com/lsernau/grh-lean/blob/72fead102f663d5c7a407964a13dc0cf9d977207/grh-compact/src/New/Columns/LocalStatement.lean\#L285-L307}{{\leanfont Data.\allowbreak{}exists\_indexDvd\_Bmu\_divisible}}}{LocalStatement.lean}{285--307}
\crow{hasTypemarkOperator}{\href{https://github.com/lsernau/grh-lean/blob/72fead102f663d5c7a407964a13dc0cf9d977207/grh-compact/src/New/Columns/MarkedPrimeType.lean\#L15-L67}{{\leanfont Data.\allowbreak{}IsMarked.\allowbreak{}hasType\_markOperator}}}{MarkedPrimeType.lean}{15--67}
\crow{heightPerRow}{\href{https://github.com/lsernau/grh-lean/blob/72fead102f663d5c7a407964a13dc0cf9d977207/grh-compact/src/New/Columns/SizeEstimates.lean\#L20-L21}{{\leanfont Data.\allowbreak{}heightPerRow}}}{SizeEstimates.lean}{20--21}
\crow{nonsplitPrimes}{\href{https://github.com/lsernau/grh-lean/blob/72fead102f663d5c7a407964a13dc0cf9d977207/grh-compact/src/New/Columns/SizeEstimates.lean\#L25-L25}{{\leanfont Data.\allowbreak{}nonsplitPrimes}}}{SizeEstimates.lean}{25--25}
\crow{primeSumeqsumlocalAsk}{\href{https://github.com/lsernau/grh-lean/blob/72fead102f663d5c7a407964a13dc0cf9d977207/grh-compact/src/New/Columns/SizeEstimates.lean\#L29-L39}{{\leanfont Data.\allowbreak{}primeSum\_eq\_sum\_localAsk}}}{SizeEstimates.lean}{29--39}
\crow{diagramWeightlebyindex}{\href{https://github.com/lsernau/grh-lean/blob/72fead102f663d5c7a407964a13dc0cf9d977207/grh-compact/src/New/Columns/SizeEstimates.lean\#L43-L48}{{\leanfont Data.\allowbreak{}diagramWeight\_le\_by\_index}}}{SizeEstimates.lean}{43--48}
\crow{normmapstarpiGsq}{\href{https://github.com/lsernau/grh-lean/blob/72fead102f663d5c7a407964a13dc0cf9d977207/grh-compact/src/New/Columns/SizeEstimates.lean\#L50-L51}{{\leanfont Data.\allowbreak{}norm\_map\_star\_piG\_sq}}}{SizeEstimates.lean}{50--51}
\crow{jetSizeleadd}{\href{https://github.com/lsernau/grh-lean/blob/72fead102f663d5c7a407964a13dc0cf9d977207/grh-compact/src/New/Columns/SizeEstimates.lean\#L55-L60}{{\leanfont Data.\allowbreak{}jetSize\_le\_add}}}{SizeEstimates.lean}{55--60}
\crow{cardstdClassesleexpmullog}{\href{https://github.com/lsernau/grh-lean/blob/72fead102f663d5c7a407964a13dc0cf9d977207/grh-compact/src/New/Columns/SizeEstimates.lean\#L64-L73}{{\leanfont Data.\allowbreak{}card\_stdClasses\_le\_exp\_mul\_log}}}{SizeEstimates.lean}{64--73}
\crow{cardstdClassesleexp}{\href{https://github.com/lsernau/grh-lean/blob/72fead102f663d5c7a407964a13dc0cf9d977207/grh-compact/src/New/Columns/SizeEstimates.lean\#L75-L79}{{\leanfont Data.\allowbreak{}card\_stdClasses\_le\_exp}}}{SizeEstimates.lean}{75--79}
\crow{allowedLoss}{\href{https://github.com/lsernau/grh-lean/blob/72fead102f663d5c7a407964a13dc0cf9d977207/grh-compact/src/New/Columns/SizeEstimates.lean\#L83-L84}{{\leanfont Data.\allowbreak{}allowedLoss}}}{SizeEstimates.lean}{83--84}
\crow{sumallowedLossmullogle}{\href{https://github.com/lsernau/grh-lean/blob/72fead102f663d5c7a407964a13dc0cf9d977207/grh-compact/src/New/Columns/SizeEstimates.lean\#L88-L103}{{\leanfont Data.\allowbreak{}sum\_allowedLoss\_mul\_log\_le}}}{SizeEstimates.lean}{88--103}
\crow{normdetCharAtdiagNle}{\href{https://github.com/lsernau/grh-lean/blob/72fead102f663d5c7a407964a13dc0cf9d977207/grh-compact/src/New/Columns/TwistedHeckeType.lean\#L31-L41}{{\leanfont Data.\allowbreak{}norm\_detCharAt\_diagN\_le}}}{TwistedHeckeType.lean}{31--41}
\crow{namelastColDatummemraises}{\href{https://github.com/lsernau/grh-lean/blob/72fead102f663d5c7a407964a13dc0cf9d977207/grh-compact/src/New/Columns/TwistedHeckeType.lean\#L46-L76}{{\leanfont Data.\allowbreak{}name\_lastColDatum\_mem\_raises}}}{TwistedHeckeType.lean}{46--76}
\crow{chainpairTranslateseqsmulchain}{\href{https://github.com/lsernau/grh-lean/blob/72fead102f663d5c7a407964a13dc0cf9d977207/grh-compact/src/New/Columns/TwistedHeckeType.lean\#L82-L106}{{\leanfont Data.\allowbreak{}chain₂\_pairTranslates\_eq\_smul\_chain}}}{TwistedHeckeType.lean}{82--106}
\crow{gaugeLeBmuatN}{\href{https://github.com/lsernau/grh-lean/blob/72fead102f663d5c7a407964a13dc0cf9d977207/grh-compact/src/New/Columns/TwistedHeckeType.lean\#L113-L135}{{\leanfont Data.\allowbreak{}gaugeLe\_Bmu\_at\_N}}}{TwistedHeckeType.lean}{113--135}
\crow{hasTypemarkOperatora}{\href{https://github.com/lsernau/grh-lean/blob/72fead102f663d5c7a407964a13dc0cf9d977207/grh-compact/src/New/Columns/TwistedHeckeType.lean\#L146-L239}{{\leanfont Columns.\allowbreak{}hasType\_markOperator}}}{TwistedHeckeType.lean}{146--239}}
% the notations of the defined functions: one slot for each explicit argument of the Lean declaration
\newcommand{\detK}[5]{\mathrm{dk}_{#1}^{#2,#3}[#4;#5]}
\newcommand{\detP}[5]{\mathrm{dp}_{#1}^{#2,#3}[#4;#5]}
\newcommand{\short}[2]{#2\text{ is }#1\text{-short}}
\newcommand{\cgen}[2]{\mathfrak{c}_{#1}(#2)}
\newcommand{\cenrep}[1]{\operatorname{cen}_{#1}}
\newcommand{\Colram}[1]{\mathrm{Col}^{\mathrm{ram}}(#1)}
\newcommand{\pairT}[2]{\operatorname{pt}_{D}(#1)_{#2}}
\newcommand{\admT}[1]{\mathrm{Adm}_{D}(#1)}
\newcommand{\good}[2]{#2\text{ is good for }(D,#1)}
\newcommand{\piPl}[1]{\operatorname{Pl}_{\pi}(#1)}
\newcommand{\colL}[2]{K^{\mathrm{col}}_{#2}(#1)}
\newcommand{\lstEq}[1]{\mathrm{lst}_{#1}}
\newcommand{\Brm}[1]{\mathrm{Br}(#1)}
\newcommand{\Col}[1]{\mathrm{Col}(#1)}
\newcommand{\Stg}[2]{\mathrm{St}_{#2}(#1)}
\newcommand{\brexp}[2]{\mathrm{be}_{#1}(#2)}
\newcommand{\bdU}[1]{\mathrm{bd}_{#1}}
\newcommand{\resY}[1]{\operatorname{restr}_{#1}}
\newcommand{\markop}[2]{\Upsilon_{D,#1}(#2)}
\newcommand{\jetexp}[3]{\mathrm{je}_{D}(#1,#2,#3)}
\newcommand{\liftloss}[2]{\mathrm{ll}_{D,p}(#1,#2)}
\newcommand{\hrow}[1]{h_{\mathrm{row}}(#1)}
\newcommand{\nsp}[1]{\mathrm{NS}(#1)}
\newcommand{\err}[2]{\mathrm{err}(#1,#2)}
\newcommand{\UnitMembers}{105}
\newcommand{\UnitInNote}{105}
\newcommand{\UnitProofs}{46}
\newcommand{\UnitSteps}{331}
\newcommand{\UnitAllSteps}{350}


\linespread{1.08}                       % 11 pt on 14.7 pt
\setlength{\parindent}{0pt}
\setlength{\parskip}{5.5pt plus 1.5pt minus 0.5pt}
\setlength{\abovedisplayskip}{7.5pt plus 2pt minus 3pt}
\setlength{\belowdisplayskip}{7.5pt plus 2pt minus 3pt}
\setlength{\abovedisplayshortskip}{3pt plus 2pt}
\setlength{\belowdisplayshortskip}{6pt plus 2pt minus 2pt}
\emergencystretch=1.5em
\clubpenalty=10000 \widowpenalty=10000 \displaywidowpenalty=10000
\relpenalty=7000 \binoppenalty=9000
\raggedbottom
\pagestyle{fancy}\fancyhf{}\renewcommand{\headrulewidth}{0pt}
\fancyfoot[C]{\sffamily\fontsize{7.6}{9}\selectfont\color{inkfaint}\thepage\,/\,\pageref*{LastPage}}
\setlist{topsep=2pt,partopsep=0pt,parsep=0pt,itemsep=4pt,leftmargin=2.1em,labelsep=0.55em}

\newcommand{\Q}{\mathbb{Q}}
\newcommand{\Z}{\mathbb{Z}}
\newcommand{\F}{\mathbb{F}}
\newcommand{\len}{\operatorname{len}}
\newcommand{\cok}{\operatorname{cok}}
\newcommand{\End}{\operatorname{End}}
\newcommand{\Hom}{\operatorname{Hom}}
\newcommand{\Gal}{\operatorname{Gal}}
\newcommand{\diag}{\operatorname{diag}}
\newcommand{\kw}[1]{\emph{#1}}
\newcommand{\m}[1]{\mbox{$#1$}}        % a short formula that must not be broken across lines
\newcommand{\head}[1]{\par\addvspace{17pt plus 4pt minus 2pt}\Needspace{6\baselineskip}%
  \noindent{\fontsize{13.2}{16}\selectfont\bfseries\color{ink}#1}\par\nopagebreak\vspace{-1pt}}
\newcommand{\sect}[2]{\refstepcounter{section}\label{#2}\par\addvspace{19pt plus 4pt minus 2pt}\Needspace{7\baselineskip}%
  \noindent\llap{{\fontsize{13.2}{16}\selectfont\bfseries\color{inkfaint}\thesection}\hspace{9pt}}%
  {\fontsize{13.2}{16}\selectfont\bfseries\color{ink}#1}\par\nopagebreak\vspace{-1pt}}
\newcommand{\subhead}[1]{\par\addvspace{9pt plus 2pt}\Needspace{4\baselineskip}\noindent\textbf{#1}\enspace\ignorespaces}
% where a statement lives in the Lean archive
\newcommand{\leantag}{{\sffamily\bfseries\fontsize{6.3}{8}\selectfont\color{leancol}\addfontfeatures{LetterSpace=7}LEAN}}
\newcommand{\sep}{\unskip\nobreak\hspace{5pt}{\color{inkfaint}\textperiodcentered}\hspace{5pt plus 2pt}\ignorespaces}
\newcommand{\leanline}[1]{\par\nopagebreak\vspace{-2.5pt}{\fontsize{8.9}{12.4}\selectfont\color{inksoft}\raggedright
  \noindent\llap{\leantag\hspace{9pt}}#1\par}\vspace{1pt}}
% a displayed statement in a tinted box: #1 colour, #2 text
\newcommand{\statement}[2]{\par\addvspace{7pt}\begingroup\setlength{\fboxsep}{8.5pt}\setlength{\fboxrule}{0.6pt}%
  \noindent\fcolorbox{#1!30!white}{#1!3!white}{\parbox{\dimexpr\linewidth-2\fboxsep-2\fboxrule\relax}{#2}}%
  \endgroup\par\addvspace{7pt}}
\newcommand{\nt}[2]{\par\noindent\hangindent=1.1em\hangafter=1 #1\enspace{\color{inksoft}#2}\par\vspace{1pt}}
\newcommand{\legdot}[1]{\tikz[baseline=-0.62ex]\fill[#1] (0,0) circle (2.5pt);\hspace{3pt}}
\newcommand{\crow}[4]{\csname CS@#1\endcsname & #2 & \textcolor{inksoft}{#3} & \textcolor{inksoft}{#4}\\}
\newcommand{\leanquoteline}[2]{\par\noindent\hangindent=\dimexpr#1\fontdimen2\font+2.2em\relax\hangafter=1
  \hspace*{#1\fontdimen2\font}#2}
\newcommand{\quoted}[1]{\par\vspace{3pt}{\ttfamily\fontsize{7.5}{10.4}\selectfont\color{leancol}\raggedright
  \setlength{\parskip}{0pt}\leanquote{#1}\par}\vspace{4pt}}
\newcommand{\bx}[3]{{\sffamily\bfseries\fontsize{7.3}{9.2}\selectfont\color{#1!82!black}#2}\\[1.5pt]#3}
\newcommand{\cs}[2]{\expandafter\newcommand\csname CS@#1\endcsname{#2}}

% theorems
\newtheoremstyle{plainx}{7pt plus 2pt}{7pt plus 2pt}{\itshape}{}{\bfseries}{.}{0.5em}{}
\newtheoremstyle{uprightx}{7pt plus 2pt}{7pt plus 2pt}{}{}{\bfseries}{.}{0.5em}{}
\theoremstyle{plainx}
\newtheorem{theorem}{Theorem}[section]
\newtheorem{lemma}[theorem]{Lemma}
\newtheorem{proposition}[theorem]{Proposition}
\newtheorem{corollary}[theorem]{Corollary}
\theoremstyle{uprightx}
\newtheorem{definition}[theorem]{Definition}
\newtheorem{setting}[theorem]{Setting}          % hypotheses shared by several statements; each of them cites it by number
\newtheorem{remark}[theorem]{Remark}
\newtheorem{example}[theorem]{Example}
\renewcommand{\qedsymbol}{{\setlength{\fboxsep}{0pt}\setlength{\fboxrule}{0.5pt}\fbox{\rule{0pt}{0.46em}\hspace{0.46em}}}}
% what holds throughout the note: the ambient objects only (implicit and instance arguments), never a hypothesis
\newenvironment{conventions}{\par\addvspace{5pt}\noindent\textbf{Conventions.}\enspace\ignorespaces}{\par\addvspace{3pt}}
\makeatletter
\renewenvironment{proof}[1][Proof]{\par\pushQED{\qed}\normalfont\topsep0pt\partopsep0pt
  \trivlist\item[\hskip\labelsep\itshape #1.]\ignorespaces}{\popQED\endtrivlist\@endpefalse}
\makeatother

% ---------------------------------------------------------------------- macros of this note
\newcommand{\lnm}[1]{{\leanfont #1}}
\newcommand{\CC}{\mathbb{C}}
\newcommand{\RR}{\mathbb{R}}
\newcommand{\NN}{\mathbb{N}}
\newcommand{\OO}{\mathcal{O}}
\newcommand{\Af}[1]{\mathbb{A}_{#1,\mathrm{f}}}
\newcommand{\cj}[1]{\overline{#1}}
\newcommand{\Aut}{\operatorname{Aut}(\CC)}
\newcommand{\GL}{\operatorname{GL}}
\newcommand{\Oj}{O_{j}}                                   % the complex numbers integral at j
\newcommand{\op}{\mathfrak{o}_{p}}                         % the valuation ring of C_p
\newcommand{\ee}{\mathrm{e}}
% objects of other units, written with all their arguments (the forms of the overview and of the note on the CM field)
\newcommand{\KP}[1]{K(#1)}
\newcommand{\chiK}[1]{\chi_{\KP{#1}}}
\newcommand{\ED}[1]{E(#1)}
\newcommand{\Dmu}[1]{\mu(#1)}
\newcommand{\NnD}[1]{N_1(#1)}
\newcommand{\PlD}[1]{\operatorname{Pl}(#1)}
\newcommand{\fD}[1]{f(#1)}
\newcommand{\kD}[1]{k(#1)}
\newcommand{\tD}[1]{t(#1)}
\newcommand{\aD}[1]{a(#1)}
\newcommand{\tops}[1]{\mathrm{tops}(#1)}
\newcommand{\thick}[1]{\mathrm{thick}(#1)}
\newcommand{\nomW}[2]{w_{#1}(#2)}
\newcommand{\bud}[1]{\mathrm{budget}(#1)}
\newcommand{\pin}[2]{[#2]_{#1}}
\newcommand{\nD}[2]{n_{#2}(#1)}
\newcommand{\wtD}[3]{\mathrm{wt}_{#1}(#2,#3)}
\newcommand{\gap}[1]{\mathrm{gap}(#1)}
\newcommand{\bone}[1]{b_1(#1)}
\newcommand{\btwo}[1]{b_2(#1)}
\newcommand{\emb}[1]{\iota_{#1}}
\newcommand{\ND}[1]{N(#1)}
\newcommand{\piN}[1]{\pi_N(#1)}
\newcommand{\gD}[1]{g(#1)}
\newcommand{\pig}[1]{\pi_g(#1)}
\newcommand{\XD}[2]{X_{#2}(#1)}
\newcommand{\YD}[2]{Y_{#2}(#1)}
\newcommand{\OmP}[2]{\Omega(#1,#2)}
\newcommand{\Ask}[3]{\mathrm{Ask}_{#1}(#2,#3)}
\newcommand{\jsz}[2]{\mathrm{st}_{#1}(#2)}
\newcommand{\dirc}[1]{\mathsf{U}(#1)}
\newcommand{\lamd}[3]{\lambda_{#1,#2}(#3)}
% objects of the unit on the CM curve (unit 150), in the forms of its note
\newcommand{\ask}[4]{\mathrm{ask}_{#1,#2}(#3,#4)}
\newcommand{\forced}[2]{\mathrm{forced}(#1,#2)}
\newcommand{\gain}[2]{\mathrm{gain}(#1,#2)}
\newcommand{\loss}[2]{\mathrm{loss}(#1,#2)}
\newcommand{\rr}[2]{\mathfrak{r}(#1,#2)}
\newcommand{\DeltaP}[2]{\Delta(#1,#2)}
\newcommand{\jc}[2]{\mathrm{jc}(#1,#2)}
\newcommand{\oz}[2]{o_0(#1,#2)}
\newcommand{\coll}[3]{\mathrm{coll}(#1,#2,#3)}
\newcommand{\cc}[7]{\mathrm{cc}_{#1}(#2,#3;#4;#5,#6;#7)}
\newcommand{\kap}[3]{\varkappa_{#1}(#2,#3)}
\newcommand{\kapD}[1]{\varkappa_{#1}}
\newcommand{\last}[3]{\mathrm{last}_{#1}(#2,#3)}
\newcommand{\lastD}[2]{\mathrm{last}_{#1}(#2,\,\cdot\,)}
\newcommand{\kept}[2]{\mathrm{kept}_{#1}(#2)}
\newcommand{\kapY}[3]{\varkappa^{Y}_{#1}(#2,#3)}
\newcommand{\kapYD}[1]{\varkappa^{Y}_{#1}}
\newcommand{\nsw}[2]{\varrho_{#1}(#2)}
\newcommand{\per}[2]{\mathrm{per}(#1,#2)}
\newcommand{\Std}[2]{\mathrm{Std}_{#1}(#2)}
\newcommand{\raises}[4]{\mathrm{Raise}_{#1}(#2,#3;#4)}
\newcommand{\intvac}[1]{\theta^{\circ}_{#1}}
\newcommand{\intafter}[2]{\mathrm{Lat}(#1;#2)}
\newcommand{\ThProd}[5]{\mathrm{ThProd}_{#1,#2}(#3,#4;#5)}
% objects of the units on the seesaw (141) and on the exchange (147)
\newcommand{\Bch}[3]{\mathrm{B}_{#1}(#2,#3)}
\newcommand{\td}[3]{\mathrm{td}_{#1}(#2,#3)}
\newcommand{\lk}[3]{\mathrm{k}_{#1,#2,#3}}
\newcommand{\lkom}[3]{\omega_{#1,#2,#3}}
\newcommand{\lkjm}[3]{\mathrm{jm}_{#1,#2,#3}}
\newcommand{\lklift}[3]{\operatorname{lift}_{#1,#2,#3}}
\newcommand{\chn}[4]{\mathrm{C}_{#1,#2}(#3)_{#4}}
\newcommand{\chnn}[4]{\mathrm{C}^{(2)}_{#1,#2}(#3)_{#4}}
\newcommand{\ind}[1]{\mathfrak{l}_{j}(#1)}
\newcommand{\mexp}[2]{A_{#1}(#2)}
\newcommand{\lc}[3]{\mathrm{lc}_{#1}(#2,#3)}
\newcommand{\lb}[3]{\mathrm{lb}_{#1}(#2,#3)}
\newcommand{\even}[1]{\Xi_{j}(#1)}
\newcommand{\Colint}[1]{\mathrm{Col}_{j}(#1)}
\newcommand{\nrm}[1]{\mathrm{nrm}(#1)}
\newcommand{\Kch}[2]{K^{\mathrm{ch}}_{#2}(#1)}
\newcommand{\Wfin}[2]{\mathcal{W}_{#2}(#1)}
% objects of the background library
\newcommand{\Alg}[3]{\mathrm{Alg}_{#1}(#2,#3)}
\newcommand{\Wy}[2]{\mathrm{Wy}_{#1}(#2)}
\newcommand{\rhow}[2]{\rho_{#1,#2}}
\newcommand{\gauge}[4]{\operatorname{gauge}^{#1,#2}_{j}(#3;#4)}
\newcommand{\Gf}[1]{G_{#1}}
\newcommand{\Kf}[1]{K_{#1}}
\newcommand{\Gam}[1]{\Gamma_{#1}}
\newcommand{\Kmax}[2]{K_{#1}[#2]}
\newcommand{\Kprin}[2]{K_{#1}(#2)}
\newcommand{\Cl}[1]{\mathrm{Cl}(#1)}
\newcommand{\Tl}[2]{T_{#1,#2}}
\newcommand{\boxes}[2]{\operatorname{boxes}(#1,#2)}
\newcommand{\wn}[1]{\|#1\|^{2}}
% the operator of Definition (ramified columns): two arguments, as Lean takes them (the census lists a third)
\newcommand{\ramop}[2]{\operatorname{ram}_{#1}(#2)}

% where each declaration of the concordance is treated: \cs{key}{place}, the place naming a \label or a lettered item
\cs{detKernel}{Definition~\ref{d:det}}
\cs{detPeriod}{Definition~\ref{d:det}}
\cs{detKernelembFinmul}{Lemma~\ref{l:det}}
\cs{detPeriodeqstdClassSumsum}{Lemma~\ref{l:det}}
\cs{detPeriodinsert}{Proposition~\ref{p:detins}}
\cs{IsShort}{Definition~\ref{d:short}}
\cs{subtranslate}{Lemma~\ref{l:sub}}
\cs{spanisShorteqtop}{Proposition~\ref{p:span}}
\cs{algFormCentralRep}{Definition~\ref{d:central}}
\cs{notisChosenofNmem}{Lemma~\ref{l:level}}
\cs{memcolumnLevel}{Lemma~\ref{l:level}}
\cs{memchosenLevelofsupportedOn}{Lemma~\ref{l:level}}
\cs{existscentralGen}{Lemma~\ref{l:cgen}}
\cs{centralGen}{Definition~\ref{d:central}}
\cs{scalarFincentralGenmem}{Lemma~\ref{l:cgenmem}}
\cs{centralRep}{Definition~\ref{d:central}}
\cs{ramifyingOp}{Definition~\ref{d:central}}
\cs{ramifiedColumns}{Definition~\ref{d:central}}
\cs{existseqsubcentralGen}{Lemma~\ref{l:cgenmem}}
\cs{sumheightmulkappale}{Lemma~\ref{l:branch}}
\cs{finrankkeronesubcentralGenle}{Lemma~\ref{l:ker}}
\cs{ramifyingOpfinrankcomapandshort}{Proposition~\ref{p:ram}}
\cs{spanshortramifiedColumns}{Theorem A}
\cs{pairTranslates}{Definition~\ref{d:good}}
\cs{admissibleTranslations}{Definition~\ref{d:good}}
\cs{IsGoodDatum}{Definition~\ref{d:good}}
\cs{notMempiPlacesofne}{Lemma~\ref{l:level}}
\cs{memintegralAfter}{Lemma~\ref{l:intafter}}
\cs{projOntranslateDatum}{Lemma~\ref{l:intafter}}
\cs{existschaineqsum}{Lemma~\ref{l:chainsum}}
\cs{translateDatummemstandardData}{Lemma~\ref{l:tdstd}}
\cs{gaugeLechain}{Proposition~\ref{p:gchain}}
\cs{normstarpiGleone}{Lemma~\ref{l:level}}
\cs{nezero}{Lemma~\ref{l:pi}}
\cs{piPlaces}{Definition~\ref{d:pi}}
\cs{mempiPlaces}{Lemma~\ref{l:pi}}
\cs{isSplitLevelpiPlaces}{Lemma~\ref{l:pi}}
\cs{cardpiPlaces}{Lemma~\ref{l:pi}}
\cs{columnLevel}{Definition~\ref{d:pi}}
\cs{principalLevellecolumnLevel}{Lemma~\ref{l:level}}
\cs{isLevelcolumnLevel}{Lemma~\ref{l:level}}
\cs{lastIndexEquiv}{Definition~\ref{d:branch}}
\cs{branching}{Definition~\ref{d:branch}}
\cs{cardX}{Lemma~\ref{l:level}}
\cs{classAutcolumnLeveleqbot}{Lemma~\ref{l:level}}
\cs{Column}{Definition~\ref{d:pi}}
\cs{Stage}{Definition~\ref{d:pi}}
\cs{branchExp}{Definition~\ref{d:branch}}
\cs{lebranchExp}{Lemma~\ref{l:branch}}
\cs{borderU}{Definition~\ref{d:branch}}
\cs{restrictToY}{Definition~\ref{d:branch}}
\cs{instUbCastRealRrLogLx}{Lemma~\ref{l:served}}
\cs{mdvdorderOf}{Lemma~\ref{l:served}}
\cs{leloss}{Lemma~\ref{l:served}}
\cs{placeOfIntmulmemtopsiff}{Lemma~\ref{l:served}}
\cs{pmeminducedPlace}{Lemma~\ref{l:ind}}
\cs{inducedPlacemul}{Lemma~\ref{l:ind}}
\cs{inducedPlacenegpowmulofodd}{Lemma~\ref{l:ind}}
\cs{existsisMarked}{Lemma~\ref{l:exmark}}
\cs{inducedPlaceeqorstariff}{Lemma~\ref{l:ind}}
\cs{cardle}{Lemma~\ref{l:markprop}}
\cs{notdvdgProdmulN}{Lemma~\ref{l:markprop}}
\cs{notMemthick}{Lemma~\ref{l:thick}}
\cs{suminducedPlacesubstareqzero}{Lemma~\ref{l:thick}}
\cs{pmem}{Lemma~\ref{l:markprop}}
\cs{valuationp}{Lemma~\ref{l:markprop}}
\cs{starne}{Lemma~\ref{l:markprop}}
\cs{maxCompactAtlecolumnLevel}{Lemma~\ref{l:markprop}}
\cs{markOperator}{Definition~\ref{d:mark}}
\cs{jetExponent}{Definition~\ref{d:mark}}
\cs{neandoddanddvd}{Lemma~\ref{l:coll}}
\cs{collisionSumle}{Lemma~\ref{l:coll}}
\cs{mulcollisionSumleloss}{Lemma~\ref{l:coll}}
\cs{collisionSumevenPlaceseqzero}{Lemma~\ref{l:coll}}
\cs{lowColumnslowBoxesbounds}{Lemma~\ref{l:low}}
\cs{lejetExponent}{Proposition~\ref{p:jet}}
\cs{sumlowColumnsle}{Lemma~\ref{l:lowsum}}
\cs{liftLoss}{Definition~\ref{d:mark}}
\cs{existshecketransfer}{Proposition~\ref{p:hecke}}
\cs{memraises}{Proposition~\ref{p:raise}}
\cs{existsmovemarkOperator}{Proposition~\ref{p:move}}
\cs{gaugeLemarkOperatorpow}{Proposition~\ref{p:markpow}}
\cs{weylNormSqBmule}{Theorem C}
\cs{hasTypeprodmarkOperatorpow}{Corollary~\ref{c:typepow}}
\cs{gaugeLeone}{Lemma~\ref{l:four}}
\cs{gaugeLeBmuofisMarked}{Lemma~\ref{l:four}}
\cs{gaugeLecollisionSum}{Lemma~\ref{l:four}}
\cs{gaugeLepowjetCeilsmul}{Lemma~\ref{l:four}}
\cs{existsindexDvdBmudivisible}{Theorem B}
\cs{hasTypemarkOperator}{Proposition~\ref{p:type}}
\cs{heightPerRow}{Definition~\ref{d:size}}
\cs{nonsplitPrimes}{Definition~\ref{d:size}}
\cs{primeSumeqsumlocalAsk}{Lemma~\ref{l:prime}}
\cs{diagramWeightlebyindex}{Lemma~\ref{l:rows}}
\cs{normmapstarpiGsq}{Lemma~\ref{l:sizes}}
\cs{jetSizeleadd}{Lemma~\ref{l:sizes}}
\cs{cardstdClassesleexpmullog}{Lemma~\ref{l:classes}}
\cs{cardstdClassesleexp}{Lemma~\ref{l:classes}}
\cs{allowedLoss}{Definition~\ref{d:size}}
\cs{sumallowedLossmullogle}{Theorem D}
\cs{normdetCharAtdiagNle}{Lemma~\ref{l:detchar}}
\cs{namelastColDatummemraises}{Lemma~\ref{l:lastcol}}
\cs{chainpairTranslateseqsmulchain}{Proposition~\ref{p:pairs}}
\cs{gaugeLeBmuatN}{Proposition~\ref{p:atN}}
\cs{hasTypemarkOperatora}{Proposition~\ref{p:typez}}

\begin{document}
\color{ink}

{\raggedright
{\sffamily\bfseries\fontsize{7.2}{9}\selectfont\color{leancol}\addfontfeatures{LetterSpace=8}READBACK NOTE \textperiodcentered{} CHAPTER V, THE COLUMNS\par}
\vspace{5pt}
{\fontsize{21.5}{25.5}\selectfont\bfseries Columns on a definite unitary group: divisibility and size of the chain of theta lifts\par}
\vspace{5pt}
{\fontsize{12.4}{16}\selectfont\itshape\color{inksoft} One unit of the Lean artifact \ArtifactName, read back into
mathematics\par}}
\vspace{8pt}
{\color{rulegrey}\hrule height 0.6pt}
\vspace{7pt}
\thispagestyle{empty}

\textbf{Summary.} A \emph{column} is an algebraic modular form on a definite unitary group of rank $\kD{D}$ over the
CM field $\ED{D}$, of a fixed weight and of a level that is principal at the chosen primes and of mirabolic type at
the primes dividing an element $\piN{D}$ of norm the level prime $\ND{D}$. A chain of $m$ theta lifts between
definite groups, followed by a restriction, sends columns to forms on a group of larger rank; by the seesaw
identity of another unit this map computes the chain of lifts of Poincar\'e series. This unit proves the three
facts about columns that Theorem~A of the handout needs. First, the columns that are \emph{ramified} at the
primes dividing $\piN{D}$ are spanned by forms that are integral at all primes and whose conjugates are all small.
Second, at every prime $p$ the chain is divisible by $p^{\gain{D}{p}}$, up to an allowed error, on a submodule of
the integral columns of small index; at an inert prime that pairs two of the top places, the divisibility is
regained on the image of scaled Hecke operators. Third, the untwisted chain enlarges the norms of all conjugates
by at most an explicit factor. The allowed errors, weighted by $\log p$, sum to at most the budget.

\statement{leancol}{\fontsize{9.6}{13.2}\selectfont\textbf{What this note is.} The artifact \ArtifactName{} is a
body of Lean code: the compact edition of a formal proof whose final theorems state the Riemann hypothesis and its
analogue for Dirichlet \emph{L}-functions. That headline is a claim. Its authors report that every file compiles
and that Lean's kernel accepts the proofs; a replay of that check is recorded in the handout when it reports;
nothing was compiled for this note, and no mathematician has read the argument through. Nothing here says that the
Riemann hypothesis is proved. This note treats one unit of the artifact, the folder of the columns: all
\UnitInNote{} of its \UnitMembers{} declarations. Every statement below was compared with its Lean text, argument
by argument. The proofs are given at the level of the structured proofs of the unit's dossier, which were read
together with the Lean text; results of other units and of the background library are stated as \emph{Inputs}
with their owners, and steps that the formal text leaves to automation and that could not be replayed without
Lean are named in §\ref{s:notes}. The Lean text read is \PinText; Appendix~A is the concordance of \ConcDecls{}
declarations.}

\sect{The statements}{s:intro}
The objects are defined in §\ref{s:setting}. In short: $D$ is a record of admissible data for parameters
$P=(N_0,d)$, with scales $x$, $x_1$, $b$ and numbers $m$, $r$; $\Col{D}$ is the space of columns and
$\Colram{D}$ its subspace of ramified columns (Definitions \ref{d:pi} and~\ref{d:central}); for a twist $\vartheta$
of~$D$ and a finitely supported function $\Psi$ on tuples of translations, $\Bch{D}{\vartheta}{\Psi}$ is the chain
of $m$ theta lifts with data given by~$\Psi$, followed by the restriction to the index set $\YD{D}{m}$ (Input~I4);
$\gain{D}{p}$ and $\loss{D}{p}$ are the divisibility asked at~$p$ and the unit of allowed loss (unit~150), and
$\bud{D}=\fD{D}\,\aD{D}\,(\log b)^3$.

\statement{arith}{\phantomsection\label{t:A}\textbf{Theorem A (short ramified columns span).}\enspace In
Setting~\ref{set:D}, the complex vector space $\Colram{D}$ is spanned by the ramified columns $\varphi$ such that
$\short{\exp(\bud{D})}{\varphi}$: the values of $\varphi$ are integral at every $p$-adic place, and
$\wn{\sigma\circ\varphi(g)}\le\exp(\bud{D})$ for every automorphism $\sigma$ of\/ $\CC$ and every point~$g$.}

\statement{arith}{\phantomsection\label{t:B}\textbf{Theorem B (divisibility of the chain at each prime).}\enspace In
Settings \ref{set:D} and~\ref{set:p}, let $\beta$ be an integer with
\[\beta\ <\ [\,p\in\nsp{D}\,]\cdot\gain{D}{p}-\err{D}{p}+1 .\]
Then there are a natural number $l\le\err{D}{p}\cdot\dim_{\CC}\Col{D}$ and an $\Oj$-submodule $S$ of\/ $\Col{D}$
such that
\begin{enumerate}[label=(\roman*)]
\item $S$ has index dividing $p^{\,l}$ in the module $\Colint{D}$ of $j$-integral columns (Input~I2), and
\item for every twist $\vartheta$ of~$D$, every finitely supported complex function $\Psi$ on the tuples of
translations of~$D$ such that $\good{\vartheta}{\Psi}$ and $p$ is prime to $\nrm{\vartheta}$, and every
$\varphi\in S$ that is a ramified column, the form $p^{-\beta}\cdot\Bch{D}{\vartheta}{\Psi}\varphi$ has $j$-gauge
at most~$1$ at~$1$.
\end{enumerate}}

\statement{arith}{\phantomsection\label{t:C}\textbf{Theorem C (size of all conjugates of the untwisted chain).}\enspace
In Setting~\ref{set:D}, let $\Psi$ be a finitely supported complex function on the tuples of translations of~$D$
such that $\good{1}{\Psi}$ for the trivial twist~$1$; let $\sigma$ be an automorphism of\/ $\CC$, $\varphi\in\Col{D}$
and $B$ a real number with $\wn{\sigma\circ\varphi(g)}\le\ee^{B}$ for every $g\in\Gf{\XD{D}{0}}$. Then for every
$h\in\Gf{\YD{D}{m}}$
\[
\wn{\sigma\circ\bigl(\Bch{D}{1}{\Psi}\varphi\bigr)(h)}\ \le\ \exp\Bigl(B+\hrow{D}+2\sum_{p\in\nsp{D}}\gain{D}{p}\,\log p+\bud{D}\Bigr).
\]}

\statement{arith}{\phantomsection\label{t:D}\textbf{Theorem D (the allowed losses fit the budget).}\enspace In
Setting~\ref{set:D}, for every finite set $U$ of natural numbers,
$\ \sum_{p\in U}\err{D}{p}\cdot\log p\ \le\ \bud{D}$.}
\leanline{\lean{spanshortramifiedColumns}\sep\lean{existsindexDvdBmudivisible}\sep\lean{weylNormSqBmule}\sep\lean{sumallowedLossmullogle}}

\textbf{How the four fit together.} The explanations of this paragraph are orientation; the deduction belongs to
the unit of the endgame, which uses these four results, and was not read for this note. A row of the handout's Theorem~A is,
after the seesaw identity, of the form $\Bch{D}{\vartheta}{\Psi}\psi$ for a ramified column~$\psi$, divided by a
rational number that collects the powers $p^{\gain{D}{p}}$ and a power of~$\ND{D}$. Theorem~A provides the
columns~$\psi$, Theorem~C bounds the size of the rows at every conjugate, and Theorem~B bounds their
denominators: at each prime the rows built from a submodule $S$ of index $p^{\,l}$ are integral after division
by $p^{\beta}$, and the exponents~$l$, weighted by $\log p$, add up to at most
$\dim\Col{D}\cdot\bud{D}$ by Theorem~D. The number $\hrow{D}$ of Theorem~C contains the term
$-\#\tops{D}\cdot\aD{D}\log\ND{D}$: the restriction divides by a power of~$\cj{\piN{D}}$, and this is how the
level prime pays for the excess of size over divisibility. The error at a prime is three units of loss; one unit
is $\fD{D}(\log b)^2\aD{D}/p$, multiplied by~$m$ when $p\equiv\pm1$ modulo~$m$.

\textbf{The mechanism at a prime.} Let $p<x_1$ be inert or ramified in $\KP{P}$. A single lift is a finite sum of
jets of theta series with complex multiplication, and by the unit on the CM curve a jet along a Schur matrix with
$L$ boxes is divisible by $p$ to an exponent of the shape $(1/(p-1)-\nsw{d}{p})L$ minus a sum of
\emph{coupling counts} (Input~I9). Proposition~\ref{p:jet} shows that this exponent, summed over the $m$ lifts,
is at least $\gain{D}{p}$ up to one unit of loss, except for the \emph{collisions}: pairs of top places exchanged
by an odd power of $-p$. A collision with a small exponent can only occur when $p$ is \emph{served}
(inert, not $\equiv\pm1$ modulo~$m$, and $-p$ of small odd order modulo~$N_0$). For such $p$ one chooses a set
$\mathcal{L}$ of primes of $\ED{D}$ above~$p$, one from each conjugate pair met by the tops, and applies at each
of them the Hecke operator of $\diag(p,\dots,p,1,\dots,1)$ scaled by a power of~$p$. On the image the
divisibility is restored (Proposition~\ref{p:markpow}), because the Hecke operator moves through every lift
(Proposition~\ref{p:hecke}) and changes the datum by a diagonal matrix that the jet bound absorbs
(Proposition~\ref{p:raise}); and the image has small index (Proposition~\ref{p:type}, whose content is a theorem
of the unit on twisted Hecke operators). At the level prime itself the data of the chain have denominator~$N$,
and a Fourier transform on the columns removes it (§\ref{s:N}).

\sect{The setting}{s:setting}
\begin{conventions}\phantomsection\label{c:conventions}
$\NN=\{0,1,2,\dots\}$. Where the formal text subtracts or divides natural numbers, the subtraction is truncated
at~$0$ and the division is rounded down; such places are marked. $[\,\mathcal{C}\,]$ is $1$ if the condition
$\mathcal{C}$ holds and $0$ otherwise. Throughout, $E$ is a CM field with its complex conjugation $x\mapsto\cj{x}$,
a CM type and Weil data (the idele and the splitting character that the Weil representation over $E$ needs; they
are those of the note on the CM field when $E=\ED{D}$). $\operatorname{Pl}(E)$ is the set of infinite places
of~$E$, and $\emb{\tau}\colon E\to\CC$ is the embedding of the CM type at $\tau\in\operatorname{Pl}(E)$. A
\emph{finite place} $w$ is a non-zero prime ideal of the ring of integers $\OO_E$; $\cj{w}$ is its conjugate,
$E_w$ the completion, and $\mathrm{q}_w$ the number of elements of the residue field. $\Af{E}$ is the ring of
finite adeles and $\widehat{\OO}_E$ its subring of integral adeles. $I$ and $I'$ are finite index sets, and $K$ denotes a subgroup of the group $\Gf{I}$ of Input~I1; a
family $\varkappa$ of \emph{column counts} is a function $(\tau,h)\mapsto\varkappa(\tau,h)\in\NN$ on
$\operatorname{Pl}(E)\times\NN$, and a family $\mathfrak{e}$ of \emph{determinant exponents} is a function
$\operatorname{Pl}(E)\to\Z$. $\Aut$ is the group of all field automorphisms of~$\CC$. For parameters
$P=(N_0,d)$, $\KP{P}=\Q(\sqrt{-d})$ and $\chiK{P}$ is its quadratic character, the Legendre symbol modulo~$d$. A \emph{$p$-adic place} is a prime number $p$ with a ring homomorphism $j\colon\CC\to\CC_p$ (Setting~\ref{set:p}). In a
statement about a commutative ring $R$ or a normed field, that structure is part of the statement.
\end{conventions}

\begin{setting}[the data]\label{set:D}
$P=(N_0,d)$ are parameters and $D$ is a record of admissible data for~$P$ (the note on the admissible data), with
real scales $x$, $x_1$, $b$ and natural numbers $m$, $r$ among its components. The functions of $D$ are written
with their argument: the exponent $\Dmu{D}$, the CM field $\ED{D}$ with its group $\PlD{D}$ of infinite places,
$\NnD{D}$, the number $\fD{D}$ of infinite places, $\kD{D}=r+1-m$, $\tD{D}=5r$, $\aD{D}$, the sets $\tops{D}$ and
$\thick{D}$, the nominal weight $\nomW{D}{\tau}$, the place $\pin{D}{\nu}\in\PlD{D}$ named by an integer~$\nu$,
the ranks $\nD{D}{i}=\kD{D}+i(\tD{D}+1)$, the weights $\wtD{D}{i}{\tau}$ and
$\bud{D}=\fD{D}\,\aD{D}\,(\log b)^3$. \emph{Size comparisons} between such quantities (inequalities that hold
because of the scale conditions $(\log b)^{100}\le x$, $x^{10^6}\le x_1$, $x_1^{10^6}\le b$ of the data) are
quoted from the note on the admissible data as ``by size''.
\end{setting}

\begin{setting}[a $p$-adic place]\label{set:p}
$p$ is a prime number and $j\colon\CC\to\CC_p$ is a ring homomorphism into the completion $\CC_p$ of an algebraic
closure of $\Q_p$, with absolute value $|\cdot|_p$, $|p|_p=p^{-1}$, and valuation ring $\op$. Put
$\Oj=\{z\in\CC:\ |j(z)|_p\le1\}$, a valuation ring of~$\CC$.
\end{setting}

\subhead{Inputs: the objects.} The definitions of this unit rest on objects of the background library and of
other units. They are recalled here, in the words of the dossier's readings, which were compared with the Lean
definitions named in the LEAN lines.

\statement{inksoft}{\textbf{Input I1 (algebraic forms; background library).}\enspace Let $I$ be a finite set.
$\Gf{I}=\{g\in\GL_I(\Af{E}):\ g^{*}g=1\}$ is the group of finite adelic points of the unitary group of the
standard definite hermitian space $E^{I}$ ($g^{*}$ the conjugate transpose); $\Gam{I}\subset\Gf{I}$ is the subgroup of
its rational points, the unitary matrices over~$E$; $\Kf{I}=\Gf{I}\cap\GL_I(\widehat{\OO}_E)$. For an integer $a$
of~$E$, $\Kprin{I}{a}$ is the principal level, the $k\in\Kf{I}$ with all entries of $k-1$ in
$a\widehat{\OO}_E$; for a finite place~$w$, $\Kmax{I}{w}$ is the group of the $k\in\Kf{I}$ whose components
away from $w$ and $\cj{w}$ are~$1$. A \emph{level} is an open subgroup $K$ of~$\Kf{I}$; $\Cl{K}$ is the finite
set of classes $\Gam{I}\backslash\Gf{I}/K$; $K$ is \emph{free} if $\Gam{I}\cap gKg^{-1}=\{1\}$ for every
$g\in\Gf{I}$.

$\Wy{I}{\varkappa}$ is the finite set of \emph{Weyl indices}: the families, over $\tau\in\operatorname{Pl}(E)$,
of standard fillings with entries in~$I$ of the diagram with $\varkappa(\tau,h)$ columns of height~$h$. The
\emph{weight representation} $\rhow{\varkappa}{\mathfrak{e}}$ sends a family $(g_\tau)$ of invertible matrices
indexed by~$I$ to $\prod_\tau\det(g_\tau)^{\mathfrak{e}(\tau)}$ times the Schur matrix of~$(g_\tau)$, a square
matrix indexed by $\Wy{I}{\varkappa}$; $\Gam{I}$ acts on $\CC^{\Wy{I}{\varkappa}}$ by
$\gamma\mapsto\rhow{\varkappa}{\mathfrak{e}}\bigl((\emb{\tau}(\cj{\gamma}))_\tau\bigr)$. The number of boxes below
height~$n'$ is $\boxes{\varkappa}{n'}=\sum_\tau\sum_{h<n'}h\,\varkappa(\tau,h)$, and the row of index~$\nu$ of
the diagram at $\tau$, for rank~$n'$, has length $\sum_{\nu<h<n'}\varkappa(\tau,h)$.

$\Alg{I}{\varkappa}{\mathfrak{e}}$ is the space of \emph{algebraic forms} of weight
$(\varkappa,\mathfrak{e})$: the functions $\varphi\colon\Gf{I}\to\CC^{\Wy{I}{\varkappa}}$ with open stabiliser
under right translation and
$\varphi(\gamma g)=\rhow{\varkappa}{\mathfrak{e}}(\gamma)\varphi(g)$ for $\gamma\in\Gam{I}$. $\Gf{I}$ acts on it
by right translation, $(\mathrm{R}(g)\varphi)(g')=\varphi(g'g)$. For a free level~$K$, evaluation at
representatives of the classes is an isomorphism from the $K$-fixed forms onto the functions
$\Cl{K}\to\CC^{\Wy{I}{\varkappa}}$; the form that takes the value $v$ at the representative of one class and $0$
at the others is a \emph{delta form}.

\emph{Gauge.} In Setting~\ref{set:p}, for $g\in\Gf{I}$ and $\xi\in\CC^{\Wy{I}{\varkappa}}$, the number
$\gauge{\varkappa}{\mathfrak{e}}{g}{\xi}$ is the largest $p$-adic absolute value of a coordinate of
$\rhow{\varkappa}{\mathfrak{e}}\bigl((g_\tau^{-1})_\tau\bigr)\,j(\xi)$, where $g_\tau\in\GL_I(\CC_p)$ is the image
of~$g$ under the continuous extension to $\Af{E}$ of the embedding $x\mapsto j(\emb{\tau}(\cj{x}))$. A form
$\varphi$ \emph{has $j$-gauge at most $c$ at} $d\in\Gf{I}$ if
$\gauge{\varkappa}{\mathfrak{e}}{gd}{\varphi(g)}\le c$ for every $g\in\Gf{I}$. The gauge is an ultrametric norm in
$\xi$ with $\gauge{\varkappa}{\mathfrak{e}}{g}{z\xi}=|j(z)|_p\gauge{\varkappa}{\mathfrak{e}}{g}{\xi}$, it does not
change when $g$ is multiplied on the right by an element that is integral and invertible at all primes above~$p$,
and for such~$g$ it is at most~$1$ exactly when all coordinates of $j(\xi)$ are integral.

\emph{Gram norm.} $\wn{\xi}$ is the real part of $\xi^{*}G\xi$, where $G$ is the integer Gram matrix of Weyl's
lattice in the weight $\varkappa$; it is positive definite, and its square root is a norm. Automorphisms of\/
$\CC$ act on vectors coordinate by coordinate.}
\leanline{in the background library: \lnm{Gf}, \lnm{stdRatPoints}, \lnm{Kf}, \lnm{principalLevel},
\lnm{maxCompactAt}, \lnm{IsLevel}, \lnm{StdClasses}, \lnm{classAut}, \lnm{WeylIndex}, \lnm{weylRep},
\lnm{Space.weightRep}, \lnm{boxes}, \lnm{diagramWeight}, \lnm{AlgForm}, \lnm{algFormRep}, \lnm{levelSpace},
\lnm{freeLevelEquiv}, \lnm{weylGauge}, \lnm{GaugeLe}, \lnm{integralFormsAt}, \lnm{embValRing}, \lnm{weylNormSq}}

\statement{inksoft}{\textbf{Input I2 (index and type; unit 149, the core).}\enspace Let $O$ be a commutative ring, $V$
an $O$-module, $M$ and $S$ submodules of~$V$ and $z\in O$. \emph{$S$ has index dividing $z$ in $M$} if for every
natural number $n$ and all $\Phi_1,\dots,\Phi_n\in M$ there is an $n\times n$ matrix $u$ over $O$ whose
determinant divides~$z$ and with $\sum_{c}u_{ac}\Phi_c\in S$ for every~$a$. An $O$-linear endomorphism $y$ of~$V$
\emph{has type $(z,s,c)$ on $M$}, for $z\in O$, a natural number $s$ and a real number~$c$, if $y(M)\subseteq M$
and there is a natural number $l\le c$ such that $y(M)+z^{s}M$ has index dividing $z^{\,l}$ in~$M$.
\emph{Types add}: if $y_1,\dots,y_n$ have types $(z,s,c_1),\dots,(z,s,c_n)$ on~$M$, then $y_1\cdots y_n$ has type
$(z,s,c_1+\dots+c_n)$ on~$M$.}
\leanline{unit 149: \lnm{IndexDvd}, \lnm{HasType}, \lnm{HasType.prod}}

\statement{inksoft}{\textbf{Input I3 (objects of the data; units 152 and 141).}\enspace In Setting~\ref{set:D}: for
each natural number $i$, $\XD{D}{i}$ is an index set with $\nD{D}{i}$ elements, one of them distinguished as the
\emph{last}; $\YD{D}{m}$ is an index set with $\nD{D}{m}-1$ elements. $\ND{D}$ is a prime number with
$\ND{D}\equiv1$ modulo $\NnD{D}\,d$ and $m\log b\le\log\ND{D}\le30\,m\log b$, and $\piN{D}$ is an integer of
$\ED{D}$ with $\piN{D}\,\cj{\piN{D}}=\ND{D}$. $\gD{D}=g_0g_1$ is the product of the two primes of the data, and
$\pig{D}$ is the integer of $\KP{P}$ in the data, with $\pig{D}\,\cj{\pig{D}}=\gD{D}$. A finite place $w$ of
$\ED{D}$ is \emph{chosen} if $\cj{\pig{D}}\in w$. For a finite set $I$, the \emph{chosen level}
$\Kch{D}{I}$ is the group of the $k\in\Kf{I}$ that are congruent to~$1$ modulo every chosen place; it
contains $\Kprin{I}{\cj{\pig{D}}}$, and if $|I|\le2\,\nD{D}{m}$, every subgroup of\/ $\Kch{D}{I}$ is free. The column
counts $\kapD{D}$, the determinant exponents $\lastD{D}{i}$, the counts $\kept{D}{h}$ and $\kapYD{D}$ are those
of the unit on the CM curve (Input~I6).

A \emph{twist} of $D$ is a continuous character $\vartheta$ of the norm-one adeles of $\ED{D}$, trivial on the
norm-one elements of $\ED{D}$ and of weight~$0$ at the infinite places, together with a finite set $Q(\vartheta)$
of finite places and characters $\vartheta_w$ of the residue fields (its \emph{tame part}) such that
$\vartheta(y)=\prod_w\vartheta_w(y\bmod w)$ for $y$ with integral finite part and trivial infinite part;
$\vartheta_w=1$ for $w\notin Q(\vartheta)$ and $\vartheta_w\ne1$ for $w\in Q(\vartheta)$; $\cj{w}\notin Q(\vartheta)$ for
$w\in Q(\vartheta)$; and $\nrm{\vartheta}=\prod_{w\in Q(\vartheta)}\mathrm{q}_w$ is prime to $\gD{D}\,\ND{D}$. The
trivial character with $Q=\emptyset$ is the trivial twist~$1$, and $\nrm{1}=1$.}
\leanline{unit 141: \lnm{Data.X}, \lnm{Data.Y}, \lnm{Data.N}, \lnm{Data.N\_modEq\_one}, \lnm{Data.log\_N\_mem\_Icc},
\lnm{Data.π}, \lnm{Data.piGBar}, \lnm{Data.IsChosen}, \lnm{Data.TameChar}, \lnm{Data.Twist},
\lnm{Data.TameChar.levelNorm}\sep unit 152: \lnm{Data.chosenLevel}, \lnm{Data.classAut\_eq\_bot\_of\_le\_chosenLevel}}

\statement{inksoft}{\textbf{Input I4 (kernels, data and the chain; background library, units 141, 147, 150).}\enspace In
Setting~\ref{set:D}, let $\vartheta$ be a twist of~$D$ and $i$ a natural number. $\Wfin{D}{i}$ is the
\emph{finite part} of the model of the Weil representation for the space of homomorphisms from $E^{\XD{D}{i}}$
to $E^{\XD{D}{i+1}}$: the complex span of the translates $\mathrm{U}(y)\,\mathrm{vac}$ of the vacuum vector by the
finite Heisenberg operators $\mathrm{U}(y)$, where $y$ is a matrix with rows $\XD{D}{i+1}$ and columns
$\XD{D}{i}$ over $\Af{\ED{D}}$. It contains the \emph{integral vacuum} $\intvac{D,i}$, a scalar multiple of the
vacuum (unit~150). For a set $\Sigma$ of finite places, $\mathrm{P}^{\Sigma}$ is the projection of $\Wfin{D}{i}$
onto the vectors that are vacua at the places of $\Sigma$ and their conjugates.

The \emph{lift kernel} $\lk{D}{\vartheta}{i}$ consists of a representation $\lkom{D}{\vartheta}{i}$ of
$\Gf{\XD{D}{i}}\times\Gf{\XD{D}{i+1}}$ on $\Wfin{D}{i}$ (the Weil representation of the pair, with the characters
of the $(i+1)$-st lift twisted by $\vartheta$) and of a linear map $\Phi\mapsto\lkjm{D}{\vartheta}{i}(\Phi)$ from
$\Wfin{D}{i}$ to the linear maps $\CC^{\Wy{\XD{D}{i}}{\kapD{D}}}\to\CC^{\Wy{\XD{D}{i+1}}{\kapD{D}}}$: the matrix of
the jets, along the Schur matrix of~$\kapD{D}$, of the Taylor series of the theta function of~$\Phi$ at the scale
$\cj{\pig{D}}\cdot\per{P}{\Dmu{D}}$. The \emph{lift} with datum~$\Phi$ is the map
$\lklift{D}{\vartheta}{i}(\Phi)\colon\Stg{D}{i}\to\Stg{D}{i+1}$,
\[
\bigl(\lklift{D}{\vartheta}{i}(\Phi)\varphi\bigr)(h)=\int_{\Gam{\XD{D}{i}}\backslash\Gf{\XD{D}{i}}}\lkjm{D}{\vartheta}{i}\bigl(\lkom{D}{\vartheta}{i}(g,h)\Phi\bigr)\varphi(g)\,\mathrm{d}g ,
\]
a finite sum over classes, the measure giving volume~$1$ to $\colL{D}{i}$. It is equivariant:
$\lklift{D}{\vartheta}{i}(\lkom{D}{\vartheta}{i}(a,b)\Phi)\,\mathrm{R}(a)=\mathrm{R}(b)\,\lklift{D}{\vartheta}{i}(\Phi)$.
For points $d\in\Gf{\XD{D}{i}}$, $d'\in\Gf{\XD{D}{i+1}}$ and a real number~$c$, in Setting~\ref{set:p},
$\raises{\lk{D}{\vartheta}{i}}{d}{d'}{c}$ is the set of the $\Phi\in\Wfin{D}{i}$ with
\[\gauge{\kapD{D}}{\lastD{D}{i+1}}{hd'}{\lkjm{D}{\vartheta}{i}(\lkom{D}{\vartheta}{i}(g,h)\Phi)\xi}\ \le\ \max(c,0)\cdot\gauge{\kapD{D}}{\lastD{D}{i}}{gd}{\xi}\]
for all $g$, $h$ and~$\xi$; it is closed under addition, and $z\Phi$ lies in
$\raises{\lk{D}{\vartheta}{i}}{d}{d'}{c'}$ when $\Phi$ lies in $\raises{\lk{D}{\vartheta}{i}}{d}{d'}{c}$ and
$|j(z)|_p\,c\le c'$ (unit~150).

A \emph{tuple of translations} of $D$ is a family $\mathsf{x}=(\mathsf{x}_i)_{i\in\NN}$ of matrices
$\mathsf{x}_i$ with rows $\XD{D}{i+1}$ and columns $\XD{D}{i}$ over $\Af{\ED{D}}$; its $i$-th datum is
$\td{D}{\mathsf{x}}{i}=\mathrm{U}(\mathsf{x}_i)\,\intvac{D,i}\in\Wfin{D}{i}$. For a family
$\Phi=(\Phi_i)_{i}$ with $\Phi_i\in\Wfin{D}{i}$, the \emph{chain} of $n$ lifts is
\[\chn{D}{\vartheta}{\Phi}{n}=\lklift{D}{\vartheta}{n-1}(\Phi_{n-1})\circ\dots\circ\lklift{D}{\vartheta}{0}(\Phi_0)\colon\ \Stg{D}{0}\to\Stg{D}{n},\]
and for a family $\Phi^{(2)}=(\Phi^{(2)}_\nu)_\nu$ with $\Phi^{(2)}_\nu\in\Wfin{D}{2\nu}\otimes\Wfin{D}{2\nu+1}$ the
\emph{chain of $n$ pairs} $\chnn{D}{\vartheta}{\Phi^{(2)}}{n}\colon\Stg{D}{0}\to\Stg{D}{2n}$ is the composite, for $\nu<n$, of
the maps that depend linearly on $\Phi^{(2)}_\nu$ and equal
$\lklift{D}{\vartheta}{2\nu+1}(\Phi')\circ\lklift{D}{\vartheta}{2\nu}(\Phi)$ for $\Phi^{(2)}_\nu=\Phi\otimes\Phi'$. For
a finitely supported complex function $\Psi$ on the tuples of translations,
\[
\Bch{D}{\vartheta}{\Psi}=\resY{D}\circ\Bigl(\sum_{\mathsf{x}}\Psi(\mathsf{x})\cdot\chn{D}{\vartheta}{\td{D}{\mathsf{x}}{\,\cdot\,}}{m}\Bigr)\Big|_{\Col{D}}
\]
with the restriction $\resY{D}$ of Definition~\ref{d:branch}. In Setting~\ref{set:p}, $\Colint{D}$ is the
$\Oj$-module of the columns that have $j$-gauge at most~$1$ at~$1$.}
\leanline{background library: \lnm{Space.finitePart}, \lnm{Space.finHeisOp}, \lnm{Space.heisFinPart},
\lnm{Space.projOn}, \lnm{Kernel}, \lnm{Kernel.lift}, \lnm{Kernel.lift\_equivariant}, \lnm{thetaKernel}\sep unit 141:
\lnm{chain}, \lnm{chain₂}, \lnm{projOnFinPart}\sep unit 147: \lnm{Data.homX}, \lnm{Data.liftKernel},
\lnm{Data.Translations}, \lnm{Data.translateDatum}, \lnm{Data.chainComb}, \lnm{Data.Bμ}, \lnm{Data.integralColumns}\sep
unit 150: \lnm{Parameters.intVacuum}, \lnm{Kernel.raises}, \lnm{Kernel.smul\_mem\_raises}}

\statement{inksoft}{\textbf{Input I5 (induced primes and marked sets; unit 147).}\enspace In Settings \ref{set:D}
and~\ref{set:p}, for $\tau\in\PlD{D}$ the prime $\ind{\tau}$ of $\ED{D}$ \emph{induced} by~$\tau$ is the set of
the integers $y$ of $\ED{D}$ with $|j(\emb{\tau}(y))|_p<1$; it lies above~$p$. A finite set $\mathcal{L}$ of finite
places of $\ED{D}$ is \emph{marked} (for $D$, $p$ and~$j$) if \textup{(M1)} $\cj{\mathfrak{l}}\notin\mathcal{L}$ for
every $\mathfrak{l}\in\mathcal{L}$; \textup{(M2)} if $\mathcal{L}$ is not empty, then $p$ is served for~$D$;
\textup{(M3)} every $\mathfrak{l}\in\mathcal{L}$ is $\ind{\tau}$ for some $\tau\in\tops{D}$; \textup{(M4)} $p$ is
small and not split for~$D$. It is a \emph{full mark} if moreover, for every $\tau\in\tops{D}$, one of
$\ind{\tau}$, $\cj{\ind{\tau}}$ lies in~$\mathcal{L}$. For a finite set $\mathcal{L}$ of finite places put
\[
\even{\mathcal{L}}=\{\tau\in\PlD{D}:\ \ind{\tau}\in\mathcal{L}\},\qquad
\mexp{D}{\mathfrak{l}}=\aD{D}\cdot\#\{\tau\in\tops{D}:\ \ind{\tau}=\mathfrak{l}\},
\]
and for a natural number~$i$ and $\tau\in\PlD{D}$
\[\lc{D}{i}{\tau}=\sum_{h<\nD{D}{i+1},\,h\ne1}\kap{D}{\tau}{h},\qquad\lb{D}{i}{\tau}=\sum_{h<\nD{D}{i+1},\,h\ne1}h\,\kap{D}{\tau}{h}.\]}
\leanline{unit 147: \lnm{Data.inducedPlace}, \lnm{Data.IsMarked}, \lnm{Data.IsFullMark}, \lnm{Data.evenPlaces},
\lnm{Data.markExponent}, \lnm{Data.lowColumns}, \lnm{Data.lowBoxes}}

\statement{inksoft}{\textbf{Input I6 (the bookkeeping of the unit on the CM curve; unit 150).}\enspace In
Setting~\ref{set:D}, for natural numbers $p$, $h$, $i$, $e$, a real number $y$, places $\tau,\tau'\in\PlD{D}$ and
a finite set $\Xi\subseteq\PlD{D}$; a pair $(e,p)$ is \emph{live} if $e\ge1$ and ($e$ is odd or $p=d$):
\begin{gather*}
\ask{D}{p}{y}{\tau}=\frac{1-\chiK{P}(p)}{2}\sum_{e:\ (e,p)\text{ live}}\Bigl(\frac{y}{p^{e}}-\nomW{D}{\pin{D}{-p^{e}}\cdot\tau}\Bigr)_{+},\\
\forced{D}{p}=\sum_{\tau\in\PlD{D}}\ask{D}{p}{\nomW{D}{\tau}}{\tau},\\
\gain{D}{p}=m\Bigl(\sum_{\tau\in\tops{D}}\ask{D}{p}{\aD{D}}{\tau}+\kD{D}\,\forced{D}{p}\Bigr),\\
\loss{D}{p}=\fD{D}\,(\log b)^2\,\frac{\aD{D}}{p}\cdot m^{\ast},
\end{gather*}
where $(\,\cdot\,)_+$ is the positive part, and $m^\ast=m$ if $p$ is short for~$D$, $m^\ast=1$ otherwise; $p$ \emph{is short for $D$} if $m$ divides $p-1$ (truncated) or
$p+1$; $p$ \emph{is small and not split for $D$} if $p<x_1$ and $\chiK{P}(p)\ne1$; $\rr{D}{p}$ is the least
natural number with $m\le p^{\rr{D}{p}}$ (for $p>1$);
$\DeltaP{D}{p}=\rr{D}{p}\,\fD{D}\,(\nD{D}{m}+\rr{D}{p})$; $\jc{D}{p}=\lceil2m\fD{D}\aD{D}/p\rceil$;
$\oz{D}{p}$ is the order of the class modulo $N_0$ of the place $\pin{D}{-p}$; $p$ \emph{is served for $D$} if
$\chiK{P}(p)=-1$, $p$ is not short for~$D$, and $\oz{D}{p}$ is odd and at most $\rr{D}{p}$. The pair $(\tau,e)$
\emph{is a collision of $p$ outside $\Xi$ for $D$} if $(e,p)$ is live and both $\tau$ and
$\pin{D}{-p^{e}}\cdot\tau$ lie in $\tops{D}\setminus\Xi$, and $\coll{D}{p}{\Xi}$ is the sum of the numbers
$\aD{D}/p^{e}$ over all the pairs $(\tau,e)$ that are collisions of $p$ outside $\Xi$ for~$D$.
For column counts $\varkappa$ on a set of places and a natural number~$n'$, with
$\mathrm{one}_\tau=[\tau\notin\Xi]\,[1<n']\,\varkappa(\tau,1)$, $\mathrm{low}_\tau=\sum_{h<n',h\ne1}\varkappa(\tau,h)$ and
$\mathrm{bx}_\tau=\sum_{h<n',h\ne1}h\,\varkappa(\tau,h)$,
\[\cc{p}{\varkappa}{n'}{\Xi}{\tau}{\tau'}{e}=\min\Bigl(\frac{\mathrm{one}_\tau}{p^{e}},\,\mathrm{low}_{\tau'}+\mathrm{one}_{\tau'}\Bigr)+\frac{\mathrm{low}_\tau}{p^{e}}+\min\Bigl(\frac{\mathrm{bx}_\tau}{p^{e}},\,\mathrm{bx}_{\tau'}\Bigr).\]
The column counts of the data are $\kap{D}{\tau}{1}=[\tau\in\tops{D}]\,\aD{D}$,
$\kap{D}{\tau}{\kD{D}}=\wtD{D}{1}{\tau}-\lfloor(\kD{D}+\tD{D}+1)/2\rfloor$ (truncated; for $\kD{D}\ne1$),
$\kap{D}{\tau}{h}=\gap{D}-(\tD{D}+1)$ for $h>\kD{D}$ with $h\equiv\kD{D}$ modulo $\tD{D}+1$, and
$\kap{D}{\tau}{h}=0$ otherwise, where $\gap{D}=2r^{8}$; $\last{D}{0}{\tau}=0$ and
$\last{D}{i}{\tau}=\lfloor\nD{D}{i-1}/2\rfloor-\wtD{D}{i}{\tau}$ for $i\ge1$;
$\kept{D}{h}=r^{8}-\lfloor(\tD{D}+1)/2\rfloor$ (truncated) if $h\ge\kD{D}$ and $h\equiv\kD{D}$ modulo
$\tD{D}+1$, and $0$ otherwise; $\kapY{D}{\tau}{h}=\kept{D}{h}+(\kap{D}{\tau}{h+1}-\kept{D}{h+1})$. Finally
$\nsw{d}{p}$ is $1/(p^2-1)$ if $\chiK{P}(p)=-1$, $1/(2(p-1))$ if $p=d$, and $1/(p-1)$ otherwise; for a prime $p$
with $\chiK{P}(p)\ne1$ and a real number~$c$,
\begin{equation}\label{e:live}
\sum_{e:\ (e,p)\text{ live}}\frac{c}{p^{e}}=\frac{2}{1-\chiK{P}(p)}\;c\,\Bigl(\frac{1}{p-1}-\nsw{d}{p}\Bigr),
\end{equation}
and the series defining $\ask{D}{p}{y}{\tau}$ converges.}
\leanline{unit 150: \lnm{Data.localAsk}, \lnm{Data.forcedPart}, \lnm{Data.gain}, \lnm{Data.loss},
\lnm{Data.IsShortPrime}, \lnm{Data.IsSmallNonsplit}, \lnm{Data.ceilLogM}, \lnm{Data.deltaP}, \lnm{Data.jetCeil},
\lnm{Data.o₀}, \lnm{Data.IsServed}, \lnm{Data.IsCollision}, \lnm{Data.collisionSum}, \lnm{couplingCount},
\lnm{Data.κ}, \lnm{Data.lastEntry}, \lnm{Data.keptColumns}, \lnm{Data.columnsY}, \lnm{hasSum\_live\_inv\_pow},
\lnm{Data.summable\_localAsk}\sep background library: \lnm{nonsplitWeight}}

\subhead{The definitions of this unit.}

\begin{definition}[the primes dividing $\piN{D}$; levels, columns, stages]\label{d:pi}
In Setting~\ref{set:D}, let $i$ be a natural number.
\begin{enumerate}[label=(\alph*)]
\item $\piPl{D}$ is the finite set of the finite places $\mathfrak{p}$ of $\ED{D}$ with $\piN{D}\in\mathfrak{p}$.
\item The \emph{column level} $\colL{D}{i}$ is the group of the elements $k$ of the chosen level
$\Kch{D}{\XD{D}{i}}$ whose last row is congruent to $(0,\dots,0,1)$ modulo~$\piN{D}$: every entry of the last row
of $k-1$ lies in $\piN{D}\,\widehat{\OO}_{\ED{D}}$.
\item The space of \emph{columns} $\Col{D}$ is the space of the forms in
$\Alg{\XD{D}{0}}{\kapD{D}}{\lastD{D}{0}}$ that are fixed under right translation by $\colL{D}{0}$.
\item The \emph{stage} $\Stg{D}{i}$ is the space $\Alg{\XD{D}{i}}{\kapD{D}}{\lastD{D}{i}}$.
\end{enumerate}
\end{definition}
\leanline{\lean{piPlaces}\sep\lean{columnLevel}\sep\lean{Column}\sep\lean{Stage}}

\begin{definition}[branching and the restriction]\label{d:branch}
In Setting~\ref{set:D}, let $\tau\in\PlD{D}$.
\begin{enumerate}[label=(\alph*)]
\item $\lstEq{D}$ is a chosen bijection from $\YD{D}{m}$ onto the set of the elements of $\XD{D}{m}$ other than
the last; both sets have $\nD{D}{m}-1$ elements.
\item $\bdU{D}\colon\Gf{\YD{D}{m}}\to\Gf{\XD{D}{m}}$ is the homomorphism $h\mapsto\diag(h,1)$, the block $h$ being
placed through $\lstEq{D}$ and the $1$ at the last index.
\item The \emph{branching exponent} is the integer
$\brexp{D}{\tau}=\sum_{0<h<\nD{D}{m}}\bigl(\kap{D}{\tau}{h}-\kept{D}{h}\bigr)+\last{D}{m}{\tau}$.
\item $\Brm{D}$ is a chosen matrix with integer entries, with rows indexed by $\Wy{\YD{D}{m}}{\kapYD{D}}$ and
columns by $\Wy{\XD{D}{m}}{\kapD{D}}$, that has the properties provided by the branching theorem of the
background library (branching from $\GL_{n+1}$ to $\GL_n$ with multiplicity one): \textup{(Br1)} over every
field of characteristic~$0$, for families $(g_\tau)$ of invertible matrices indexed by $\YD{D}{m}$ and of units
$(s_\tau)$,
\[\Brm{D}\cdot\rhow{\kapD{D}}{\lastD{D}{m}}\bigl((\diag(g_\tau,s_\tau))_\tau\bigr)=\prod_{\tau}s_\tau^{\,\brexp{D}{\tau}}\cdot\rhow{\kapYD{D}}{\lastD{D}{m}}\bigl((g_\tau)_\tau\bigr)\cdot\Brm{D};\]
\textup{(Br2)} $\wn{\Brm{D}\xi}\le L^{6L}\,\wn{\xi}$ for every vector $\xi$ and every natural number
$L\ge\sum_\tau\sum_{\kD{D}\le h<\nD{D}{m}-1}h\,\kapY{D}{\tau}{h}$; \textup{(Br3)} a spanning property of the
library, which this unit does not use.
\item The \emph{restriction} $\resY{D}\colon\Stg{D}{m}\to\Alg{\YD{D}{m}}{\kapYD{D}}{\lastD{D}{m}}$ is the
$\CC$-linear map
\[
\bigl(\resY{D}\varphi\bigr)(h)=\prod_{\tau'\in\PlD{D}}\emb{\tau'}\bigl(\cj{\piN{D}}\bigr)^{-\brexp{D}{\tau'}}\cdot\Brm{D}\,\varphi\bigl(\bdU{D}(h)\bigr)\qquad(h\in\Gf{\YD{D}{m}}).
\]
\end{enumerate}
\end{definition}
\leanline{\lean{lastIndexEquiv}\sep\lean{borderU}\sep\lean{branchExp}\sep\lean{branching}\sep\lean{restrictToY}\sep
the branching theorem is \lnm{exists\_branching\_matrix} of the background library; the catalogue of the handout
records it as a known result (branching from $\GL(n+1)$ to $\GL(n)$ with multiplicity one), \emph{reference to be
confirmed}}

That $\resY{D}\varphi$ is again an algebraic form follows from (Br1), applied to a rational point $\gamma$ of
$\Gf{\YD{D}{m}}$ and the units $s_\tau=1$; the formal definition contains this verification.

\begin{definition}[short forms]\label{d:short}
Let $E$ be a CM field with a CM type, $I$ a finite set, $\varkappa$ column counts and $\mathfrak{e}$ determinant
exponents. For a real number $B$ and a form $\varphi\in\Alg{I}{\varkappa}{\mathfrak{e}}$ we say that
$\short{B}{\varphi}$ if
\begin{enumerate}[label=(\roman*)]
\item for every prime number $p$ and every ring homomorphism $j\colon\CC\to\CC_p$, $\varphi$ has $j$-gauge at
most~$1$ at~$1$; and
\item $\wn{\sigma\circ\varphi(g)}\le B$ for every $\sigma\in\Aut$ and every $g\in\Gf{I}$.
\end{enumerate}
\end{definition}
\leanline{\lean{IsShort}}

\begin{lemma}[a central element at a prime dividing $\piN{D}$]\label{l:cgen}
In Setting~\ref{set:D}, let $\mathfrak{p}\in\piPl{D}$. There is a finite adele $c$ of $\ED{D}$ with
$c\,\cj{c}=1$ such that: $c$ is integral; $c_w=1$ at every finite place $w$ with $w\ne\mathfrak{p}$ and
$\cj{w}\ne\mathfrak{p}$; and every natural number $n$ for which $c_{\mathfrak{p}}^{\,n}-1$ has
$\mathfrak{p}$-adic absolute value less than~$1$ is divisible by $\ND{D}-1$.
\end{lemma}
\leanline{\lean{existscentralGen}\sep proof in §\ref{s:short}}

\begin{definition}[central translations and ramified columns]\label{d:central}
\begin{enumerate}[label=(\alph*)]
\item Let $E$ be a CM field with a CM type, $I$ a finite set, $\varkappa$ column counts, $\mathfrak{e}$
determinant exponents, $J$ a level of~$\Gf{I}$ and $\chi$ a character of~$J$. The \emph{central representation}
is the representation $z\mapsto\mathrm{R}(z\cdot1_I)$, by right translation by scalar matrices, of the group of
the finite adeles $z$ of $E$ with $z\cj{z}=1$ on the space of the forms
$\varphi\in\Alg{I}{\varkappa}{\mathfrak{e}}$ with $\mathrm{R}(k)\varphi=\chi(k)\varphi$ for all $k\in J$. (Scalar
matrices commute with~$J$, so this space is stable.)
\item In Setting~\ref{set:D}, for $\mathfrak{p}\in\piPl{D}$, $\cgen{D}{\mathfrak{p}}$ is a chosen adele $c$ with
the properties of Lemma~\ref{l:cgen}.
\item In Setting~\ref{set:D}, $\cenrep{D}$ is the central representation of (a) on $\Col{D}$, that is, for
$E=\ED{D}$, $I=\XD{D}{0}$, the weight $(\kapD{D},\lastD{D}{0})$, $J=\colL{D}{0}$ and $\chi=1$.
\item In Setting~\ref{set:D}, for a finite set $\Sigma\subseteq\piPl{D}$, the endomorphism
$\ramop{D}{\Sigma}=\prod_{\mathfrak{p}\in\Sigma}\bigl(1-\cenrep{D}(\cgen{D}{\mathfrak{p}})\bigr)$ of\/ $\Col{D}$; the
factors commute.
\item In Setting~\ref{set:D}, the space of \emph{ramified columns} $\Colram{D}$ is the image of
$\ramop{D}{\piPl{D}}$.
\end{enumerate}
\end{definition}
\leanline{\lean{algFormCentralRep}\sep\lean{centralGen}\sep\lean{centralRep}\sep\lean{ramifyingOp}\sep\lean{ramifiedColumns}}

\begin{definition}[admissible translations, paired translates, good data]\label{d:good}
In Setting~\ref{set:D}:
\begin{enumerate}[label=(\alph*)]
\item For a natural number $q$, $\admT{q}$ is the set of the tuples of translations $\mathsf{x}$ such that, for
every~$i$, every entry of $\pig{D}\cdot\mathsf{x}_i$ is integral at every finite place $w$ with $q\notin w$.
\item For a family $v=(v_i)_{i\in\NN}$ with $v_i\in\Wfin{D}{i}$ and a natural number~$\nu$,
\[\pairT{v}{\nu}=\sum_{y}\mathrm{U}\bigl(\mathrm{col}(-y/N)\bigr)v_{2\nu}\otimes\mathrm{U}\bigl(\mathrm{row}(-\cj{y/N})\bigr)v_{2\nu+1}\ \in\ \Wfin{D}{2\nu}\otimes\Wfin{D}{2\nu+1},\]
where $N=\ND{D}$; $y$ runs over the vectors indexed by $\XD{D}{2\nu+1}$ with entries in
$(\Z/N)^{\piPl{D}}$; $y/N$ is the vector of finite adeles whose component at $\mathfrak{p}\in\piPl{D}$ is the least
non-negative representative of $y_{\mathfrak{p}}$ divided by~$N$, and $0$ at the other places;
$\mathrm{col}(z)$ is the matrix with rows $\XD{D}{2\nu+1}$ and columns $\XD{D}{2\nu}$ whose last column is $z$ and
whose other columns are~$0$; and $\mathrm{row}(z)$ is the matrix with rows $\XD{D}{2\nu+2}$ and columns
$\XD{D}{2\nu+1}$ whose last row is $z$ and whose other rows are~$0$.
\item Let $\vartheta$ be a twist of~$D$ and $\Psi$ a finitely supported complex function on the tuples of
translations of~$D$. We say that $\good{\vartheta}{\Psi}$ if the following three conditions hold.

\textup{(G1)} The support of $\Psi$ lies in $\admT{\nrm{\vartheta}\cdot\ND{D}}$; for every $\mathsf{x}$, every
prime number $\ell$ prime to $\nrm{\vartheta}$ and every ring homomorphism $\iota\colon\CC\to\CC_\ell$,
$|\iota(\Psi(\mathsf{x}))|_\ell\le1$; and if $Q(\vartheta)=\emptyset$, then
$\sum_{\mathsf{x}}|\sigma(\Psi(\mathsf{x}))|\le\exp(x^{60})$ for every $\sigma\in\Aut$.

\textup{(G2)} (\emph{joint invariance}) For every family $(a_i)_{i\in\NN}$ with $a_i\in\colL{D}{i}$ and $a_i=1$
for $i\ge m$, the tensor
$\sum_{\mathsf{x}}\Psi(\mathsf{x})\cdot\bigotimes_{i<m}\td{D}{\mathsf{x}}{i}$ in $\bigotimes_{i<m}\Wfin{D}{i}$ is
fixed by $\bigotimes_{i<m}\lkom{D}{\vartheta}{i}(a_i,a_{i+1})$; in the formal text: the two tensors have the same
image under every linear map.

\textup{(G3)} (\emph{the shape at $N$}) For every form $\varphi\in\Stg{D}{0}$,
\[
\sum_{\mathsf{x}}\Psi(\mathsf{x})\cdot\chn{D}{\vartheta}{\td{D}{\mathsf{x}}{\,\cdot\,}}{m}\,\varphi
=\sum_{\mathsf{x}}\Psi(\mathsf{x})\cdot\chnn{D}{\vartheta}{\nu\mapsto\pairT{\mathrm{P}^{\piPl{D}}\td{D}{\mathsf{x}}{\,\cdot\,}}{\nu}}{m/2}\,\varphi .
\]
\end{enumerate}
\end{definition}
\leanline{\lean{admissibleTranslations}\sep\lean{pairTranslates}\sep\lean{IsGoodDatum}\sep with
\lnm{IsAdmissibleCoeffs}, \lnm{IsJointlyInvariant}, \lnm{translateSum} of unit 141}

In (G3) the formal text writes the number of pairs as $\lfloor m/2\rfloor-1+1$ and the number of lifts on the
left as twice that; $m$ is a power of~$2$ with $m\ge2$, so these are $m/2$ and~$m$. That good data exist, and
that the rows of the handout are of the form $\Bch{D}{\vartheta}{\Psi}\psi$ for good~$\Psi$, is proved in the
unit on the seesaw; this unit only uses the three conditions.

\begin{definition}[the jet exponent, the loss of a lift, the mark operators]\label{d:mark}
In Setting~\ref{set:D}:
\begin{enumerate}[label=(\alph*)]
\item For natural numbers $p$ and $i$ and a finite set $\Xi\subseteq\PlD{D}$, put $n'=\nD{D}{i+1}$,
$L=\boxes{\kapD{D}}{n'}$ and $l=\lfloor\log_pL\rfloor$ (which is $0$ if $p\le1$ or $L=0$). The \emph{jet
exponent} is the real number
\begin{multline*}
\jetexp{p}{i}{\Xi}=\Bigl(\frac{1}{p-1}-\nsw{d}{p}\Bigr)L-\fD{D}\,n'^{\,3}\,(l+1)\\
+\frac{\chiK{P}(p)-1}{2}\sum_{\substack{0\le e\le l\\ e\text{ odd or }p=d}}\ \sum_{\tau\in\PlD{D}}\cc{p}{\kapD{D}}{n'}{\Xi}{\tau}{\pin{D}{-p^{e}}\cdot\tau}{e}.
\end{multline*}
\item In Setting~\ref{set:p}, for a finite set $\mathcal{L}$ of finite places of $\ED{D}$ and a natural
number~$i$, the \emph{loss of the lift $i$} is
$\liftloss{\mathcal{L}}{i}=\rr{D}{p}\sum_{\tau\in\even{\mathcal{L}}}\lc{D}{i}{\tau}-\jetexp{p}{i}{\even{\mathcal{L}}}$.
\item In Setting~\ref{set:p}, for a marked set $\mathcal{L}$ and $\mathfrak{l}\in\mathcal{L}$, the \emph{mark
operator} $\markop{\mathcal{L}}{\mathfrak{l}}$ is the $\Oj$-linear endomorphism
$\varphi\mapsto p^{\mexp{D}{\mathfrak{l}}}\cdot\Tl{\mathfrak{l}}{\rr{D}{p}}\varphi$ of\/ $\Col{D}$, where
$\Tl{\mathfrak{l}}{s}$ is the Hecke operator of the double coset of $\diag(\varpi1_s,1)$ at the split
place~$\mathfrak{l}$ (Input~I8); columns are fixed by $\Kmax{\XD{D}{0}}{\mathfrak{l}}$, and the operator preserves
$\Col{D}$ (Lemma~\ref{l:markprop}).
\end{enumerate}
\end{definition}
\leanline{\lean{jetExponent}\sep\lean{liftLoss}\sep\lean{markOperator}}

For $p=d$ the sum over $e$ in (a) contains $e=0$; the
formal text has it so.

\begin{definition}[the height per row, the non-split primes, the allowed loss]\label{d:size}
In Setting~\ref{set:D}:
\begin{enumerate}[label=(\alph*)]
\item $\hrow{D}=m\Bigl(\sum_{\tau\in\tops{D}}\bigl(\jsz{D}{\aD{D}}-\Ask{D}{\aD{D}}{\tau}\bigr)+\kD{D}\,\dirc{D}\Bigr)-\#\tops{D}\cdot\aD{D}\cdot\log\ND{D}$,
where for a natural number~$n$, $\jsz{D}{n}=\log n!-n\log\bigl(\OmP{P}{\Dmu{D}}/\gD{D}\bigr)$, for a real $y$ and a
place~$\tau$,
$\Ask{D}{y}{\tau}=\sum_{\nu=1}^{\aD{D}}\lamd{d}{1}{\nu}\bigl(y/\nu-\nomW{D}{\pin{D}{-\nu}\cdot\tau}\bigr)_{+}$, and
$\dirc{D}=\sum_{\tau\in\PlD{D}}\bigl(\jsz{D}{\wtD{D}{0}{\tau}}-\Ask{D}{\nomW{D}{\tau}}{\tau}\bigr)$ is the cost of
one direction (objects of unit~152 and of the overview).
\item $\nsp{D}$ is the finite set of the prime numbers $p<\lceil x_1\rceil$ with $\chiK{P}(p)\ne1$.
\item For a natural number~$p$, the \emph{allowed loss} $\err{D}{p}$ is $3\,\loss{D}{p}$ if $p\in\nsp{D}$;
$4\,m\,\fD{D}\,b$ if $p\notin\nsp{D}$ and $p=\ND{D}$; and $0$ otherwise.
\end{enumerate}
\end{definition}
\leanline{\lean{heightPerRow}\sep\lean{nonsplitPrimes}\sep\lean{allowedLoss}\sep unit 152: \lnm{Data.jetSize},
\lnm{Data.primeSum}, \lnm{Data.dirCost}}

\sect{Levels and columns}{s:levels}
\begin{lemma}[the primes dividing $\piN{D}$]\label{l:pi}
In Setting~\ref{set:D}:
\begin{enumerate}[label=(\alph*)]
\item $\piN{D}\ne0$.
\item A finite place $w$ of $\ED{D}$ lies in $\piPl{D}$ if and only if $\piN{D}\in w$.
\item For every $\mathfrak{p}\in\piPl{D}$: $\cj{\mathfrak{p}}\notin\piPl{D}$, the residue field of $\mathfrak{p}$ has
$\ND{D}$ elements, and $\ND{D}$ is a uniformiser at~$\mathfrak{p}$. For every finite place $w$ of $\ED{D}$:
$\ND{D}\in w$ if and only if $w\in\piPl{D}$ or $\cj{w}\in\piPl{D}$.
\item $\#\piPl{D}=\fD{D}$.
\end{enumerate}
\end{lemma}
\begin{proof}
(a) and (b) follow from $\piN{D}\cj{\piN{D}}=\ND{D}\ne0$ and the definition. (c) Since $\ND{D}\equiv1$ modulo
$\NnD{D}\,d$, the prime $\ND{D}$ splits completely in $\ED{D}$ (unit~152): every place above it has residue
field of $\ND{D}$ elements and is different from its conjugate, and a place above it contains exactly one of
$\piN{D}$, $\cj{\piN{D}}$. (d) The same input gives a bijection between $\PlD{D}$ and the primes that contain
$\piN{D}$.
\end{proof}
\leanline{\lean{nezero}\sep\lean{mempiPlaces}\sep\lean{isSplitLevelpiPlaces}\sep\lean{cardpiPlaces}\sep inputs of
unit 152: \lnm{Parameters.q\_eq\_of\_modEq\_one}, \lnm{Parameters.star\_mem\_iff},
\lnm{Parameters.nonempty\_infinitePlace\_equiv\_finPlace}}

\begin{lemma}[the column levels]\label{l:level}
In Setting~\ref{set:D}, let $i$ be a natural number.
\begin{enumerate}[label=(\alph*)]
\item $\Kprin{\XD{D}{i}}{\cj{\pig{D}}\,\piN{D}}\subseteq\colL{D}{i}$.
\item $\colL{D}{i}$ is a level: an open subgroup of\/ $\Kf{\XD{D}{i}}$.
\item $\XD{D}{i}$ has $\nD{D}{i}$ elements.
\item $\colL{D}{0}$ is free: $\Gam{\XD{D}{0}}\cap g\,\colL{D}{0}\,g^{-1}=\{1\}$ for every $g\in\Gf{\XD{D}{0}}$.
\item An element $k$ of\/ $\Gf{\XD{D}{i}}$ lies in $\colL{D}{i}$ if and only if $k\in\Kch{D}{\XD{D}{i}}$ and, for
every $\mathfrak{p}\in\piPl{D}$ and every index~$c$, the entry of $k_{\mathfrak{p}}-1$ in the last row and the
column~$c$ lies in the maximal ideal at~$\mathfrak{p}$.
\item If $w$ is a finite place of $\ED{D}$ with $\ND{D}\in w$, then neither $w$ nor $\cj{w}$ is chosen.
\item If $S$ is a set of finite places of $\ED{D}$ such that $\ND{D}\in w$ for every $w\in S$, then every element
of\/ $\Kf{\XD{D}{i}}$ whose components away from the places of $S$ and their conjugates are~$1$ lies in
$\Kch{D}{\XD{D}{i}}$.
\item If $\ell$ is a prime number with $\ell\ne\ND{D}$ and $w$ a finite place of $\ED{D}$ with $\ell\in w$, then
$\ND{D}\notin w$ and $\cj{w}\notin\piPl{D}$.
\item In Setting~\ref{set:p}, $|j(\cj{\pig{D}})|_p\le1$.
\end{enumerate}
\end{lemma}
\begin{proof}
(a) An element of the principal level of $\cj{\pig{D}}\piN{D}$ is congruent to~$1$ modulo the chosen places and
modulo~$\piN{D}$. (b) $\colL{D}{i}$ lies in $\Kf{\XD{D}{i}}$ and contains the open group of~(a). (c) is the
definition of the index sets. (d) $\colL{D}{0}$ is a subgroup of the chosen level and
$\nD{D}{0}\le2\nD{D}{m}$ (Input~I3). (e) unfolds Definition~\ref{d:pi}(b) place by place: by
Lemma~\ref{l:pi}(c), $\piN{D}$ is a uniformiser at the places of $\piPl{D}$ and a unit elsewhere. (f) A chosen
place lies above $g_0$ or $g_1$, and these primes differ from $\ND{D}$ by size. (g) follows from (f): such an
element is $1$ at every chosen place. (h) A prime ideal contains only one prime number; the second claim is
Lemma~\ref{l:pi}(c). (i) $\pig{D}$ is an algebraic integer, and ring homomorphisms preserve integrality.
\end{proof}
\leanline{\lean{principalLevellecolumnLevel}\sep\lean{isLevelcolumnLevel}\sep\lean{cardX}\sep\lean{classAutcolumnLeveleqbot}\sep
\lean{memcolumnLevel}\sep\lean{notisChosenofNmem}\sep\lean{memchosenLevelofsupportedOn}\sep\lean{notMempiPlacesofne}\sep\lean{normstarpiGleone}}

\begin{lemma}[two bounds for the weight]\label{l:branch}
In Setting~\ref{set:D}, let $\tau\in\PlD{D}$.
\begin{enumerate}[label=(\alph*)]
\item $\kap{D}{\tau}{1}-2b\le\brexp{D}{\tau}$, as real numbers.
\item $\sum_{h<\#\XD{D}{0}}h\,\kap{D}{\tau}{h}\le\aD{D}$.
\end{enumerate}
\end{lemma}
\begin{proof}
(a) In $\brexp{D}{\tau}$ the term $h=1$ of the sum is $\kap{D}{\tau}{1}$, the other terms are non-negative, and
$\last{D}{m}{\tau}\ge-\wtD{D}{m}{\tau}\ge-2b$ by the bound $\wtD{D}{i}{\tau}\le2b$ for $i\le2m$ of unit~152.
(b) $\XD{D}{0}$ has $\kD{D}$ elements, and among the heights $h<\kD{D}$ only $h=1$ has columns, namely
$\aD{D}$ at the tops and none elsewhere.
\end{proof}
\leanline{\lean{lebranchExp}\sep\lean{sumheightmulkappale}\sep input of unit 152: \lnm{Data.w\_le\_two\_mul\_b}}

\sect{Short forms span}{s:short}
\begin{lemma}[differences of translates]\label{l:sub}
Let $E$ be a CM field with a CM type, $I$ a finite set, $\varkappa$ column counts and
$\mathfrak{e}$ determinant exponents. Let $B$ be a real number, $k\in\Kf{I}$ and
$\varphi\in\Alg{I}{\varkappa}{\mathfrak{e}}$. If $\short{B}{\varphi}$, then
$\short{4B}{(\varphi-\mathrm{R}(k)\varphi)}$.
\end{lemma}
\begin{proof}
At a $p$-adic place the forms of gauge at most~$1$ at~$1$ form a module, and the translate
$\mathrm{R}(k)\varphi$ is one of them: its value at $g$ is $\varphi(gk)$, and the gauge of a vector at $g$ equals
its gauge at $gk$, because the gauge is invariant under $\Kf{I}$ (Input~I1). For $\sigma\in\Aut$ and a point~$g$ put
$\xi=\sigma\circ\varphi(g)$, $\xi'=\sigma\circ\varphi(gk)$; by the parallelogram law for the Gram form,
$\wn{\xi-\xi'}\le\wn{\xi-\xi'}+\wn{\xi+\xi'}=2\wn{\xi}+2\wn{\xi'}\le4B$.
\end{proof}
\leanline{\lean{subtranslate}}

\statement{inksoft}{\textbf{Input I7 (short bases, delta forms, Schur matrices; background library).}\enspace Let $E$
be a CM field with $f$ infinite places, $I$ a finite set and $N'\ge1$ a natural number that lies in the different of
$E$ over~$\Q$. \textup{(i)} For every $g\in\Gf{I}$ there is an invertible matrix $Y(g)$ over $E$, indexed
by~$I$, such that $g^{-1}Y(g)$ has entries in $\widehat{\OO}_E$ and every entry of $Y(g)$ has absolute value at
most $6N'(f|I|)^2$ at every embedding $E\to\CC$. \textup{(ii)} For a free level~$K$ the delta form with value
$v$ at the representative $g_c$ of a class has, at every point, $j$-gauge at most
$\gauge{\varkappa}{\mathfrak{e}}{g_c}{v}$, and its values $\xi$ satisfy
$\wn{\sigma\circ\xi}\le\wn{\sigma\circ v}$ for $\sigma\in\Aut$. \textup{(iii)} The \emph{evaluation vectors}
$\mathrm{ev}(y)\in\CC^{\Wy{I}{\varkappa}}$, the values at a point $y$ of the basis polynomials of the weight
$\varkappa$, span $\CC^{\Wy{I}{\varkappa}}$ when $y$ runs over the points whose coordinates are roots of unity; they
have integer polynomial coordinates, and
$\rhow{\varkappa}{0}(z)\,\mathrm{ev}(y)$ is an evaluation vector at the transformed point. \textup{(iv)} For
families $z=(z_\tau)$ of complex matrices indexed by~$I$ there are real numbers $c_{\tau,\nu}\ge0$, with
$\sum_{\nu}c_{\tau,\nu}$ at most the sum of the squares of the absolute values of the entries of $z_\tau$,
such that the Schur matrix of $z$ multiplies $\wn{\,\cdot\,}$ by at most
$\prod_{\tau}\prod_{\nu<|I|}c_{\tau,\nu}^{\,l(\tau,\nu)}$, where $l(\tau,\nu)$ is the length of the row of
index~$\nu$ at~$\tau$.}
\leanline{background library: \lnm{Space.exists\_short\_basis\_stdSpace}, \lnm{weylGauge\_deltaForm\_le},
\lnm{weylNormSq\_deltaForm\_le}, \lnm{PolyRep.ext\_weylEval}, \lnm{weylNormSq\_schur\_le}; the catalogue classes
them as elementary and gives no reference for them: \emph{reference to be confirmed}; the header of the file of
the last one names the contravariant form on Weyl's lattice}

\begin{proposition}[short forms span the forms of a free level]\label{p:span}
Let $E$ be a CM field with a CM type, $I$ a non-empty finite set, $\varkappa$ column counts and
$\mathfrak{e}$ determinant exponents. Let $K$ be a level of\/ $\Gf{I}$ that is free, $N'\ge1$ a natural number that
lies in the different of the ring of integers of $E$ over~$\Z$, $M$ a natural number with
\[|I|\cdot\#\operatorname{Pl}(E)+\boxes{\varkappa}{|I|}\ \le\ M ,\]
and $B$ a real number with $\bigl(6N'(\#\operatorname{Pl}(E)\cdot|I|)^2\bigr)^{6M}\le B$. Then the space of the
$K$-fixed forms in $\Alg{I}{\varkappa}{\mathfrak{e}}$ is spanned over $\CC$ by its elements $\varphi$ with
$\short{B}{\varphi}$.
\end{proposition}
\begin{proof}
Put $R=6N'(f|I|)^2$ with $f=\#\operatorname{Pl}(E)$. By Input~I1 every $K$-fixed form is a sum of delta forms, so
it suffices to treat the delta forms at one class, with representative~$g$. Take the matrices $Y(\,\cdot\,)$ of
Input~I7(i). For a ring homomorphism $\phi$ from $\CC$ to a normed field, write $\phi_\tau$ for the embedding
$y\mapsto\phi(\emb{\tau}(\cj{y}))$ of~$E$ and put
\[A_\phi=\prod_{\tau}\det\phi_\tau\bigl(Y(g^{\mathfrak{e}(\tau)})\bigr)\cdot\rhow{\varkappa}{0}\bigl((\phi_\tau(Y(g)))_\tau\bigr).\]
The matrix $A=A_{\mathrm{id}}$ is invertible, so by Input~I7(iii) the vectors $A\,\mathrm{ev}(y)$, $y$ with
coordinates roots of unity, span $\CC^{\Wy{I}{\varkappa}}$; and
$\phi\circ(A\,\mathrm{ev}(y))=A_\phi\,\mathrm{ev}(\phi\circ y)$. It remains to show that the delta form with
value $A\,\mathrm{ev}(y)$ is $B$-short. \emph{At a $p$-adic place $j$:} by Input~I7(ii) its gauge is at most
$\gauge{\varkappa}{\mathfrak{e}}{g}{A\,\mathrm{ev}(y)}$, the largest absolute value of a coordinate of
\[\prod_\tau\det\bigl((g^{\mathfrak{e}(\tau)})_\tau^{-1}\,j_\tau Y(g^{\mathfrak{e}(\tau)})\bigr)\cdot\rhow{\varkappa}{0}\bigl((g_\tau^{-1}\,j_\tau Y(g))_\tau\bigr)\,\mathrm{ev}(j\circ y).\]
The matrices $g^{-1}Y(g)$ are integral, so the determinants and the Schur matrix have integral entries, and
$\mathrm{ev}(j\circ y)$ is integral because $j$ maps roots of unity to units. \emph{At $\sigma\in\Aut$:} by
Input~I7(ii) the values $\xi$ of the delta form have
$\wn{\sigma\circ\xi}\le\wn{A_\sigma\,\mathrm{ev}(\sigma\circ y)}$. Each determinant has absolute value at most
$(|I|R)^{|I|}\le R^{2|I|}$, the sums $\sum_\nu c_{\tau,\nu}$ of Input~I7(iv) are at most $|I|^2R^2\le R^4$,
the row lengths at all places add up to $\boxes{\varkappa}{|I|}$, and the evaluation vector at a point with
coordinates of absolute value~$1$ is bounded in terms of $|I|$ and the boxes. The formal text collects these
bounds into $R^{6M}$, using the hypothesis on~$M$; this is at most~$B$.
\end{proof}
\leanline{\lean{spanisShorteqtop}}

\begin{proof}[Proof of Lemma~\ref{l:cgen}]
By Lemma~\ref{l:pi}(c) the residue field at $\mathfrak{p}$ has $\ND{D}$ elements, so its group of units is cyclic
of order $\ND{D}-1$. Choose a unit $y$ of the ring of integers of the completion at~$\mathfrak{p}$ whose residue
generates this group, regard it as an idele with component $y$ at $\mathfrak{p}$ and $1$ elsewhere, and put
$c=y/\cj{y}$. Then $c\,\cj{c}=1$, $c$ is integral, and $c$ is $1$ away from $\mathfrak{p}$ and
$\cj{\mathfrak{p}}$. Since $\cj{\mathfrak{p}}\ne\mathfrak{p}$, the component of $c$ at $\mathfrak{p}$ is~$y$; if
$y^{\,n}-1$ lies in the maximal ideal, the residue of $y$ has order dividing~$n$, so $\ND{D}-1$ divides~$n$.
\end{proof}

\begin{lemma}[the central elements act on columns]\label{l:cgenmem}
In Setting~\ref{set:D}, let $\mathfrak{p}\in\piPl{D}$.
\begin{enumerate}[label=(\alph*)]
\item For every finite set $I$, the scalar matrix $\cgen{D}{\mathfrak{p}}\cdot1_I$ lies in
$\Kmax{I}{\mathfrak{p}}$.
\item For every $\varphi\in\Colram{D}$ there is a column $\varphi'\in\Col{D}$ with
$\varphi=\varphi'-\mathrm{R}\bigl(\cgen{D}{\mathfrak{p}}\cdot1_{\XD{D}{0}}\bigr)\varphi'$.
\end{enumerate}
\end{lemma}
\begin{proof}
(a) $\cgen{D}{\mathfrak{p}}$ is integral of norm one and is $1$ away from $\mathfrak{p}$, $\cj{\mathfrak{p}}$
(Lemma~\ref{l:cgen}). (b) $\varphi=\ramop{D}{\piPl{D}}\varphi_0$ for a column $\varphi_0$; split off the factor at
$\mathfrak{p}$ and take $\varphi'=\ramop{D}{\piPl{D}\setminus\{\mathfrak{p}\}}\varphi_0$.
\end{proof}
\leanline{\lean{scalarFincentralGenmem}\sep\lean{existseqsubcentralGen}}

\begin{lemma}[few columns are fixed by a central element]\label{l:ker}
In Setting~\ref{set:D}, for every $\mathfrak{p}\in\piPl{D}$,
\[\dim_{\CC}\ker\bigl(1-\cenrep{D}(\cgen{D}{\mathfrak{p}})\bigr)\ \le\ (\ND{D}-1)^{-1}\cdot\dim_{\CC}\Col{D}.\]
\end{lemma}
\begin{proof}
Let $z$ be the scalar matrix of $\cgen{D}{\mathfrak{p}}$ and $K$ the group generated by $\colL{D}{0}$ and~$z$. By
Lemma~\ref{l:cgenmem}(a) and Lemma~\ref{l:level}(g), $K$ lies in the chosen level; so $K$ is a level, and both
$K$ and $\colL{D}{0}$ are free (Input~I3). A column fixed by $z$ is fixed by~$K$. For free levels the dimension of
the space of fixed forms is the number of classes times the dimension of the weight (Input~I1), and each class
of $K$ consists of $[K:\colL{D}{0}]$ classes of $\colL{D}{0}$. Hence the kernel has dimension at most
$\dim\Col{D}/[K:\colL{D}{0}]$. Some power $z^{\,n}$ with $1\le n\le[K:\colL{D}{0}]$ lies in $\colL{D}{0}$; by
Lemma~\ref{l:level}(e) its component at $\mathfrak{p}$ is congruent to~$1$ modulo~$\mathfrak{p}$, so
$\ND{D}-1$ divides~$n$ by Lemma~\ref{l:cgen}. Therefore $[K:\colL{D}{0}]\ge\ND{D}-1$.
\end{proof}
\leanline{\lean{finrankkeronesubcentralGenle}}

\begin{proposition}[the ramifying operator]\label{p:ram}
In Setting~\ref{set:D}, let $\Sigma$ be a finite subset of $\piPl{D}$ and $W$ a complex subspace of\/ $\Col{D}$.
Then
\[\dim_{\CC}\ramop{D}{\Sigma}^{-1}(W)\ \le\ \dim_{\CC}W+\#\Sigma\cdot(\ND{D}-1)^{-1}\cdot\dim_{\CC}\Col{D},\]
and for every real number $B$ and every $\varphi\in\Col{D}$ with $\short{B}{\varphi}$,
\[\short{4^{\#\Sigma}B}{\ramop{D}{\Sigma}\varphi}.\]
\end{proposition}
\begin{proof}
Induction on $\Sigma$; for $\Sigma=\emptyset$ the operator is the identity. Adding a place $\mathfrak{p}$ composes
with $1-\cenrep{D}(\cgen{D}{\mathfrak{p}})$. The preimage of a subspace under a linear map has dimension at most
the dimension of the subspace plus that of the kernel, which Lemma~\ref{l:ker} bounds; and by
Lemmas \ref{l:cgenmem}(a) and~\ref{l:sub} the new factor multiplies the bound of shortness by~$4$.
\end{proof}
\leanline{\lean{ramifyingOpfinrankcomapandshort}}

For $W=0$ and $\Sigma=\piPl{D}$ the first claim says that the kernel of $\ramop{D}{\piPl{D}}$ has dimension at most
$\fD{D}\dim\Col{D}/(\ND{D}-1)$: the ramified columns have nearly full dimension. The unit on the filling (Theorem~B of the handout) uses the
proposition for its count of dimensions.

\begin{proof}[Proof of Theorem A]
Put $B'=\exp(\bud{D})/4^{\fD{D}}$. Apply Proposition~\ref{p:span} to $E=\ED{D}$, $I=\XD{D}{0}$, the weight
$(\kapD{D},\lastD{D}{0})$, the level $\colL{D}{0}$, which is free by Lemma~\ref{l:level}(d), the number
$N'=\NnD{D}\,d$, which lies in the different of $\ED{D}$ (unit~152), and
$M=\kD{D}\,\fD{D}+\fD{D}\,\aD{D}$, which is admissible by Lemma~\ref{l:branch}(b). The inequality
$(6N'(\fD{D}\kD{D})^2)^{6M}\le B'$ holds by size. So $\Col{D}$ is spanned by columns $\varphi$ with
$\short{B'}{\varphi}$. The operator $\ramop{D}{\piPl{D}}$ maps $\Col{D}$ onto $\Colram{D}$, and by
Proposition~\ref{p:ram} and Lemma~\ref{l:pi}(d) it maps such $\varphi$ to ramified columns that are
$4^{\fD{D}}B'$-short, that is, $\exp(\bud{D})$-short. The image of a spanning set spans the image.
\end{proof}
\leanline{\lean{spanshortramifiedColumns}\sep input of unit 152: \lnm{Parameters.cycloConductor\_mem\_different}}

\sect{The chain as a sum over classes}{s:chain}
\begin{setting}[a good datum at a $p$-adic place]\label{set:good}
Settings \ref{set:D} and~\ref{set:p} hold; $\vartheta$ is a twist of~$D$, and $\Psi$ is a finitely supported complex
function on the tuples of translations of~$D$ such that $\good{\vartheta}{\Psi}$; and $p$ is prime to
$\nrm{\vartheta}$.
\end{setting}

\statement{inksoft}{\textbf{Input I8 (standard data, Hecke operators, sums over classes; units 150, 141 and the
background library).}\enspace \textup{(i)} In Settings \ref{set:D} and~\ref{set:p}, for an integer $\delta$ of
$\KP{P}$ and a set $S$ of finite places of $\ED{D}$, $\intafter{\delta}{S}$ is the group of the matrices $y$ over
$\Af{\ED{D}}$ (of any fixed shape) such that every entry of $\delta y$ is integral at every place of~$S$; and
$\Std{\delta}{D,i}$ is the $\Oj$-submodule, of the model that contains $\Wfin{D}{i}$, generated by
the vectors $\mathrm{U}(y)\intvac{D,i}$ with $y\in\intafter{\delta}{\{w:\ p\in w\}}$. If $\delta\ne0$ and
$\Phi\in\Std{\delta}{D,i}$, then for every complex number~$c$, with the kernel taken at the scale
$c\cdot\per{P}{\Dmu{D}}$, $\Phi$ multiplies gauges at the point~$1$ by at most
$\bigl(|j(c)|_p/|j(\cj{\delta}\,d)|_p\bigr)^{L}$, $L=\boxes{\kapD{D}}{\nD{D}{i+1}}$; for the lift kernel
($c=\cj{\pig{D}}$) this reads
$\Phi\in\raises{\lk{D}{\vartheta}{i}}{1}{1}{(|j(\cj{\pig{D}})|_p/|j(\cj{\delta}d)|_p)^{L}}$.

\textup{(ii)} Let $w$ be a finite place of $E$ with $\cj{w}\ne w$ and $I$ a finite set. The group $\GL_I(E_w)$
embeds in $\Gf{I}$, an element $g$ going to the point with component $g$ at~$w$, the component at $\cj{w}$ that
unitarity forces, and $1$ elsewhere; for a family $(g_w)_{w\in S}$ over a set $S$ of such places, no two
conjugate, $\mathrm{at}(g)$ is the product of these points. For a natural number~$s$, the \emph{basic double
coset of index $s$} at $w$ is the set of the $g\in\GL_I(E_w)$ such that $g$ and $\varpi g^{-1}$ are integral
and $\det g$ has valuation~$s$ ($\varpi$ a uniformiser); for $s\le|I|$ it is the double coset of
$\diag(\varpi1_s,1)$ under the integral invertible matrices. On the vectors of a representation of $\Gf{I}$ that
are fixed by $\Kmax{I}{w}$, the Hecke operator $\Tl{w}{s}$ is the sum of the translates by representatives of
the left cosets of the integral invertible matrices in this double coset; it preserves the vectors fixed by a
group that contains $\Kmax{I}{w}$ and commutes with every operator that commutes with the group at~$w$.

\textup{(iii)} In Setting~\ref{set:D}, let $\vartheta$ be a twist, $\Psi$ a finitely supported function on the
tuples of translations, and $\Phi=(\Phi_{\mathsf{x},i})$ a family with $\Phi_{\mathsf{x},i}\in\Wfin{D}{i}$. The
combination $\sum_{\mathsf{x}}\Psi(\mathsf{x})\bigotimes_{i<m}\Phi_{\mathsf{x},i}$ is \emph{jointly invariant}
if it satisfies condition \textup{(G2)} of Definition~\ref{d:good} with $\Phi_{\mathsf{x},i}$ in place of
$\td{D}{\mathsf{x}}{i}$. Suppose that it is, that $\varphi$ is a column, and that
$\mathcal{C}_0,\dots,\mathcal{C}_m$ are properties of pairs (point, vector) such that $\mathcal{C}_0(g,\varphi(g))$
holds for all~$g$ and $\mathcal{C}_i(g,\xi)$ implies
$\mathcal{C}_{i+1}\bigl(h,\lkjm{D}{\vartheta}{i}(\lkom{D}{\vartheta}{i}(g,h)\Phi_{\mathsf{x},i})\xi\bigr)$ for
$\mathsf{x}$ in the support, $i<m$ and all $g,h$. Then for every point~$h$, the vector
$\sum_{\mathsf{x}}\Psi(\mathsf{x})\bigl(\chn{D}{\vartheta}{\Phi_{\mathsf{x}}}{m}\varphi\bigr)(h)$ is a sum of
$\prod_{i<m}\#\Cl{\colL{D}{i}}$ vectors of the form
$\sum_{\mathsf{x}}\Psi(\mathsf{x})\beta(\bigotimes_{i<m}\Phi_{\mathsf{x},i})$, with $\beta$ linear and
$\mathcal{C}_m\bigl(h,\beta(\bigotimes_{i<m}\Phi_{\mathsf{x},i})\bigr)$ for every $\mathsf{x}$ in the support.}
\leanline{unit 150: \lnm{Parameters.integralAfter}, \lnm{Parameters.standardData},
\lnm{Parameters.mem\_raises\_of\_standardData}\sep background library: \lnm{atSplit}, \lnm{atSplits},
\lnm{basicDoubleCoset}, \lnm{heckeSplit}, \lnm{map\_heckeSplit}\sep unit 141: \lnm{IsJointlyInvariant},
\lnm{IsJointlyInvariant.sum\_chain\_mem\_nsmul}}

The reason for (iii) is that a lift is a finite sum over the classes of its level, each term being linear in the
datum; $m$ lifts give a sum over tuples of classes, and joint invariance makes the sum of the terms independent
of the representatives, which is what allows one linear map $\beta$ for all~$\mathsf{x}$. The levels are free
because they lie in the chosen level (Input~I3).

\begin{lemma}[integrality of translations]\label{l:intafter}
In Setting~\ref{set:D}:
\begin{enumerate}[label=(\alph*)]
\item Let $q$ and $i$ be natural numbers, $\mathsf{x}\in\admT{q}$, $\delta$ an integer of $\KP{P}$, and $S$ a
set of finite places of $\ED{D}$ such that every $w\in S$ satisfies $q\notin w$ and
($\delta=\pig{D}$ or $\cj{w}$ is not chosen). Then $\mathsf{x}_i\in\intafter{\delta}{S}$.
\item Let $\mathsf{x}$ be a tuple of translations, $i$ a natural number and $S$ a set of finite places of
$\ED{D}$ such that $\mathsf{x}_i\in\intafter{1}{S\cup\cj{S}}$, where $\cj{S}$ is the set of the conjugates of the
places of~$S$. Then $\mathrm{P}^{S}\td{D}{\mathsf{x}}{i}=\td{D}{\mathsf{x}}{i}$.
\end{enumerate}
\end{lemma}
\begin{proof}
(a) At $w\in S$ the entries of $\pig{D}\mathsf{x}_i$ are integral, by the definition of $\admT{q}$. If
$\delta=\pig{D}$ this is the claim. Otherwise $\cj{\pig{D}}\notin\cj{w}$, that is, $\pig{D}\notin w$: so
$\pig{D}$ is a unit at~$w$, the entries of $\mathsf{x}_i$ are integral at~$w$, and so are those of
$\delta\mathsf{x}_i$. (b) The datum is the translate of a multiple of the vacuum by a matrix that is integral
at $S\cup\cj{S}$, and such a translate is again a vacuum at these places (a property of the projection, from
the library).
\end{proof}
\leanline{\lean{memintegralAfter}\sep\lean{projOntranslateDatum}}

\begin{lemma}[the data of a good datum are standard]\label{l:tdstd}
In Setting~\ref{set:good}, assume $p\ne\ND{D}$. Let $\mathsf{x}$ be in the support of~$\Psi$, $i$ a natural
number and $\delta$ an integer of $\KP{P}$ such that $\delta=\pig{D}$ or $p$ does not divide $\gD{D}$. Then
$\td{D}{\mathsf{x}}{i}\in\Std{\delta}{D,i}$.
\end{lemma}
\begin{proof}
By (G1), $\mathsf{x}\in\admT{\nrm{\vartheta}\ND{D}}$. Let $w$ be a place above~$p$. Then
$\nrm{\vartheta}\ND{D}\notin w$, because $p$ is prime to $\nrm{\vartheta}$ and $p\ne\ND{D}$. If $p$ does not
divide $\gD{D}$, then $\cj{w}$, which also lies above~$p$, is not chosen. So Lemma~\ref{l:intafter}(a) gives
$\mathsf{x}_i\in\intafter{\delta}{\{w:\ p\in w\}}$, and the datum is one of the generators of
$\Std{\delta}{D,i}$.
\end{proof}
\leanline{\lean{translateDatummemstandardData}}

\begin{lemma}[the chain is a plain sum]\label{l:chainsum}
In Setting~\ref{set:D}, let $\vartheta$ be a twist of~$D$, $\Psi$ a finitely supported complex function on the
tuples of translations of~$D$, and $\Phi=(\Phi_{\mathsf{x},i})$ a family of data
$\Phi_{\mathsf{x},i}\in\Wfin{D}{i}$ such that $\sum_{\mathsf{x}}\Psi(\mathsf{x})\bigotimes_{i<m}\Phi_{\mathsf{x},i}$
is jointly invariant. For every natural number $i$ let $\mathrm{N}_i$ be a real function of a point of\/
$\Gf{\XD{D}{i}}$ and a vector, and $r_i\ge0$ a real number, such that for every $\mathsf{x}$ in the support
of~$\Psi$, every $i<m$ and all $g$, $h$, $\xi$
\[\mathrm{N}_{i+1}\bigl(h,\lkjm{D}{\vartheta}{i}(\lkom{D}{\vartheta}{i}(g,h)\Phi_{\mathsf{x},i})\xi\bigr)\ \le\ r_i\cdot\mathrm{N}_i(g,\xi).\]
Let $\varphi\in\Col{D}$ and let $B$ be a real number with $\mathrm{N}_0(g,\varphi(g))\le B$ for all~$g$, and let
$g$ be a point of\/ $\Gf{\XD{D}{m}}$. Then there are linear maps $\beta_c$, for
$1\le c\le\prod_{i<m}\#\Cl{\colL{D}{i}}$, from $\bigotimes_{i<m}\Wfin{D}{i}$ to
$\CC^{\Wy{\XD{D}{m}}{\kapD{D}}}$ such that
$\mathrm{N}_m\bigl(g,\beta_c(\bigotimes_{i<m}\Phi_{\mathsf{x},i})\bigr)\le\prod_{i<m}r_i\cdot B$ for every $c$ and
every $\mathsf{x}$ in the support, and
\[\sum_{c}\sum_{\mathsf{x}}\Psi(\mathsf{x})\,\beta_c\Bigl(\bigotimes_{i<m}\Phi_{\mathsf{x},i}\Bigr)=\Bigl(\sum_{\mathsf{x}}\Psi(\mathsf{x})\,\chn{D}{\vartheta}{\Phi_{\mathsf{x}}}{m}\varphi\Bigr)(g).\]
\end{lemma}
\begin{proof}
Input~I8(iii) with the properties $\mathcal{C}_i(g,\xi)$: $\mathrm{N}_i(g,\xi)\le\prod_{i'<i}r_{i'}\cdot B$. They
propagate by the hypothesis on the kernels and $r_i\ge0$.
\end{proof}
\leanline{\lean{existschaineqsum}}

\begin{proposition}[gains of the lifts add along the chain]\label{p:gchain}
In Setting~\ref{set:good}, for every natural number $i$ let $t_i$ be a $\CC$-linear endomorphism of\/
$\Wfin{D}{i}$ such that
$\lkom{D}{\vartheta}{i}(a,a')\,t_i\,\td{D}{\mathsf{x}}{i}=t_i\,\lkom{D}{\vartheta}{i}(a,a')\,\td{D}{\mathsf{x}}{i}$
for every $\mathsf{x}$ in the support of~$\Psi$, every $a\in\colL{D}{i}$ and every $a'\in\colL{D}{i+1}$; let
$d_i\in\Gf{\XD{D}{i}}$ be points and $r_i$ real numbers such that
\[t_i\,\td{D}{\mathsf{x}}{i}\in\raises{\lk{D}{\vartheta}{i}}{d_i}{d_{i+1}}{p^{\,r_i}}\qquad\text{for every }\mathsf{x}\text{ in the support of }\Psi\text{ and every }i<m.\]
Let $\varphi\in\Col{D}$ have $j$-gauge at most~$1$ at~$d_0$. Then the form
$\sum_{\mathsf{x}}\Psi(\mathsf{x})\cdot\chn{D}{\vartheta}{(t_i\td{D}{\mathsf{x}}{i})_i}{m}\varphi$ has $j$-gauge
at most $p^{\sum_{i<m}r_i}$ at~$d_m$.
\end{proposition}
\begin{proof}
The combination of the data $t_i\td{D}{\mathsf{x}}{i}$ is jointly invariant: it is the image of the tensor of
(G2) under $\bigotimes_it_i$, which commutes with the action of the levels on the data in the support. Apply
Lemma~\ref{l:chainsum} with $\mathrm{N}_i(g,\xi)=\gauge{\kapD{D}}{\lastD{D}{i}}{gd_i}{\xi}$, the numbers
$p^{\,r_i}$ and $B=1$. The value of the form at a point is a sum of terms $\Psi(\mathsf{x})\beta_c(\dots)$; the
gauge is ultrametric, the gauge of each $\beta_c(\dots)$ is at most $p^{\sum r_i}$, and
$|j(\Psi(\mathsf{x}))|_p\le1$ by (G1), since $p$ is prime to $\nrm{\vartheta}$.
\end{proof}
\leanline{\lean{gaugeLechain}}

\sect{Served primes, induced primes and marked sets}{s:marks}
\begin{lemma}[orders and losses]\label{l:served}
In Setting~\ref{set:D}:
\begin{enumerate}[label=(\alph*)]
\item For every prime number $p$, $\rr{D}{p}\le\log x_1$.
\item If $p$ is a prime number that is served for $D$, then $m$ divides the order of the place $\pin{D}{-p}$ in
the group $\PlD{D}$.
\item For every natural number $p$, $\ \fD{D}\,(\log b)^2\cdot\aD{D}/p\ \le\ \loss{D}{p}$.
\item For natural numbers $p$ and $e$ with $e$ odd and every $\tau\in\tops{D}$: the place
$\pin{D}{-p^{e}}\cdot\tau$ lies in $\tops{D}$ if and only if $\oz{D}{p}$ divides~$e$.
\end{enumerate}
\end{lemma}
\begin{proof}
(a) $m\le2^{\Dmu{D}}\le p^{\Dmu{D}}$, so $\rr{D}{p}\le\Dmu{D}$, and $\Dmu{D}\le\log x_1$ by size. (c) The last
factor of $\loss{D}{p}$ is $m\ge1$ or~$1$. (d) Since $e$ is odd, $\pin{D}{-p^{e}}=\pin{D}{-p}^{\,e}$; the tops
are the places of class $1$ modulo~$N_0$, so the product is a top exactly when the class of $\pin{D}{-p}$ has
order dividing~$e$. (b) The data have $m=2^{c+1}$ and $\Dmu{D}=22(c+1)$ for a natural number~$c$, and
$2^{\Dmu{D}+1}$ divides $\NnD{D}$. Let $n\ge1$ be the order of $\pin{D}{-p}$. The class of this place modulo
$2^{\Dmu{D}+1}$ is that of $-p$ (unit~152; $p$ is prime to $N_0$ because $\chiK{P}(p)\ne1$), so
$(-p)^{n}\equiv1$ modulo $2^{\Dmu{D}+1}$. By lifting the exponent at~$2$ (background library),
$2^{\Dmu{D}}$ divides $(-p-1)n$ or $(-p+1)n$. Since $p$ is not short, $2^{c+1}$ divides neither $p+1$ nor $p-1$;
so $2^{\Dmu{D}-c}$ divides~$n$, and $\Dmu{D}-c\ge c+1$.
\end{proof}
\leanline{\lean{instUbCastRealRrLogLx}\sep\lean{mdvdorderOf}\sep\lean{leloss}\sep\lean{placeOfIntmulmemtopsiff}\sep
background library: \lnm{two\_pow\_dvd\_or\_of\_pow\_eq\_one}\sep unit 152: \lnm{Parameters.placeClass\_placeOfInt}}

\begin{lemma}[induced primes]\label{l:ind}
In Settings \ref{set:D} and~\ref{set:p}, let $\tau,\tau'\in\PlD{D}$.
\begin{enumerate}[label=(\alph*)]
\item $p\in\ind{\tau}$.
\item $\ind{\tau\cdot\tau'}=\tau'\cdot\ind{\tau}$, where a place acts on the primes through the automorphism of
$\ED{D}$ over $\KP{P}$ that it names.
\item If $\chiK{P}(p)=-1$ and $e$ is an odd natural number, then
$\ind{\pin{D}{-p^{e}}\cdot\tau}=\cj{\ind{\tau}}$.
\item If $\chiK{P}(p)=-1$, then $\ind{\tau}$ is equal to $\ind{\tau'}$ or to $\cj{\ind{\tau'}}$ if and only if
$\tau=\pin{D}{-p}^{\,n}\cdot\tau'$ for some natural number~$n$.
\end{enumerate}
\end{lemma}
\begin{proof}
(a) $|j(p)|_p<1$. (b) The embedding at $\tau\tau'$ is the embedding at $\tau$ composed with the automorphism
named by~$\tau'$ (unit~152). (c), (d) For $p$ inert in $\KP{P}$ and prime to $\NnD{D}\,d$, the decomposition of
$p$ in $\ED{D}$ is described in the note on the CM field: a place $v$ maps a prime above $p$ to itself or to its
conjugate exactly when $v$ is a power of $\pin{D}{-p}$, to the conjugate for odd powers. With (b) this gives
both claims.
\end{proof}
\leanline{\lean{pmeminducedPlace}\sep\lean{inducedPlacemul}\sep\lean{inducedPlacenegpowmulofodd}\sep\lean{inducedPlaceeqorstariff}\sep
unit 152: \lnm{Parameters.phi\_mul}, \lnm{Parameters.smul\_eq\_star\_iff\_inert}, \lnm{Parameters.smul\_eq\_iff\_inert}}

\begin{lemma}[a full mark exists at a served prime]\label{l:exmark}
In Settings \ref{set:D} and~\ref{set:p}, assume that $p<x_1$ and that $p$ is served for~$D$. Then there is a
finite set $\mathcal{L}$ of finite places of $\ED{D}$ that is a full mark.
\end{lemma}
\begin{proof}
Fix a well-ordering of the finite places and let $\mathcal{L}$ be the set of the primes $\mathfrak{l}=\ind{\tau}$,
$\tau\in\tops{D}$, that are smaller than their conjugate. Then (M1) and (M3) hold by construction, and (M2),
(M4) by hypothesis, a served prime being inert. For the covering, let $\tau$ be a top. By Lemma~\ref{l:served}(d)
the place $\tau'=\pin{D}{-p}^{\,\oz{D}{p}}\cdot\tau$ is a top, and it induces $\cj{\ind{\tau}}$ by
Lemma~\ref{l:ind}(c), $\oz{D}{p}$ being odd. Moreover $\ind{\tau}\ne\cj{\ind{\tau}}$: by
Lemma~\ref{l:served}(b) the order of $\pin{D}{-p}$ is even, so its odd powers do not fix a prime above~$p$. Hence
one of the two primes is the smaller and lies in~$\mathcal{L}$.
\end{proof}
\leanline{\lean{existsisMarked}}

\begin{setting}[a marked set]\label{set:mark}
Settings \ref{set:D} and~\ref{set:p} hold, and $\mathcal{L}$ is a finite set of finite places of $\ED{D}$ that is
marked for $D$, $p$ and~$j$ (Input~I5).
\end{setting}

\begin{lemma}[marked sets]\label{l:markprop}
In Setting~\ref{set:mark}:
\begin{enumerate}[label=(\alph*)]
\item $\#\mathcal{L}\le\fD{D}$.
\item $p$ does not divide $\gD{D}$, $p\ne\ND{D}$, and $p$ does not divide $\gD{D}\,\ND{D}$.
\item For every $\mathfrak{l}\in\mathcal{L}$: $p\in\mathfrak{l}$; $p$ is a uniformiser at~$\mathfrak{l}$;
$\cj{\mathfrak{l}}\ne\mathfrak{l}$; and $\Kmax{\XD{D}{0}}{\mathfrak{l}}\subseteq\colL{D}{0}$.
\end{enumerate}
\end{lemma}
\begin{proof}
(a) $\mathcal{L}$ lies in the image of $\tau\mapsto\ind{\tau}$. (b) $p<x_1$, and the primes $g_0$, $g_1$, $\ND{D}$
are larger than $x_1$ by size. (c) The first claim is Lemma~\ref{l:ind}(a). $p$ is unramified in $\ED{D}$, being
prime to $\NnD{D}\,d$ (an inert or ramified $p$ does not divide $\NnD{D}$, and if $\mathcal{L}\ne\emptyset$ then
$p$ is served, so $p\ne d$). The third claim is (M1). For the last, by (b) the element $\cj{\pig{D}}\piN{D}$ is a unit at $\mathfrak{l}$ and
$\cj{\mathfrak{l}}$, so $\Kmax{\XD{D}{0}}{\mathfrak{l}}$ lies in the principal level of this element, and
Lemma~\ref{l:level}(a) applies.
\end{proof}
\leanline{\lean{cardle}\sep\lean{notdvdgProdmulN}\sep\lean{pmem}\sep\lean{valuationp}\sep\lean{starne}\sep\lean{maxCompactAtlecolumnLevel}}

\begin{lemma}[marked primes avoid the thick places; a balance]\label{l:thick}
In Setting~\ref{set:mark}:
\begin{enumerate}[label=(\alph*)]
\item If $\tau\in\PlD{D}$ is such that $\ind{\tau}\in\mathcal{L}$ or $\cj{\ind{\tau}}\in\mathcal{L}$, then
$\tau\notin\thick{D}$.
\item For every natural number $i$ and every real function $G$ of a natural number and a finite place,
\[\sum_{\tau\in\PlD{D}}\Bigl([\ind{\tau}\in\mathcal{L}]\,G\bigl(\wtD{D}{i}{\tau},\ind{\tau}\bigr)-[\cj{\ind{\tau}}\in\mathcal{L}]\,G\bigl(\wtD{D}{i}{\tau},\cj{\ind{\tau}}\bigr)\Bigr)=0 .\]
\end{enumerate}
\end{lemma}
\begin{proof}
(a) $\mathcal{L}$ is not empty, so $p$ is served: $\chiK{P}(p)=-1$ and $\oz{D}{p}$ is odd. By (M3) and
Lemma~\ref{l:ind}(d), $\tau=\pin{D}{-p}^{\,n}\cdot\tau_0$ for a top $\tau_0$ and some~$n$, so the class of $\tau$
modulo $N_0$ is that of $(-p)^{n}$. A thick place has class~$-1$ modulo~$N_0$ (a condition of the data). If
$\tau$ were thick, raising $(-p)^{n}\equiv-1$ to the odd power $\oz{D}{p}$ would give $1\equiv-1$ modulo~$N_0$.
(b) If $\mathcal{L}$ is empty both sums vanish. Otherwise $p$ is inert; substitute
$\tau\mapsto\pin{D}{-p}\cdot\tau$ in the second sum. By Lemma~\ref{l:ind}(c) the prime induced by the new place
is $\cj{\ind{\tau}}$, so the condition becomes $\ind{\tau}\in\mathcal{L}$ and the place argument of $G$ becomes
$\ind{\tau}$; and the weight is unchanged, because by (a) neither $\tau$ nor $\pin{D}{-p}\cdot\tau$ is thick,
and $\wtD{D}{i}{\,\cdot\,}$ takes one value on the places that are not thick.
\end{proof}
\leanline{\lean{notMemthick}\sep\lean{suminducedPlacesubstareqzero}}

\sect{The jet exponent}{s:jet}
\begin{lemma}[collisions]\label{l:coll}
In Setting~\ref{set:D}, let $\Xi$ be a finite subset of\/ $\PlD{D}$.
\begin{enumerate}[label=(\alph*)]
\item Let $p$ and $e$ be natural numbers and $\tau\in\PlD{D}$. If $(\tau,e)$ is a collision of $p$ outside $\Xi$
for~$D$, then $p\ne d$, $e$ is odd, and $\oz{D}{p}$ divides~$e$.
\item Let $p$ be a prime number and $o$ a natural number such that $o\le e$ for every collision $(\tau,e)$ of
$p$ outside $\Xi$ for~$D$. Then $\coll{D}{p}{\Xi}\le2\,\fD{D}\,\aD{D}/p^{\,o}$.
\item Let $p$ be a prime number that is small and not split for~$D$ and is not served for~$D$. Then
$m\cdot\coll{D}{p}{\emptyset}\le\loss{D}{p}$.
\item In Setting~\ref{set:p}, if $\mathcal{L}$ is a full mark, then $\coll{D}{p}{\even{\mathcal{L}}}=0$.
\end{enumerate}
\end{lemma}
\begin{proof}
(a) Both $\tau$ and $\pin{D}{-p^{e}}\cdot\tau$ are tops, so they have the same class modulo~$N_0$; by a lemma
of unit~150 this excludes $p=d$ (no power of $d$ is $\equiv-1$ modulo~$N_0$). Then $e$ is odd because $(e,p)$ is
live, and Lemma~\ref{l:served}(d) applies. (b) For each $\tau$ the terms are $\aD{D}/p^{e}$ with $e\ge o$; they
are bounded by the geometric series $\sum_{n\ge0}(\aD{D}/p^{\,o})2^{-n}=2\aD{D}/p^{\,o}$, since $p\ge2$; sum over
the $\fD{D}$ places. (c) By $(\log b)^2\ge2$, $\loss{D}{p}\ge2\fD{D}(\aD{D}/p)\cdot(m\text{ or }1)$. If $p$ is short,
use (b) with $o=1$. If not, every collision has $e\ge\rr{D}{p}+1$: by (a), $e$ is odd, $\oz{D}{p}$ divides~$e$,
hence $\oz{D}{p}$ is odd, and $\chiK{P}(p)=-1$; if $\oz{D}{p}\le e\le\rr{D}{p}$, then $p$ would be served. So
$\coll{D}{p}{\emptyset}\le2\fD{D}\aD{D}/p^{\rr{D}{p}+1}$ by (b), and $m\le p^{\rr{D}{p}}$. (d) Suppose
$(\tau,e)$ is a collision outside $\even{\mathcal{L}}$. Then $\mathcal{L}$ is not empty, since it covers the
top~$\tau$; so $p$ is inert, $e$ is odd by (a), and $\pin{D}{-p^{e}}\cdot\tau$ induces $\cj{\ind{\tau}}$
(Lemma~\ref{l:ind}(c)). As $\mathcal{L}$ covers~$\tau$, one of $\tau$, $\pin{D}{-p^{e}}\cdot\tau$ lies in
$\even{\mathcal{L}}$, a contradiction.
\end{proof}
\leanline{\lean{neandoddanddvd}\sep\lean{collisionSumle}\sep\lean{mulcollisionSumleloss}\sep\lean{collisionSumevenPlaceseqzero}\sep
unit 150: \lnm{Data.ne\_d\_of\_placeClass\_eq}}

\begin{lemma}[low columns and low boxes]\label{l:low}
In Setting~\ref{set:D}, let $i<m$ be a natural number and $\tau\in\PlD{D}$. Then
\[\kD{D}\cdot\nomW{D}{\tau}\le\lb{D}{i}{\tau},\qquad\lc{D}{i}{\tau}\le\nomW{D}{\tau}+x^{60},\qquad\lb{D}{i}{\tau}\le\kD{D}\cdot\nomW{D}{\tau}+x^{60}.\]
\end{lemma}
\begin{proof}
The term $h=\kD{D}$ of $\lb{D}{i}{\tau}$ is
$\kD{D}\bigl(\wtD{D}{1}{\tau}-\lfloor(\kD{D}+\tD{D}+1)/2\rfloor\bigr)\ge\kD{D}\,\nomW{D}{\tau}$ (unit~152). For the
upper bounds, each $\kap{D}{\tau}{h}$ with $h\ne1$ is at most $[h=\kD{D}]\wtD{D}{1}{\tau}+\gap{D}$ (unit~150),
$\wtD{D}{1}{\tau}$ exceeds $\nomW{D}{\tau}$ by a small amount (unit~152), and the sums of the remaining terms
over $h<\nD{D}{i+1}\le\nD{D}{m}$ are at most $x^{60}$ by size.
\end{proof}
\leanline{\lean{lowColumnslowBoxesbounds}\sep unit 152: \lnm{Data.nominalWeight\_le\_w\_sub},
\lnm{Data.w\_le\_nominalWeight\_add}\sep unit 150: \lnm{Data.κ\_le}}

\begin{proposition}[the jet exponent is at least the gain of one lift, up to collisions]\label{p:jet}
In Setting~\ref{set:D}, let $p$ be a prime number that is small and not split for~$D$, $i<m$ a natural number
and $\Xi$ a finite subset of\/ $\PlD{D}$. Then
\begin{multline*}
\frac{\gain{D}{p}}{m}-\coll{D}{p}{\Xi}+\frac{1-\chiK{P}(p)}{2}\cdot\#\bigl(\tops{D}\cap\Xi\bigr)\cdot\frac{\aD{D}}{x_1}\\
-\fD{D}\,(\log b)^2\,\frac{\aD{D}}{p}\cdot\frac1m\ \le\ \jetexp{p}{i}{\Xi}.
\end{multline*}
\end{proposition}
\begin{proof}
Write $c=(1-\chiK{P}(p))/2\in\{\tfrac12,1\}$, $n'=\nD{D}{i+1}$, $L=\boxes{\kapD{D}}{n'}$,
$l=\lfloor\log_pL\rfloor$, $\eta=\aD{D}/(pm)$, and $\tau_e=\pin{D}{-p^{e}}\cdot\tau$. Since $n'>1$, the columns
of height one contribute $\aD{D}$ at each top, so $L=\sum_\tau\bigl(\kap{D}{\tau}{1}+\lb{D}{i}{\tau}\bigr)$, and
by \eqref{e:live}
\[\Bigl(\frac1{p-1}-\nsw{d}{p}\Bigr)L=c\sum_{\tau}\ \sum_{e:\ (e,p)\text{ live}}\frac{\kap{D}{\tau}{1}+\lb{D}{i}{\tau}}{p^{e}}.\]
By the definitions of Input~I6,
$\gain{D}{p}/m=c\sum_\tau\sum_{e\text{ live}}\bigl([\tau\in\tops{D}](\aD{D}/p^{e}-\nomW{D}{\tau_e})_++\kD{D}(\nomW{D}{\tau}/p^{e}-\nomW{D}{\tau_e})_+\bigr)$.
For a place $\tau$ and a natural number $e$ let $\mathrm{lv}_e$ be $1$ if $(e,p)$ is live and $0$ otherwise,
and $\mathrm{cl}_{\tau,e}$ be $1$ if $(\tau,e)$ is a collision of $p$ outside $\Xi$ for $D$ and $0$ otherwise.
The claim follows by summation over $\tau$ and~$e$ from the following inequality, for each $\tau$ and~$e$:
\begin{multline}\label{e:term}
[e\le l]\Bigl([e\text{ odd or }p=d]\;\cc{p}{\kapD{D}}{n'}{\Xi}{\tau}{\tau_e}{e}-4\eta\Bigr)+[e=1]\,[\tau\in\tops{D}\cap\Xi]\,\frac{\aD{D}}{x_1}\\
+\mathrm{lv}_e\Bigl([\tau\in\tops{D}]\bigl(\tfrac{\aD{D}}{p^{e}}-\nomW{D}{\tau_e}\bigr)_++\kD{D}\bigl(\tfrac{\nomW{D}{\tau}}{p^{e}}-\nomW{D}{\tau_e}\bigr)_+\Bigr)\\
\le\ \mathrm{lv}_e\,\frac{\kap{D}{\tau}{1}+\lb{D}{i}{\tau}}{p^{e}}+\mathrm{cl}_{\tau,e}\,\frac{\aD{D}}{p^{e}} .
\end{multline}
Indeed, multiply \eqref{e:term} by $c$ and sum: the right side gives
$(1/(p-1)-\nsw{d}{p})L+c\,\coll{D}{p}{\Xi}$, the left side gives
$c\sum\cc{p}{\kapD{D}}{n'}{\Xi}{\tau}{\tau_e}{e}-4c\,\eta\,\fD{D}(l+1)+\gain{D}{p}/m+c\#(\tops{D}\cap\Xi)\aD{D}/x_1$,
and therefore
\[\jetexp{p}{i}{\Xi}\ \ge\ \frac{\gain{D}{p}}{m}+c\,\#(\tops{D}\cap\Xi)\frac{\aD{D}}{x_1}-c\,\coll{D}{p}{\Xi}-\bigl(4c\,\eta+n'^{\,3}\bigr)\fD{D}(l+1).\]
Now $c\le1$, the collision sum is non-negative, $n'^{\,3}\le x^{60}\le\eta$ by size (here $p<x_1$ is used), and
$5(l+1)\le(\log b)^2$, because $l\le\log L/\log2$ and $L\le\gD{D}\le b^{22}$ by Lemma~\ref{l:low}
and by size. This gives the proposition.

\emph{Proof of \eqref{e:term}.} The facts used are: $\aD{D}/x_1\le\nomW{D}{\tau'}\le\aD{D}/x$ for every
place~$\tau'$, with $\nomW{D}{\tau'}=\aD{D}/x_1$ when $\tau'$ is not thick; Lemma~\ref{l:low}; and the sizes
$x^{60}\le\eta$, $\aD{D}/(xp)+x^{60}/p\le\eta$. Put $\tau'=\tau_e$ and recall that in the coupling count
$\mathrm{one}_\tau=[\tau\in\tops{D}\setminus\Xi]\,\aD{D}$, the count of the other columns is $\lc{D}{i}{\tau}$ and their boxes
are $\lb{D}{i}{\tau}$.

\emph{The pair $(e,p)$ is live and $e\le l$.} Three bounds add up to \eqref{e:term}.
(1) $\lc{D}{i}{\tau}/p^{e}\le\eta$.
(2) $\min\bigl(\lb{D}{i}{\tau}/p^{e},\lb{D}{i}{\tau'}\bigr)+\kD{D}\bigl(\nomW{D}{\tau}/p^{e}-\nomW{D}{\tau'}\bigr)_+\le\lb{D}{i}{\tau}/p^{e}+x^{60}$:
if the positive part vanishes this is clear; otherwise bound the minimum by
$\lb{D}{i}{\tau'}\le\kD{D}\nomW{D}{\tau'}+x^{60}$ and use $\kD{D}\nomW{D}{\tau}\le\lb{D}{i}{\tau}$.
(3) $\min\bigl(\mathrm{one}_\tau/p^{e},\lc{D}{i}{\tau'}+\mathrm{one}_{\tau'}\bigr)+[\tau\in\tops{D}]\bigl(\aD{D}/p^{e}-\nomW{D}{\tau'}\bigr)_++[e=1][\tau\in\tops{D}\cap\Xi]\aD{D}/x_1$
is at most $\kap{D}{\tau}{1}/p^{e}+x^{60}$ plus $\aD{D}/p^{e}$ if $(\tau,e)$ is a collision. If $\tau$ is not a
top, the left side is~$0$. If $\tau\in\tops{D}\cap\Xi$, then $\mathrm{one}_\tau=0$ and
$(\aD{D}/p^{e}-\nomW{D}{\tau'})_+\le\aD{D}/p^{e}-\aD{D}/x_1$ whenever it is positive; for $e=1$ it is, since
$p<x_1$. If $\tau\in\tops{D}\setminus\Xi$ and $(\tau,e)$ is a collision, both terms are at most $\aD{D}/p^{e}$.
If $\tau\in\tops{D}\setminus\Xi$ and it is not a collision, then $\mathrm{one}_{\tau'}=0$, and either the positive part
vanishes, or the minimum is at most $\lc{D}{i}{\tau'}\le\nomW{D}{\tau'}+x^{60}$ and the sum is at most
$\aD{D}/p^{e}+x^{60}$. Altogether the loss is at most $\eta+2x^{60}\le4\eta$.

\emph{The pair $(e,p)$ is live and $e>l$.} Only the second and third groups of the left side occur, and
$\kD{D}(\nomW{D}{\tau}/p^{e}-\nomW{D}{\tau'})_+\le\lb{D}{i}{\tau}/p^{e}$,
$[\tau\in\tops{D}]\bigl((\aD{D}/p^{e}-\nomW{D}{\tau'})_++[e=1][\tau\in\Xi]\aD{D}/x_1\bigr)\le\kap{D}{\tau}{1}/p^{e}$,
as before.

\emph{The pair is not live.} Then $e=0$, or $e$ is even and $p\ne d$, and the right side is~$0$. For $e\ge2$
the left side is $-4\eta$ or~$0$. For $e=0$ and $p\ne d$ likewise. For $e=0$ and $p=d$ the place
$\tau'=\pin{D}{-1}\cdot\tau$ has the class of $\tau$ times~$-1$ modulo~$N_0$, so at most one of $\tau$, $\tau'$ is
a top and at most one is thick; hence the first minimum is at most $\lc{D}{i}{\tau'}$ or~$0$, the second is at
most $\lb{D}{i}{\,\cdot\,}$ at a place that is not thick, which is at most $\kD{D}\aD{D}/x_1+x^{60}$, and with
$p=d$ these three terms are at most $\eta$ each by size. So the coupling count is at most $3\eta$.
\end{proof}
\leanline{\lean{lejetExponent}\sep unit 150: \lnm{hasSum\_live\_inv\_pow}, \lnm{Data.summable\_localAsk}\sep unit 152:
\lnm{Data.nominalWeight\_mem\_Icc}, \lnm{Parameters.placeClass\_of\_thick}}

\begin{lemma}[the cost of the marks is within the loss]\label{l:lowsum}
In Setting~\ref{set:mark}, let $i<m$ be a natural number. Then
\[\#\mathcal{L}\cdot\DeltaP{D}{p}+\rr{D}{p}\sum_{\tau\in\even{\mathcal{L}}}\lc{D}{i}{\tau}\ \le\ \frac{\loss{D}{p}}{m}+\#\bigl(\tops{D}\cap\even{\mathcal{L}}\bigr)\cdot\frac{\aD{D}}{x_1}.\]
\end{lemma}
\begin{proof}
It suffices to show $\rr{D}{p}\,\lc{D}{i}{\tau}\le\aD{D}/(pm)+[\tau\in\tops{D}]\,\aD{D}/x_1$ for
$\tau\in\even{\mathcal{L}}$: then the left side is at most
$\fD{D}\rr{D}{p}\fD{D}(\nD{D}{m}+\rr{D}{p})+\fD{D}\aD{D}/(pm)+\#(\tops{D}\cap\even{\mathcal{L}})\aD{D}/x_1$ by
Lemma~\ref{l:markprop}(a), and the first two terms are at most $\fD{D}(\log b)^2\aD{D}/(pm)\le\loss{D}{p}/m$ by
size and Lemma~\ref{l:served}(c). Let $\tau\in\even{\mathcal{L}}$. By Lemma~\ref{l:thick}(a), $\tau$ is not
thick, so $\lc{D}{i}{\tau}\le\aD{D}/x_1+x^{60}$ by Lemma~\ref{l:low}. If $p\le m$, then
$\rr{D}{p}(\aD{D}/x_1+x^{60})\le\aD{D}/(pm)$ by size, using Lemma~\ref{l:served}(a). If $m\le p$, then
$\rr{D}{p}=1$; $p$ is served because $\mathcal{L}$ is not empty, so $\oz{D}{p}\le1$, that is, $-p\equiv1$
modulo~$N_0$; by (M3) and Lemma~\ref{l:ind}(d), $\tau$ is a power of $\pin{D}{-p}$ times a top, hence a top; and
$x^{60}\le\aD{D}/(pm)$ by size.
\end{proof}
\leanline{\lean{sumlowColumnsle}}

\sect{Hecke operators through the chain}{s:hecke}
\statement{inksoft}{\textbf{Input I9 (divisibility of jets; unit 150, and the Hecke transfer of the background
library).}\enspace \textup{(i)} (\emph{Gauges from a jet bound.}) In Settings \ref{set:D} and~\ref{set:p}, let
$\vartheta$ be a twist, $i$ a natural number, $\delta$ an integer of $\KP{P}$ and $\Phi'\in\Std{\delta}{D,i}$.
Let $y,y'\in\Gf{\XD{D}{i+1}}$ be integral and invertible at all primes above~$p$, and let $y''$ be a point
of\/ $\Gf{\XD{D}{i+1}}$ whose image under $x\mapsto j(\emb{\tau}(\cj{x}))$ is, for every $\tau$, a diagonal
matrix $\diag(u_\tau)$. Let $B$ be a real number such that for every $\Phi\in\Std{\delta}{D,i}$ and all Weyl
indices $c'$, $c$ the entry $(c',c)$ of the jet matrix of~$\Phi$ at the scale of the lift kernel, multiplied by
$\prod_\tau(\prod_{a}|u_\tau(a)|_p)^{\last{D}{i+1}{\tau}}$ and by the product of $|u_\tau(a)|_p$ over the
entries $a$ of the tableau $c'$ at~$\tau$, has $p$-adic absolute value at most~$B$. Then
$\lkom{D}{\vartheta}{i}(1,yy''y')\Phi'\in\raises{\lk{D}{\vartheta}{i}}{1}{1}{B}$.

\textup{(ii)} (\emph{Jets of standard data are governed by theta products.}) For $\Phi\in\Std{\delta}{D,i}$, a
complex number $c$ and a $\CC$-linear form $\mathrm{lf}$ on power series: if $|j(\mathrm{lf}(\Theta))|_p\le r$ for
every power series $\Theta$ whose image under $j$ lies in the set
$\ThProd{P}{\Dmu{D}}{\delta}{|j(c\,s_1)|_p}{\XD{D}{i+1}\times\XD{D}{i}}$ of theta products of the CM curve
($s_1$ the period factor), then $|j(\mathrm{lf}(T))|_p\le r$ for the Taylor series $T$ of the theta function of
$\Phi$ at the scale $c\cdot\per{P}{\Dmu{D}}$. If $p$ does not divide $\NnD{D}$, then $|j(s_1)|_p=1$.

\textup{(iii)} (\emph{Normal form and divisibility of Schur jets.}) Let $\chiK{P}(p)\ne1$. Every theta product
$\Theta$ in $\ThProd{P}{\Dmu{D}}{1}{r}{I}$ with $r\le1$ is \emph{integral on the null cone} in every degree,
for the maps $\tau\mapsto\pin{D}{-p^{e}}\cdot\tau$, the exponents $e\le\lfloor\log_pL\rfloor$ with $e$ odd or
$p=d$, and elements $\varpi$ of $\CC_p$ and $\varpi_0$ of $\op$ with $|\varpi|_p=p^{-\nsw{d}{p}}$ and
$|\varpi_0|_p=p^{(\chiK{P}(p)-1)/2}$. For a series $\Theta$ that is integral on the null cone up to the degree
$L=\boxes{\varkappa}{n'}$, a finite set $\Xi$ of places and a family of \emph{charges}, natural numbers
$\mathrm{ch}(\tau,h,a)$ attached to a place, a height and an index, such that for every exponent $e$ as above,
every place $\tau$, height $h$ and index~$a$
\begin{equation}\label{e:charge}
\tau\in\Xi\ \text{ or }\ (h=1\text{ and }\pin{D}{-p^{e}}\cdot\tau\in\Xi)\quad\Longrightarrow\quad\mathrm{ch}(\tau,1,a)\ne0\ \text{ or }\ \mathrm{ch}(\pin{D}{-p^{e}}\cdot\tau,h,a)\ne0,
\end{equation}
the entry $(c',c)$ of the matrix of jets of $\Theta$ along the Schur matrix of $\varkappa$ has absolute value
at most $p$ to the power: the total charge of the boxes of the tableau $c'$, minus
\[\bigl(\tfrac1{p-1}-\nsw{d}{p}\bigr)L-\tfrac{1-\chiK{P}(p)}{2}\sum_{e,\tau}\cc{p}{\varkappa}{n'}{\Xi}{\tau}{\pin{D}{-p^{e}}\cdot\tau}{e}-\#\PlD{D}\,n'^{\,3}\bigl(\lfloor\log_pL\rfloor+1\bigr).\]
For $\varkappa=\kapD{D}$ and $n'=\nD{D}{i+1}$ the subtracted number is $\jetexp{p}{i}{\Xi}$.

\textup{(iv)} (\emph{Hecke transfer.}) Let $w$ be a finite place with $\cj{w}\ne w$; consider a theta lift from
rank $n$ to rank $n+k'$ whose two characters are unramified at~$w$ and whose datum is a vacuum at $w$
and~$\cj{w}$. Then the lift of a form $\varphi$ fixed by the group at~$w$ is fixed by the group at~$w$, and the
Hecke polynomials $\mathrm{H}_w(X)=\sum_{n''}(-X)^{n''}\mathrm{q}_w^{\,n''(n''-1)/2}\,\Tl{w}{n''}$ satisfy
\[\mathrm{H}_w(X)\,\mathrm{lift}(\varphi)=\prod_{l<k'}\bigl(1-\mathrm{q}_w^{\,l}\,\alpha X\bigr)\cdot\mathrm{lift}\bigl(\mathrm{H}_w(\mathrm{q}_w^{\,k'}\alpha\alpha'^{-1}X)\varphi\bigr),\]
where $\alpha$ and $\alpha'$ are the values at a uniformiser of~$w$ of the Hecke characters of the two
characters, normalised by $\mathrm{q}_w^{\,n/2}$ and $\mathrm{q}_w^{\,(n+k')/2}$.}
\leanline{unit 150: \lnm{Parameters.mem\_raises\_of\_jet\_bound}, \lnm{Parameters.norm\_jetSeries\_le},
\lnm{Parameters.norm\_iota\_periodFactor\_eq\_one}, \lnm{Parameters.thetaProducts\_normal\_form},
\lnm{norm\_schurJet\_le}\sep background library: \lnm{hecke\_transfer\_of\_pairing}, a form of
\lnm{hecke\_transfer\_of\_equivariant}, which the catalogue records as the Eichler commutation relation through
the theta correspondence (S.~Rallis, \emph{Langlands' functoriality and the Weil representation}, Amer.\ J.\
Math.\ 104 (1982); R.~Howe, \emph{$\theta$-series and invariant theory}, Proc.\ Sympos.\ Pure Math.\ 33 (1979);
T.~Tamagawa, Ann.\ of Math.\ 77 (1963)), place not confirmed, relation: a special case (split places of unitary
groups, unramified characters, the unramified vacuum)}

\begin{proposition}[Hecke transfer along the chain]\label{p:hecke}
In Settings \ref{set:good} and~\ref{set:mark}, let $\mathfrak{l}\in\mathcal{L}$ with
$\cj{\mathfrak{l}}\ne\mathfrak{l}$, and let $i$ be a natural number. There are complex numbers $c_n$,
$n\in\NN$, with $c_n\in\Oj$ for all $n$ if $i<m$, such that for every $\mathsf{x}$ in the support of~$\Psi$:
\begin{enumerate}[label=(\alph*)]
\item $\lkom{D}{\vartheta}{i}(1,a)\,\td{D}{\mathsf{x}}{i}=\td{D}{\mathsf{x}}{i}$ for every
$a\in\Kmax{\XD{D}{i+1}}{\mathfrak{l}}$; and
\item for every form $\varphi\in\Stg{D}{0}$ that is fixed by $\Kmax{\XD{D}{0}}{\mathfrak{l}}$, with
$C=\chn{D}{\vartheta}{\td{D}{\mathsf{x}}{\,\cdot\,}}{i+1}$,
\[\sum_{n\le\rr{D}{p}}c_n\,\Tl{\mathfrak{l}}{n}\bigl(C\varphi\bigr)=p^{\DeltaP{D}{p}}\cdot C\bigl(\Tl{\mathfrak{l}}{\rr{D}{p}}\varphi\bigr).\]
\end{enumerate}
\end{proposition}
\begin{proof}
By Lemma~\ref{l:markprop}, $p$ does not divide $\gD{D}\ND{D}$ and $p\in\mathfrak{l}$; by hypothesis $p$ is prime
to $\nrm{\vartheta}$. So $\mathfrak{l}$ and $\cj{\mathfrak{l}}$ are not chosen, do not divide $\piN{D}$ and do not
lie in $Q(\vartheta)$ or its conjugate, and the characters of all the lifts are unramified at~$\mathfrak{l}$.
Since $\mathcal{L}$ is not empty, $p$ is served; by Lemma~\ref{l:served}(b) the order of $\pin{D}{-p}$ is even,
say $2o'$, and the residue field of $\mathfrak{l}$ has $\mathrm{q}=p^{2o'}$ elements (the decomposition of an inert
prime, unit~152). For $\mathsf{x}$ in the support, Lemma~\ref{l:intafter} shows that every datum
$\td{D}{\mathsf{x}}{i'}$ is a vacuum at $\mathfrak{l}$ and~$\cj{\mathfrak{l}}$; an element $a$ as in (a) then
acts on it through a value of an unramified character, which is~$1$. This is (a).

For (b), Input~I9(iv) applies to each of the lifts of~$C$, with $k'=\tD{D}+1$. By induction on the number $e$
of lifts, for the chain $C_e$ of the first $e$ lifts and all complex $U$,
\[\mathrm{H}_{\mathfrak{l}}\bigl(-U\,y_e\bigr)\,C_e\varphi=\prod_{e'<e}\ \prod_{l\le\tD{D}}\bigl(1+U\,\mathrm{q}^{\,l}\,p^{\,o'\nD{D}{i}}\,\alpha_{e'}\bigr)\cdot C_e\bigl(\mathrm{H}_{\mathfrak{l}}(-U\,y_0)\varphi\bigr),\qquad y_e=p^{\,o'(\nD{D}{i+1}-\nD{D}{e})},\]
where $\alpha_{e'}$ is the value at the uniformiser of the character of the lift~$e'$. These values are
integral at~$j$: the character of the lift $e'$ has weights $-\wtD{D}{e'}{\tau}$, so $|j(\alpha_{e'})|_p$ is a
product over $\tau$ of powers of~$p$ whose exponents are $\wtD{D}{e'}{\tau}$ at the places inducing
$\mathfrak{l}$ and $-\wtD{D}{e'}{\tau}$ at those inducing $\cj{\mathfrak{l}}$, up to a common sign, and by
Lemma~\ref{l:thick}(b), with $G(c,w)=-c\,[w=\mathfrak{l}]$, the exponents cancel. Taking $e=i+1$, the product is
a polynomial $\Pi(U)$ with coefficients in $\Oj$ and constant term~$1$. Comparing the coefficients of
$U^{\rr{D}{p}}$ after multiplication by the inverse power series $\Pi(U)^{-1}$, which again has coefficients in
$\Oj$, expresses $C(\Tl{\mathfrak{l}}{\rr{D}{p}}\varphi)$, times a power of $p$ coming from $y_0$ and from
$\mathrm{q}^{\rr{D}{p}(\rr{D}{p}-1)/2}$, as a combination of the $\Tl{\mathfrak{l}}{n}(C\varphi)$,
$n\le\rr{D}{p}$, with coefficients that are elements of $\Oj$ times powers of~$p$; the identity needs
$\rr{D}{p}\le\kD{D}$, which holds by size. Multiplying by $p^{\DeltaP{D}{p}}$ and dividing by the power of $p$ on
the left gives the numbers~$c_n$. For $i<m$ they are integral: the exponents of $p$ that are divided out are at
most $2o'\,\rr{D}{p}\,(\nD{D}{m}+\rr{D}{p})$, and $2o'\le\fD{D}$ because the order of an element of $\PlD{D}$ is at
most $\fD{D}$; this is the definition of $\DeltaP{D}{p}$.
\end{proof}
\leanline{\lean{existshecketransfer}\sep unit 152: \lnm{Parameters.q\_eq\_inert}\sep unit 141: \lnm{Data.hasWeight\_liftChar}}

\begin{proposition}[Hecke translates of the data raise gauges by a controlled amount]\label{p:raise}
In Settings \ref{set:good} and~\ref{set:mark}, let $\mathsf{x}$ be in the support of~$\Psi$, $i$ a natural
number, and $g=(g_{\mathfrak{l}})_{\mathfrak{l}\in\mathcal{L}}$ a family with
$g_{\mathfrak{l}}\in\GL_{\XD{D}{i+1}}(\ED{D}_{\mathfrak{l}})$ such that every $g_{\mathfrak{l}}$ lies in the basic
double coset of some index $n\le\rr{D}{p}$ at~$\mathfrak{l}$. Then
\[\lkom{D}{\vartheta}{i}\bigl(1,\mathrm{at}(g)\bigr)\,\td{D}{\mathsf{x}}{i}\ \in\ \raises{\lk{D}{\vartheta}{i}}{1}{1}{p^{\,\sum_{\mathfrak{l}\in\mathcal{L}}\mexp{D}{\mathfrak{l}}+\liftloss{\mathcal{L}}{i}}}.\]
\end{proposition}
\begin{proof}
By Lemma~\ref{l:markprop}, $p$ is a uniformiser at each $\mathfrak{l}\in\mathcal{L}$, so $g_{\mathfrak{l}}$ is the product of an integral invertible matrix $y_{\mathfrak{l}}$, the diagonal matrix
with entries $p$ at the indices of a set $J_{\mathfrak{l}}$ and $1$ elsewhere, and an integral invertible matrix
$y'_{\mathfrak{l}}$; the set $J_{\mathfrak{l}}$ has at most $\rr{D}{p}$ elements. Thus $\mathrm{at}(g)=y\,y''\,y'$
with $y,y'\in\Kf{\XD{D}{i+1}}$ and a point $y''$ whose image at~$\tau$ is the diagonal matrix with entries
$p^{\,\epsilon_\tau(a)}$,
\[\epsilon_\tau(a)=[\ind{\tau}\in\mathcal{L},\,a\in J_{\ind{\tau}}]-[\cj{\ind{\tau}}\in\mathcal{L},\,a\in J_{\cj{\ind{\tau}}}]\]
(the second term is the component at the conjugate place). By Lemma~\ref{l:tdstd} with $\delta=1$ (recall
$p\nmid\gD{D}$ and $p\ne\ND{D}$), the datum lies in $\Std{1}{D,i}$, and Input~I9(i) reduces the claim to a bound
for the jets of all $\Phi\in\Std{1}{D,i}$.

Let $\mathrm{ch}(\tau,h,a)=1$ if ($\ind{\tau}\in\mathcal{L}$, $h=1$ and $a\notin J_{\ind{\tau}}$) or
($\cj{\ind{\tau}}\in\mathcal{L}$ and $a\in J_{\cj{\ind{\tau}}}$), and $0$ otherwise. Condition~\eqref{e:charge}
holds for $\Xi=\even{\mathcal{L}}$: if $\tau\in\Xi$ and $a\notin J_{\ind{\tau}}$, then $\mathrm{ch}(\tau,1,a)=1$;
if $\tau\in\Xi$ and $a\in J_{\ind{\tau}}$, then the place $\pin{D}{-p^{e}}\cdot\tau$ induces
$\cj{\ind{\tau}}$ (Lemma~\ref{l:ind}(c); $p$ is inert and $e$ is odd), so its charge at~$a$ is~$1$; the case
$h=1$, $\pin{D}{-p^{e}}\cdot\tau\in\Xi$ is symmetric. By Input~I9(ii) and (iii), applied with $c=\cj{\pig{D}}$,
for which $|j(c\,s_1)|_p\le1$ by Lemma~\ref{l:level}(i), every entry $(c',c)$ of the jet matrix of
$\Phi\in\Std{1}{D,i}$ has absolute value at most $p$ to the power (charges of $c'$)
$-\jetexp{p}{i}{\even{\mathcal{L}}}$.

It remains to count the exponent of $p$ in the bound $B$ of Input~I9(i). The determinant factors contribute
$\sum_\tau\last{D}{i+1}{\tau}\cdot\bigl(-\sum_a\epsilon_\tau(a)\bigr)$, and
$\last{D}{i+1}{\tau}=\lfloor\nD{D}{i}/2\rfloor-\wtD{D}{i+1}{\tau}$; by Lemma~\ref{l:thick}(b) with
$G(c,w)=\#J_w\cdot(\lfloor\nD{D}{i}/2\rfloor-c)$ this sum is~$0$. For a box of $c'$ at~$\tau$ with entry~$a$,
the factor $|u_\tau(a)|_p$ and the charge together contribute at most $1$ if $\ind{\tau}\in\mathcal{L}$ and the
box lies in a column of height one; at most $1$ if $\ind{\tau}\in\mathcal{L}$, the box lies in a higher column
and $a\in J_{\ind{\tau}}$; and at most~$0$ in the other cases (a charge $1$ at a place inducing a conjugate of a
marked prime is cancelled by $|u_\tau(a)|_p=p$). The entries of a column of a tableau are distinct, so a
higher column at $\tau\in\even{\mathcal{L}}$ contributes at most $\#J_{\ind{\tau}}\le\rr{D}{p}$. Summing: the
columns of height one give at most $\sum_{\tau\in\even{\mathcal{L}}}\kap{D}{\tau}{1}=\sum_{\mathfrak{l}}\mexp{D}{\mathfrak{l}}$,
the others at most $\rr{D}{p}\sum_{\tau\in\even{\mathcal{L}}}\lc{D}{i}{\tau}$. With the jet exponent this is
$\sum_{\mathfrak{l}}\mexp{D}{\mathfrak{l}}+\liftloss{\mathcal{L}}{i}$.
\end{proof}
\leanline{\lean{memraises}}

\begin{proposition}[moving mark operators into the data]\label{p:move}
In Settings \ref{set:good} and~\ref{set:mark}, let $i$ be a natural number and
$\mathfrak{l}_1,\dots,\mathfrak{l}_n$ a list of distinct elements of~$\mathcal{L}$. There is a $\CC$-linear
endomorphism $t$ of\/ $\Wfin{D}{i}$ such that for every $\mathsf{x}$ in the support of~$\Psi$, with
$\Phi=t\,\td{D}{\mathsf{x}}{i}$:
\begin{enumerate}[label=(\alph*)]
\item $\lkom{D}{\vartheta}{i}(a,a')\Phi=t\,\lkom{D}{\vartheta}{i}(a,a')\td{D}{\mathsf{x}}{i}$ for all
$a\in\Gf{\XD{D}{i}}$ and $a'\in\Kf{\XD{D}{i+1}}$;
\item for every column $\varphi\in\Col{D}$,
\[\lklift{D}{\vartheta}{i}(\Phi)\,\chn{D}{\vartheta}{\td{D}{\mathsf{x}}{\,\cdot\,}}{i}\,\varphi=\chn{D}{\vartheta}{\td{D}{\mathsf{x}}{\,\cdot\,}}{i+1}\,\markop{\mathcal{L}}{\mathfrak{l}_1}\cdots\markop{\mathcal{L}}{\mathfrak{l}_n}\varphi ;\]
\item if $i<m$: for every family $g=(g_{\mathfrak{l}})_{\mathfrak{l}\in\mathcal{L}}$ with $g_{\mathfrak{l}}=1$
for $\mathfrak{l}$ in the list and $g_{\mathfrak{l}}$ in the basic double coset of some index at most
$\rr{D}{p}$ for the other~$\mathfrak{l}$,
\[\lkom{D}{\vartheta}{i}\bigl(1,\mathrm{at}(g)\bigr)\Phi\ \in\ \raises{\lk{D}{\vartheta}{i}}{1}{1}{p^{\,\sum_{c=1}^{n}(\DeltaP{D}{p}-\mexp{D}{\mathfrak{l}_c})+\sum_{\mathfrak{l}\in\mathcal{L}}\mexp{D}{\mathfrak{l}}+\liftloss{\mathcal{L}}{i}}}.\]
\end{enumerate}
\end{proposition}
\begin{proof}
Induction on the list. For the empty list take $t=1$; (a), (b) are trivial and (c) is
Proposition~\ref{p:raise}. Let the claim hold for $\mathfrak{l}_2,\dots,\mathfrak{l}_n$ with the operator~$t'$,
and write $\mathfrak{l}=\mathfrak{l}_1$. Take the numbers $c_n$ of Proposition~\ref{p:hecke} for $\mathfrak{l}$
and~$i$, put $z=p^{\mexp{D}{\mathfrak{l}}}/p^{\DeltaP{D}{p}}$ and
\[t=\Bigl(\sum_{n\le\rr{D}{p}}c_n\,z\,\Tl{\mathfrak{l}}{n}\Bigr)\circ t',\]
the Hecke operators acting on $\Wfin{D}{i}$ through $a'\mapsto\lkom{D}{\vartheta}{i}(1,a')$; by (a) for $t'$ and
Proposition~\ref{p:hecke}(a), $t'\td{D}{\mathsf{x}}{i}$ is fixed by $\Kmax{\XD{D}{i+1}}{\mathfrak{l}}$. (a) The
Hecke operators at $\mathfrak{l}$ commute with $\lkom{D}{\vartheta}{i}(a,a')$ for $a'\in\Kf{\XD{D}{i+1}}$
(Input~I8(ii)). (b) By the equivariance of the lift (Input~I4), Hecke operators on the datum become Hecke
operators on the lifted form. With $C_i$, $C_{i+1}$ the chains of the statement and
$\varphi'=\markop{\mathcal{L}}{\mathfrak{l}_2}\cdots\markop{\mathcal{L}}{\mathfrak{l}_n}\varphi$, which is fixed
by $\Kmax{\XD{D}{0}}{\mathfrak{l}}$ by Lemma~\ref{l:markprop}(c),
\[\lklift{D}{\vartheta}{i}(\Phi)\,C_i\varphi=z\sum_{n}c_n\Tl{\mathfrak{l}}{n}\bigl(C_{i+1}\varphi'\bigr)=z\,p^{\DeltaP{D}{p}}\,C_{i+1}\bigl(\Tl{\mathfrak{l}}{\rr{D}{p}}\varphi'\bigr)=C_{i+1}\,\markop{\mathcal{L}}{\mathfrak{l}}\varphi',\]
by the induction hypothesis, Proposition~\ref{p:hecke}(b) and the definition of the mark operator. (c) Expand
each $\Tl{\mathfrak{l}}{n}$ as a sum over coset representatives; each term is
$\lkom{D}{\vartheta}{i}(1,\mathrm{at}(g'))\,t'\td{D}{\mathsf{x}}{i}$ for a family $g'$ that agrees with $g$
except at~$\mathfrak{l}$, where it is a representative of a basic double coset of index $n\le\rr{D}{p}$. The
induction hypothesis applies to~$g'$, and the scalar $c_nz$ has $|j(c_nz)|_p\le p^{\DeltaP{D}{p}-\mexp{D}{\mathfrak{l}}}$
because $c_n\in\Oj$ for $i<m$; conclude by Input~I4.
\end{proof}
\leanline{\lean{existsmovemarkOperator}}

\begin{proposition}[the chain after the mark operators]\label{p:markpow}
In Settings \ref{set:good} and~\ref{set:mark}, let $Y$ be the product of the operators
$\markop{\mathcal{L}}{\mathfrak{l}}$ over all $\mathfrak{l}\in\mathcal{L}$, in a chosen order, and let
$v\in\Colint{D}$. Then the form $\Bch{D}{\vartheta}{\Psi}\bigl(Y^{m}v\bigr)$ has $j$-gauge at most
$p^{\,\sum_{i<m}(\#\mathcal{L}\cdot\DeltaP{D}{p}+\liftloss{\mathcal{L}}{i})}$ at~$1$.
\end{proposition}
\begin{proof}
For each $i$ take the operator $t_i$ of Proposition~\ref{p:move} for the list of all elements
of~$\mathcal{L}$. By (b) and induction on~$i$, the chain of $i$ lifts with the data $t_{i'}\td{D}{\mathsf{x}}{i'}$
applied to $v$ equals the chain with the data $\td{D}{\mathsf{x}}{i'}$ applied to $Y^{i}v$. By (a) the $t_i$
commute with the levels, and by (c) for the family $g=1$ the datum $t_i\td{D}{\mathsf{x}}{i}$ lies in
$\raises{\lk{D}{\vartheta}{i}}{1}{1}{p^{\,\#\mathcal{L}\DeltaP{D}{p}+\liftloss{\mathcal{L}}{i}}}$. So
Proposition~\ref{p:gchain}, with $d_i=1$, bounds the gauge of the chain combination at $Y^{m}v$. The restriction
$\resY{D}$ does not increase gauges at~$1$ when $p\ne\ND{D}$ (unit~141), and $p\ne\ND{D}$ by
Lemma~\ref{l:markprop}(b).
\end{proof}
\leanline{\lean{gaugeLemarkOperatorpow}\sep unit 141: \lnm{Data.GaugeLe.restrictToY}}

\sect{The image of the mark operators has small index}{s:type}
\statement{inksoft}{\textbf{Input I10 (twisted Hecke operators; unit 144).}\enspace Let $n\ge1$ be a natural number,
$W$ a valuation ring and $R$ a commutative ring. \emph{Twisting data} $T$ for $(n,W,R)$ are a prime number
$p_T$, a natural number $o$, ring homomorphisms $\phi_s\colon W\to R$ for $s\in\Z/n$ (the \emph{letters}), a
parity (even or odd) for each $s\in\Z/n$, and an element of $R$ with square $p_T$, such that $5\le p_T^{\,o}$,
$p_T$ divides $\phi_{s+1}(y)-\phi_s(y)^{p_T^{o}}$ for all $s$ and~$y$, the parities of $s$ and $s+1$
differ, and the residue field of $W$ has $(p_T^{\,o})^{n}$ elements. For natural numbers $q'$, $o$, $n$, $k'$,
$a'$, $r'$, $l'$ and real numbers $c_1$, $c_2$, $c_3$, the \emph{Hecke sizes} are the conditions: $0<r'$,
$2r'\le k'$, $c_1\le1\le c_2$, $2l'=n$, $n\,a'k'(l'+2)+a'+k'<q'^{\,l'}$,
$6(2(k'+k'^2)o\,n^2+2)+2\le c_2$, $2(n+2)k'/c_2\le c_1$,
$(c_2+n(k'+1))\,k'\,q'^{\,3r'}+k'+2\le a'/c_2$ and $n\,a'\,(q'^{-2r'}+r'^2/(k'q'))\le c_3$. The theorem of
unit~144 on the elementary divisors of a Frobenius-twisted Hecke operator says: under the Hecke sizes and if
$p_T$ is not a unit of~$R$, an operator given on the image $M$ of an $R$-linear map from functions on a finite
set of classes by the twisted formula of that unit, with an injective source map, has type
$(p_T,s,(c_3+c_1s+c_2)\cdot n_M)$ on $M$ for every natural number~$s$, where $n_M$ is the number of classes times
the number of Weyl indices.}
\leanline{unit 144: \lnm{TwistingData}, \lnm{HeckeSizes}, \lnm{TwistingData.hasType\_twisted\_hecke}}

\begin{proposition}[the type of a mark operator, from twisting data]\label{p:typez}
In Setting~\ref{set:mark}, let $\mathfrak{l}\in\mathcal{L}$, let $W$ be the valuation ring of $\ED{D}$
at~$\mathfrak{l}$, $n\ge1$ a natural number, and $T$ twisting data for $(n,W,\Oj)$ with prime $p_T$, exponent
$o$ and letters $\phi_s$, such that $p_T$ is not a unit of\/ $\Oj$ and $p_T=p$. Let $t$ be a bijection from
$\Z/n$ onto the set of the places $\tau\in\tops{D}$ with $\ind{\tau}=\mathfrak{l}$ or
$\ind{\tau}=\cj{\mathfrak{l}}$, such that $\phi_s$ is the restriction of $\emb{t(s)}$ to $W$ if $s$ is even
and of $y\mapsto\emb{t(s)}(\cj{y})$ if $s$ is odd; assume
$\mexp{D}{\mathfrak{l}}=\aD{D}\cdot\#\{s\in\Z/n:\ s\text{ even}\}$, and let $l'$ be a natural number and $c_1$,
$c_2$, $c_3$ real numbers such that the Hecke sizes hold for
$(p^{\,o},o,n,\kD{D},\aD{D},\rr{D}{p},l';c_1,c_2,c_3)$. Then for every natural number~$s$ the mark operator
$\markop{\mathcal{L}}{\mathfrak{l}}$ has type $(p,s,(c_3+c_1s+c_2)\cdot\dim_{\CC}\Col{D})$ on $\Colint{D}$.
\end{proposition}
\begin{proof}
This is a translation into the terms of Input~I10. By Lemma~\ref{l:level}(d) and Input~I1, evaluation at
representatives $g_c$ of the classes of $\colL{D}{0}$, chosen integral above~$p$, identifies $\Col{D}$ with the
functions on the classes with values in $\CC^{\Wy{\XD{D}{0}}{\kapD{D}}}$, and $\Colint{D}$ with the functions
with values in $\Oj^{\Wy{\XD{D}{0}}{\kapD{D}}}$ (Input~I1, last property of the gauge); this is the map of
Input~I10, and $\dim\Col{D}$ is the number~$n_M$. The Hecke operator $\Tl{\mathfrak{l}}{\rr{D}{p}}$ on a form fixed
by $\Kmax{\XD{D}{0}}{\mathfrak{l}}$ is the sum, over the cosets $\beta$ of the parahoric subgroup of
$\GL_{\kD{D}}(W)$ of block sizes $(\rr{D}{p},\kD{D}-\rr{D}{p})$, of the translates by
$\beta\cdot\diag(p1_{\rr{D}{p}},1)$. Write $g_c\cdot\mathrm{at}(\beta\,\diag(p1_{\rr{D}{p}},1))=\gamma\,g_{c'}\,a$
with $\gamma$ rational and $a$ in the level; then the translate of the form at $g_c$ is
$\rhow{\kapD{D}}{\lastD{D}{0}}(\gamma)$ applied to its value at~$g_{c'}$. At the places $t(s)$ this matrix is,
since the weight there is a single row of length $\aD{D}$, a symmetric power of the image of $\gamma$ under the
letter $\phi_s$, or of its inverse transpose for odd~$s$ (unitarity); together with the factor
$p^{\mexp{D}{\mathfrak{l}}}=p^{\aD{D}\#\{s\text{ even}\}}$ this is the twisted formula. At the other places
$\tau$, either $\tau$ is not a top and the weight is trivial below the rank $\kD{D}$, or $\tau$ induces neither
$\mathfrak{l}$ nor $\cj{\mathfrak{l}}$ and $\gamma$ is integral and invertible there, so the matrix is integral
and invertible. The source map $(c,\beta)\mapsto c'$ together with the coset is injective because the level is
free. The formal proof checks these points in 35 steps, of which the last fifteen (the identification of the
letter matrices and of the powers of~$p$) were read and not replayed.
\end{proof}
\leanline{\lean{hasTypemarkOperatora}}

\begin{proposition}[the type of a mark operator]\label{p:type}
In Setting~\ref{set:mark}, let $s$ be a natural number and $\mathfrak{l}\in\mathcal{L}$. Then
$\markop{\mathcal{L}}{\mathfrak{l}}$ has type
\[\Bigl(p,\ s,\ \bigl(128\,\mexp{D}{\mathfrak{l}}\cdot\frac{\log b}{m\,p}+\frac{s}{x^{14}}+x^{40}\bigr)\cdot\dim_{\CC}\Col{D}\Bigr)\quad\text{on }\Colint{D}.\]
\end{proposition}
\begin{proof}
$\mathcal{L}$ is not empty, so $p$ is served: $p$ is inert, $o=\oz{D}{p}$ is odd and $o\le\rr{D}{p}$; moreover
$p\ge13$ (unit~152). Let $\tau_0$ be a top inducing $\mathfrak{l}$ (M3) and $v=\pin{D}{-p}^{\,o}$, of order~$n$;
then $o\,n$ is the order of $\pin{D}{-p}$, which $m$ divides (Lemma~\ref{l:served}(b)), so $n$ is even. Put
$t(s)=v^{s}\tau_0$ for $s\in\Z/n$. By Lemmas \ref{l:ind}(d) and~\ref{l:served}(d), $t$ is a bijection onto the
tops that induce $\mathfrak{l}$ or $\cj{\mathfrak{l}}$, and by Lemma~\ref{l:ind}(c) consecutive places induce
conjugate primes; call $s$ even if $t(s)$ induces~$\mathfrak{l}$. The letters $\phi_s$ of
Proposition~\ref{p:typez} map $W$ into $\Oj$, and $\phi_{s+1}\equiv\phi_s^{\,p^{o}}$ modulo~$p$: both
sides are ring homomorphisms from $W$ to $\Oj/p$ that agree on a root of unity of order $\NnD{D}$ and on
$\sqrt{-d}$, by the Frobenius criterion of unit~152. The residue field of $\mathfrak{l}$ has
$p^{\,o\,n}$ elements (unit~152, as $on$ is even), a square root of $p$ lies in~$\Oj$, and
$\mexp{D}{\mathfrak{l}}=\aD{D}\cdot\#\{s\text{ even}\}=\aD{D}\,n/2$. These are twisting data. The Hecke sizes
hold with $l'=n/2$, $c_1=x^{-14}$, $c_2=x^{40}$ and $c_3=128\,\mexp{D}{\mathfrak{l}}\log b/(mp)$: the formal
text derives the three conditions that are not immediate size comparisons from $p\ge9$, from
$o\,n\ge m\ge2\log b$, from $p^{\,o\cdot3\rr{D}{p}}\le b^{1/2}$ (as $p<x_1$) and from
$p^{\,2\rr{D}{p}}\ge m^2$. Proposition~\ref{p:typez} gives the claim.
\end{proof}
\leanline{\lean{hasTypemarkOperator}\sep unit 152: \lnm{Parameters.frob\_ext}, \lnm{Parameters.q\_eq\_inert},
\lnm{Parameters.thirteen\_le\_of\_chiK\_ne\_one}}

\begin{corollary}[the type of the product]\label{c:typepow}
In Setting~\ref{set:mark}, let $Y$ be the product of the operators $\markop{\mathcal{L}}{\mathfrak{l}}$ over all
$\mathfrak{l}\in\mathcal{L}$, in a chosen order. Then $Y^{m}$ has type
$\bigl(p,\ \jc{D}{p},\ 3\,\loss{D}{p}\cdot\dim_{\CC}\Col{D}\bigr)$ on $\Colint{D}$.
\end{corollary}
\begin{proof}
Types add (Input~I2), so by Proposition~\ref{p:type} the operator $Y^{m}$ has type $(p,\jc{D}{p},c\cdot\dim\Col{D})$ with
\[c=m\sum_{\mathfrak{l}\in\mathcal{L}}\Bigl(128\,\mexp{D}{\mathfrak{l}}\frac{\log b}{mp}+\frac{\jc{D}{p}}{x^{14}}+x^{40}\Bigr).\]
Now the sum of the numbers $\mexp{D}{\mathfrak{l}}$ is at most $\aD{D}\,\#\tops{D}\le\aD{D}\fD{D}$, the set
$\mathcal{L}$ has at most $\fD{D}$ elements, $\jc{D}{p}\le2m\fD{D}\aD{D}/p+1$, and
$\fD{D}(\log b)^2\aD{D}/p\le\loss{D}{p}$ (Lemma~\ref{l:served}(c)); by size, $c\le3\loss{D}{p}$.
\end{proof}
\leanline{\lean{hasTypeprodmarkOperatorpow}}

\sect{The level prime}{s:N}
\begin{setting}[the level prime as $p$]\label{set:N}
Setting~\ref{set:D} holds, and Setting~\ref{set:p} holds with $p=\ND{D}$: $j\colon\CC\to\CC_{\ND{D}}$ is a ring
homomorphism; we write $N$ for $\ND{D}$ and $|\cdot|_N$ for the absolute value of\/ $\CC_N$. $\vartheta$ is a twist of~$D$. The set $\mathcal{N}$ of the conjugates $\cj{\ind{\tau}}$,
$\tau\in\PlD{D}$, of the induced primes is a \emph{split level} for $\ND{D}$ in the sense of unit~141 (the
properties of Lemma~\ref{l:pi}(c)), and this fact is an argument of the statements below.
\end{setting}

\statement{inksoft}{\textbf{Input I11 (names at the level prime; unit 141).}\enspace In Setting~\ref{set:N}:
for a function $F$ on matrices over $\prod_{w\in\mathcal{N}}\ED{D}_w$ and a datum $v$, the \emph{name}
$F\otimes v$ is the datum obtained from $v$ by replacing its component at the places of $\mathcal{N}$ and their
conjugates by~$F$; $\mathbf{1}_{\mathrm{lc}}$ is the indicator function of the integral matrices whose last
column is congruent to $(0,\dots,0,1)$ modulo~$N$. For a natural number~$i$, $d_i\in\Gf{\XD{D}{i}}$ is the point
with component $\diag(1,\dots,1,N)$ at the places of $\mathcal{N}\cap\piPl{D}$ and $1$ at the other places of
$\mathcal{N}$; it has the property that $\resY{D}$ maps forms of $j$-gauge at most $c$ at $d_m$ to forms of
$j$-gauge at most $c$ at~$1$. $\mathrm{F}$ is the \emph{Fourier operator} on $\Stg{D}{0}$: a sum, over the
vectors $y$ modulo $N$ with unit last entry, of an additive character of the last entry times the translate
by an integral matrix with last row and last column $y$, followed by the translate by the inverse of
$\diag(1,\dots,1,N)$, all through the embedding of the general linear group at $\mathcal{N}$ conjugated by~$d_0$.
A pair of general linear elements at $\mathcal{N}$ acts on a name through the function, up to the values
$\det\nolimits_{\chi,n}(g)=\prod_{w\in\mathcal{N}}\chi_w(\det g_w)\,|\det g_w|_w^{-n/2}$ of the characters $\chi$ of
the lifts. \textup{(i)} A chain of pairs of lifts whose steps are equivariant up to characters equals a scalar
times the chain of the modified data applied to a power of the Fourier operator; the scalar is a product of
values $\det_{\chi,n}$ at the matrices $\diag(1,\dots,1,N)$, of an integer, and of inverses of such values.
\textup{(ii)} There is a natural number $e'\le4\,m\,\fD{D}\,b\cdot\dim_{\CC}\Col{D}$ such that the module of the
$\varphi\in\Col{D}$ for which $\mathrm{F}^{m/2}\varphi$ is a column of $j$-gauge at most~$1$ at $d_0$ has index
dividing $N^{e'}$ in $\Colint{D}$. \textup{(iii)} The components at $\mathcal{N}$ of the elements of the column
levels, conjugated by the points $d_i$, are integral with last column congruent to $(0,\dots,0,1)$ modulo~$N$.
\textup{(iv)} For a modulus $\mathfrak{m}$ of $\ED{D}$ and an integer $g'$ with $g'\in\mathfrak{m}$, the name of a
function modulo $\mathfrak{m}$ with integral values, applied to a datum $u$, is a finite combination
$\sum_yc_y\mathrm{U}(y)\mathrm{P}u$ of translates of the projection of~$u$, with $\cj{g'}y$ integral and
$\mathrm{N}(g')^{\,\mathrm{rank}}c_y$ integral (unit~150).}
\leanline{unit 141: \lnm{IsSplitLevel}, \lnm{Data.nonInducedPlaces}, \lnm{IsSplitLevel.name},
\lnm{IsSplitLevel.lastColDatum}, \lnm{Data.diagPointN}, \lnm{Data.GaugeLe.restrictToY\_of\_diag},
\lnm{IsSplitLevel.fourierOp}, \lnm{Data.nonInducedRep}, \lnm{detCharAt}, \lnm{IsSplitLevel.compUpTo\_pairs\_eq\_smul},
\lnm{Data.exists\_indexDvd\_fourierOp\_pow}, \lnm{Data.atPlaces\_conj\_mem\_mirabolic}\sep unit 150:
\lnm{isAdmissibleComb\_name}}

\begin{lemma}[character values at the level prime]\label{l:detchar}
In Setting~\ref{set:N}, let $i\le2m$, $n$ and $i'$ be natural numbers and $S$ a set of finite places. Let $\chi$
be the character of the lift $i$ twisted by $\vartheta$, after base change, and $\delta_S$ the matrix over
$\prod_{w\in\mathcal{N}}\ED{D}_w$, indexed by the elements of $\XD{D}{i'}$ other than the last, with component
$\diag(1,\dots,1,N)$ at the places of $S$ and $1$ elsewhere. Then
$|j(\det\nolimits_{\chi,n}(\delta_S))|_N\le\ND{D}^{\,2\fD{D}b}$.
\end{lemma}
\begin{proof}
The norm increases with $S$ (unit~141), so let $S$ be everything. The character has weights
$-\wtD{D}{i}{\tau}$, so $|j(\chi(u))|_N$ is the product over $\tau$ of the absolute values of the components
of~$u$ at the induced primes to the powers $\pm\wtD{D}{i}{\tau}$; for the determinant of $\delta_S$ each
component is $N$ or~$1$. Hence each factor is at most $N^{\wtD{D}{i}{\tau}}\le N^{2b}$, by the bound of unit~152
for $i\le2m$, and there are $\fD{D}$ factors.
\end{proof}
\leanline{\lean{normdetCharAtdiagNle}}

\begin{lemma}[the indicator of the last column raises gauges little]\label{l:lastcol}
In Setting~\ref{set:N}, let $\mathsf{x}$ be a tuple of translations and $i<m$ a natural number. Then
\[\mathbf{1}_{\mathrm{lc}}\otimes\td{D}{\mathsf{x}}{i}\ \in\ \raises{\lk{D}{\vartheta}{i}}{1}{1}{\ND{D}^{\,\fD{D}b}}.\]
\end{lemma}
\begin{proof}
Let $c_0=(1+\sqrt{-d})/2$. The residue field of the prime induced at the unit place has $N$ elements, so $c_0$
is congruent to an integer $n$ modulo it; put $y_0=c_0-n\in\OO_{\KP{P}}$. Then $j(y_0)$ is not a unit, and
$j(\cj{y_0}\,d)$ is a unit: $\cj{y_0}=1-2n-y_0$ and $(2c_0-1)^2=-d$ is not zero modulo~$N$. By Input~I11(iv),
with a modulus whose places are those of $\mathcal{N}$ and $g'=\cj{y_0}$, the name is a combination of
translates $\mathrm{U}(y)\mathrm{P}\td{D}{\mathsf{x}}{i}$ with $y_0y$ integral and coefficients whose
multiples by $\mathrm{N}(y_0)^{\,\mathrm{rank}}$ are integral; each translate lies in $\Std{y_0}{D,i}$ (the
projection of the datum is a generator, and translations by such $y$ preserve the module, unit~150). So
$\mathrm{N}(y_0)^{\,\mathrm{rank}}$ times the name lies in $\Std{y_0}{D,i}$ and, by Input~I8(i), raises gauges by
at most $(|j(\cj{\pig{D}})|_N/|j(\cj{y_0}d)|_N)^{L}\le1$. Dividing by $\mathrm{N}(y_0)^{\,\mathrm{rank}}$ costs
$|\mathrm{N}(y_0)|_N^{-\mathrm{rank}}$; the rank of the space of homomorphisms is
$\nD{D}{i}\nD{D}{i+1}\le\nD{D}{m}^2$ and the norm is taken in a field of degree $2\fD{D}$, and by size the result
is at most $N^{\fD{D}b}$.
\end{proof}
\leanline{\lean{namelastColDatummemraises}\sep unit 150: \lnm{Parameters.finHeisOp\_mem\_standardData}}

\begin{proposition}[the paired chain at ramified columns]\label{p:pairs}
In Setting~\ref{set:N} there is a complex number $z$ with $|j(z)|_N\le\ND{D}^{\,2m\fD{D}b}$ such that for every
family $v=(v_i)_{i\in\NN}$ with $v_i\in\Wfin{D}{i}$ and every $\varphi\in\Colram{D}$
\[\chnn{D}{\vartheta}{\nu\mapsto\pairT{\mathrm{P}^{\piPl{D}}v}{\nu}}{m/2}\,\varphi=z\cdot\chn{D}{\vartheta}{\bigl(\lkom{D}{\vartheta}{i}(d_i,d_{i+1})(\mathbf{1}_{\mathrm{lc}}\otimes v_i)\bigr)_i}{m}\,\mathrm{F}^{m/2}\varphi .\]
\end{proposition}
\begin{proof}
The sum over $y$ in Definition~\ref{d:good}(b) is a sum of names: translating by $-y/N$ in the last column, or
its conjugate in the last row, multiplies the component at $\mathcal{N}$ by a character of the last column,
respectively row, and $\diag(1,\dots,1,N)$ splits off the places of $\piPl{D}$. Each pair of lifts is then
equivariant, up to the values $\det_{\chi,n}$, under the general linear groups at~$\mathcal{N}$; the characters
are unramified at the places above~$N$. Input~I11(i) gives the identity for ramified~$\varphi$ (ramified columns
are ramified in the sense of unit~141, by Lemma~\ref{l:cgenmem}(b)). The scalar $z$ is a product, over the $m/2$
pairs, of two values bounded by Lemma~\ref{l:detchar}, of an integer and of an inverse value, which is bounded in
the same way; so $|j(z)|_N\le(N^{2\fD{D}b})^{2\cdot m/2}$.
\end{proof}
\leanline{\lean{chainpairTranslateseqsmulchain}}

\begin{proposition}[the chain at the level prime]\label{p:atN}
In Setting~\ref{set:N}, let $\Psi$ be a finitely supported complex function on the tuples of translations
of~$D$ such that $\good{\vartheta}{\Psi}$, and let $\varphi\in\Colram{D}$ be such that $\mathrm{F}^{m/2}\varphi$ is
a column of $j$-gauge at most~$1$ at~$d_0$. Then $\Bch{D}{\vartheta}{\Psi}\varphi$ has $j$-gauge at most
$\ND{D}^{\,3m\fD{D}b}$ at~$1$.
\end{proposition}
\begin{proof}
By (G3) and Proposition~\ref{p:pairs} with $v_i=\td{D}{\mathsf{x}}{i}$, the chain combination at $\varphi$ is $z$
times the chain combination with the data
$t_i\td{D}{\mathsf{x}}{i}$, $t_i=\lkom{D}{\vartheta}{i}(d_i,d_{i+1})(\mathbf{1}_{\mathrm{lc}}\otimes\,\cdot\,)$, at
$\mathrm{F}^{m/2}\varphi$. The operators $t_i$ commute with the levels on these data: by Input~I11(iii) the
level elements, conjugated by the $d_i$, preserve the last column modulo~$N$, hence the indicator. By
Lemma~\ref{l:lastcol} and the definition of the set of raising data,
$t_i\td{D}{\mathsf{x}}{i}\in\raises{\lk{D}{\vartheta}{i}}{d_i}{d_{i+1}}{N^{\fD{D}b}}$. Now $N$ is prime to
$\nrm{\vartheta}$ (Input~I3), so Proposition~\ref{p:gchain} applies and gives gauge at most $N^{m\fD{D}b}$
at~$d_m$; multiply by $|j(z)|_N\le N^{2m\fD{D}b}$ and restrict (Input~I11).
\end{proof}
\leanline{\lean{gaugeLeBmuatN}}

\sect{Proof of Theorem B}{s:B}
\begin{lemma}[four gauge bounds for the chain]\label{l:four}
In Setting~\ref{set:good}, let $v\in\Colint{D}$.
\begin{enumerate}[label=(\alph*)]
\item If $p\ne\ND{D}$ and $p\ne d$, then $\Bch{D}{\vartheta}{\Psi}v$ has $j$-gauge at most~$1$ at~$1$.
\item If $\mathcal{L}$ is a full mark and $Y$ is the product of the operators
$\markop{\mathcal{L}}{\mathfrak{l}}$ over all $\mathfrak{l}\in\mathcal{L}$, in a chosen order, then
$\Bch{D}{\vartheta}{\Psi}(Y^{m}v)$ has $j$-gauge at most $p^{\,2\loss{D}{p}-\gain{D}{p}}$ at~$1$.
\item If $p$ is small and not split for $D$, then $\Bch{D}{\vartheta}{\Psi}v$ has $j$-gauge at most $p^{\,c}$
at~$1$, where $c=m\cdot\coll{D}{p}{\emptyset}+\loss{D}{p}-\gain{D}{p}$.
\item If $p$ is small and not split for $D$, then $\Bch{D}{\vartheta}{\Psi}\bigl(p^{\jc{D}{p}}\cdot v\bigr)$ has
$j$-gauge at most $p^{\,\loss{D}{p}-\gain{D}{p}}$ at~$1$.
\end{enumerate}
\end{lemma}
\begin{proof}
(a) By Lemma~\ref{l:tdstd} with $\delta=\pig{D}$ every datum lies in $\Std{\pig{D}}{D,i}$, so by Input~I8(i) it
lies in $\raises{\lk{D}{\vartheta}{i}}{1}{1}{c}$ with
$c=(|j(\cj{\pig{D}})|_p/|j(\cj{\pig{D}}d)|_p)^{L}=|j(d)|_p^{-L}=1$, as $p\ne d$. Proposition~\ref{p:gchain}
with $t_i=1$, $d_i=1$, $r_i=0$ and the restriction (unit~141, $p\ne\ND{D}$) give the claim.

(b) By Proposition~\ref{p:markpow} the exponent is at most
$\sum_{i<m}\bigl(\#\mathcal{L}\DeltaP{D}{p}+\rr{D}{p}\sum_{\tau\in\even{\mathcal{L}}}\lc{D}{i}{\tau}-\jetexp{p}{i}{\even{\mathcal{L}}}\bigr)$.
The place~$1$ is a top, so $\mathcal{L}$ is not empty and $p$ is served, $\chiK{P}(p)=-1$. By
Lemma~\ref{l:lowsum}, Proposition~\ref{p:jet} with $\Xi=\even{\mathcal{L}}$ and Lemma~\ref{l:coll}(d), each
summand is at most
\[\frac{\loss{D}{p}}{m}-\frac{\gain{D}{p}}{m}+\fD{D}(\log b)^2\frac{\aD{D}}{pm}\ \le\ \frac{2\loss{D}{p}-\gain{D}{p}}{m}:\]
the two terms $\#(\tops{D}\cap\even{\mathcal{L}})\,\aD{D}/x_1$ cancel, and the last step is
Lemma~\ref{l:served}(c).

(c) The empty set is marked, by (M4); for it $Y=1$, $\even{\emptyset}=\emptyset$ and
$\liftloss{\emptyset}{i}=-\jetexp{p}{i}{\emptyset}$. By Propositions \ref{p:markpow} and~\ref{p:jet}, the exponent
is at most $\sum_{i<m}\bigl(-\gain{D}{p}/m+\coll{D}{p}{\emptyset}+\fD{D}(\log b)^2\aD{D}/(pm)\bigr)$, and
Lemma~\ref{l:served}(c) applies again.

(d) $\Bch{D}{\vartheta}{\Psi}$ is linear, so the gauge of (c) is multiplied by $p^{-\jc{D}{p}}$. By
Lemma~\ref{l:coll}(b) with $o=1$, $m\,\coll{D}{p}{\emptyset}\le2m\fD{D}\aD{D}/p\le\jc{D}{p}$.
\end{proof}
\leanline{\lean{gaugeLeone}\sep\lean{gaugeLeBmuofisMarked}\sep\lean{gaugeLecollisionSum}\sep\lean{gaugeLepowjetCeilsmul}}

\begin{proof}[Proof of Theorem B]
If a form has $j$-gauge at most $p^{\,c}$ at~$1$ with $c\le-\beta$, then its multiple by $p^{-\beta}$ has
$j$-gauge at most~$1$ at~$1$. Three cases, following the definition of $\err{D}{p}$.

\emph{$p\in\nsp{D}$.} Then $p$ is a prime below $x_1$ with $\chiK{P}(p)\ne1$, $\err{D}{p}=3\loss{D}{p}$,
$\loss{D}{p}\ge1$ by size and Lemma~\ref{l:served}(c), and the hypothesis reads
$\beta<\gain{D}{p}-3\loss{D}{p}+1$; so $2\loss{D}{p}-\gain{D}{p}\le-\beta$. If $p$ is served, take a full mark
$\mathcal{L}$ (Lemma~\ref{l:exmark}), the operator $Y$ of Lemma~\ref{l:four}(b) and
\[S=Y^{m}\,\Colint{D}+p^{\jc{D}{p}}\,\Colint{D}.\]
By Corollary~\ref{c:typepow} there is a natural number $l\le3\loss{D}{p}\dim\Col{D}$ such that $S$ has index
dividing $p^{\,l}$ in $\Colint{D}$. An element of $S$ is $Y^{m}u+p^{\jc{D}{p}}v$ with $u,v\in\Colint{D}$, and by
Lemma~\ref{l:four}(b) and~(d) both parts are mapped to forms of gauge at most $p^{-\beta}$. If $p$ is not
served, take $l=0$ and $S=\Colint{D}$: by Lemma~\ref{l:four}(c) and Lemma~\ref{l:coll}(c) the gauge is at most
$p^{\,2\loss{D}{p}-\gain{D}{p}}$. In this case the columns need not be ramified.

\emph{$p\notin\nsp{D}$ and $p=\ND{D}$.} Then $\err{D}{p}=4m\fD{D}b$ and $\beta<1-4m\fD{D}b$, so
$3m\fD{D}b\le-\beta$, as $m\fD{D}b\ge1$. Take for $S$ the module of Input~I11(ii) and $l=e'$. For a ramified
column $\varphi\in S$, Proposition~\ref{p:atN} gives gauge at most $N^{3m\fD{D}b}$. (Here $\ND{D}$ is prime to
$\nrm{\vartheta}$ for every twist, and Setting~\ref{set:N} holds with the given~$j$.)

\emph{The other primes.} Then $\err{D}{p}=0$, $\beta\le0$, and we take $l=0$, $S=\Colint{D}$. Since
$p\ne\ND{D}$, and $p\ne d$ because $d$ is a prime below $x_1$ with $\chiK{P}(d)=0$ and so belongs to
$\nsp{D}$, Lemma~\ref{l:four}(a) applies.
\end{proof}
\leanline{\lean{existsindexDvdBmudivisible}}

\sect{Sizes: Theorems C and D}{s:size}
\begin{lemma}[the prime sum, prime by prime]\label{l:prime}
In Setting~\ref{set:D}, let $U$ be a natural number with $\lceil x_1\rceil\le U$, $y$ a real number with
$y\le\aD{D}$ and $\tau\in\PlD{D}$. Then
\[\Ask{D}{y}{\tau}=\sum_{p<U,\ p\text{ prime}}2\,\ask{D}{p}{y}{\tau}\cdot\log p .\]
\end{lemma}
\begin{proof}
The weight is $\lamd{d}{1}{\nu}=\Lambda(\nu)(1-\chiK{P}(\nu))$, with $\Lambda$ the von Mangoldt function. The term
of index $\nu\ge\lceil x_1\rceil$ of $\Ask{D}{y}{\tau}$ vanishes, because $y/\nu\le\aD{D}/x_1$ and every nominal
weight is at least $\aD{D}/x_1$; and $\lceil x_1\rceil\le\aD{D}$ by size. So the sum may be taken over all
$\nu\ge1$, and since $\Lambda$ is supported on the prime powers it is
$\sum_p\sum_{e\ge1}\log p\,(1-\chiK{P}(p)^{e})\,(y/p^{e}-\nomW{D}{\pin{D}{-p^{e}}\cdot\tau})_+$. In each of the
cases $p=d$, $\chiK{P}(p)=1$, $\chiK{P}(p)=-1$ one has
$1-\chiK{P}(p)^{e}=2\cdot\tfrac{1-\chiK{P}(p)}{2}\cdot[e\text{ odd or }p=d]$ for $e\ge1$. The primes $p\ge\lceil
x_1\rceil$ contribute~$0$.
\end{proof}
\leanline{\lean{primeSumeqsumlocalAsk}}

\begin{lemma}[row lengths]\label{l:rows}
In Setting~\ref{set:D}, let $i<m$ and $\nu$ be natural numbers and $\tau\in\PlD{D}$. The row of index~$\nu$ at
$\tau$ of the diagram of\/ $\kapD{D}$ for the rank $\nD{D}{i+1}$ has length
\[\sum_{\nu<h<\nD{D}{i+1}}\kap{D}{\tau}{h}\ \le\ \bigl([\nu=0]\,\kap{D}{\tau}{1}+[0<\nu<\kD{D}]\,\wtD{D}{0}{\tau}\bigr)+\bigl([\nu=0]\,(\bone{D}+\btwo{D})+r^{16}\bigr),\]
where $\bone{D}$ and $\btwo{D}$ are the least multiples of $\NnD{D}$ that are at least $\aD{D}/x$ and
$\aD{D}/x_1$.
\end{lemma}
\begin{proof}
Bound each $\kap{D}{\tau}{h}$ with $h\ne1$ by $[h=\kD{D}]\wtD{D}{1}{\tau}+\gap{D}$ (unit~150); here
$\wtD{D}{1}{\tau}=\wtD{D}{0}{\tau}+\gap{D}$, $\wtD{D}{0}{\tau}\le\bone{D}+\btwo{D}+r^{8}$, the height $\kD{D}$ occurs only
for $\nu<\kD{D}$, and $(\nD{D}{m}+2)\gap{D}\le r^{16}$ by size.
\end{proof}
\leanline{\lean{diagramWeightlebyindex}}

\begin{lemma}[two sizes]\label{l:sizes}
In Setting~\ref{set:D}:
\begin{enumerate}[label=(\alph*)]
\item For every $\sigma\in\Aut$, $|\sigma(\cj{\pig{D}})|^2=\gD{D}$.
\item For natural numbers $n$, $n'$, $e$ with $n\le n'+e$,
$\ \jsz{D}{n}\le\jsz{D}{n'}+\bigl(\log(n'+e+1)-\log(\OmP{P}{\Dmu{D}}/\gD{D})\bigr)\cdot e$.
\end{enumerate}
\end{lemma}
\begin{proof}
(a) $\sigma$ maps $\KP{P}$ to itself, by the identity or by conjugation, and $\pig{D}\cj{\pig{D}}=\gD{D}$.
(b) $n!\le(n'+e)!\le n'!\,(n'+e+1)^{e}$, and $-\log(\OmP{P}{\Dmu{D}}/\gD{D})\ge0$ by size, so
$-n\log(\dots)\le-(n'+e)\log(\dots)$.
\end{proof}
\leanline{\lean{normmapstarpiGsq}\sep\lean{jetSizeleadd}}

\begin{lemma}[the number of classes]\label{l:classes}
In Setting~\ref{set:D}, let $I$ be a finite set and $K$ a subgroup of\/ $\Gf{I}$.
\begin{enumerate}[label=(\alph*)]
\item If $G\ge1$ is a natural number, $|I|\le2\,\nD{D}{m}$ and $\Kprin{I}{G}\subseteq K$, then
$\#\Cl{K}\le\exp\bigl(x^{29}(1+\log G)\bigr)$.
\item If $\Kprin{I}{\cj{\pig{D}}\,\piN{D}}\subseteq K$ and $|I|\le\nD{D}{m}$, then $\#\Cl{K}\le\exp(x^{30})$.
\end{enumerate}
\end{lemma}
\begin{proof}
(a) The classes of $K$ map onto those of $\Kf{I}$ with fibres of at most $[\Kf{I}:K]\le[\Kf{I}:\Kprin{I}{G}]$
elements; this index is at most a power of the norm of~$G$, with exponent $|I|^2$ (background library), and the
number of classes of $\Kf{I}$ is at most $(1+2(6N'(\fD{D}|I|)^2)^3)^{4\fD{D}|I|^2}$ for $N'=\NnD{D}\,d$, a number
in the different (unit~149, from the short bases of Input~I7). Take logarithms and compare by size. (b) Apply
(a) with $G=\gD{D}\,\ND{D}$, which $\cj{\pig{D}}\piN{D}$ divides, and $x^{29}(1+\log(\gD{D}\ND{D}))\le x^{30}$
by size.
\end{proof}
\leanline{\lean{cardstdClassesleexpmullog}\sep\lean{cardstdClassesleexp}\sep unit 149: \lnm{Space.card\_classes\_le}\sep
background library: \lnm{relIndex\_principalLevel\_le}}

\statement{inksoft}{\textbf{Input I12 (sizes of one lift and of the restriction; units 143 and 141).}\enspace In
Setting~\ref{set:D}, let $\sigma\in\Aut$. \textup{(i)} For every form $\varphi\in\Stg{D}{m}$ and every
$h\in\Gf{\YD{D}{m}}$,
$\wn{\sigma\circ(\resY{D}\varphi)(h)}\le\exp\bigl(x^{70}-\sum_{\tau}\brexp{D}{\tau}\cdot\log\ND{D}\bigr)\cdot\wn{\sigma\circ\varphi(\bdU{D}(h))}$.
\textup{(ii)} Let $i$ be a natural number, $\Phi=\mathrm{U}(y)\intvac{D,i}$ a translate of the integral vacuum,
$g$, $h$ points and $L\ge1$ a natural number. Then for every vector~$\xi$, with $l(\tau,\nu)$ the length of the
row $\nu$ at $\tau$ for the rank $\nD{D}{i+1}$,
\[\sqrt{\wn{\sigma\circ\lkjm{D}{1}{i}(\lkom{D}{1}{i}(g,h)\Phi)\xi}}\ \le\ \exp\Bigl(c_i+\tfrac12\sum_{\tau}\sum_{\nu<\nD{D}{i}}\bigl(\tfrac{l(\tau,\nu)}{L}+\jsz{D}{l(\tau,\nu)}\bigr)\Bigr)\cdot\sqrt{\wn{\sigma\circ\xi}},\]
where $c_i$ is the constant of the parameters times $(\NnD{D}\,\nD{D}{i}\nD{D}{i+1}+1)^2$, plus
$\fD{D}\,\nD{D}{i}\nD{D}{i+1}\log(2L+2)$.}
\leanline{unit 141: \lnm{Data.weylNormSq\_restrictToY\_le}\sep unit 143:
\lnm{Parameters.norm\_jet\_le\_of\_isAdmissibleComb}, second part, for an admissible combination of one term}

\begin{proof}[Proof of Theorem C]
\emph{The claim without prime sums.} By Lemma~\ref{l:prime} with $U=\lceil x_1\rceil$ and the definitions of
$\gain{D}{p}$ and $\forced{D}{p}$, and because $\gain{D}{p}=0$ when $\chiK{P}(p)=1$,
\[2\sum_{p\in\nsp{D}}\gain{D}{p}\log p=m\Bigl(\sum_{\tau\in\tops{D}}\Ask{D}{\aD{D}}{\tau}+\kD{D}\sum_{\tau}\Ask{D}{\nomW{D}{\tau}}{\tau}\Bigr).\]
These are exactly the prime sums of $\hrow{D}$, so the exponent in Theorem~C is
\[B+m\Bigl(\#\tops{D}\cdot\jsz{D}{\aD{D}}+\kD{D}\sum_{\tau}\jsz{D}{\wtD{D}{0}{\tau}}\Bigr)-\#\tops{D}\cdot\aD{D}\log\ND{D}+\bud{D},\]
a statement about sizes only.

\emph{The restriction.} By Input~I12(i) and Lemma~\ref{l:branch}(a), the restriction costs the factor
$\exp(x^{70}-\sum_\tau(\kap{D}{\tau}{1}-2b)\log\ND{D})=\exp(x^{70}+2b\fD{D}\log\ND{D}-\#\tops{D}\aD{D}\log\ND{D})$.

\emph{The chain.} Apply Lemma~\ref{l:chainsum} to the data $\td{D}{\mathsf{x}}{i}$, jointly invariant by (G2),
with $\mathrm{N}_i(g,\xi)=\sqrt{\wn{\sigma\circ\xi}}$ and with $r_i$ the factor of Input~I12(ii) for
$L=\aD{D}+1$. The value of the chain combination at a point is a sum of at most
$\prod_{i<m}\#\Cl{\colL{D}{i}}\le\exp(m\,x^{30})$ terms (Lemmas \ref{l:classes}(b) and~\ref{l:level}(a)) for each
$\mathsf{x}$, each of norm at most $|\sigma(\Psi(\mathsf{x}))|\cdot\prod_ir_i\cdot\ee^{B/2}$; the norm
$\sqrt{\wn{\sigma\circ\,\cdot\,}}$ satisfies the triangle inequality, and
$\sum_{\mathsf{x}}|\sigma(\Psi(\mathsf{x}))|\le\exp(x^{60})$ by (G1), since $Q(1)=\emptyset$.

\emph{The exponent.} It remains to bound $\prod_ir_i^2$. In stage $i<m$ the rows are bounded by
Lemma~\ref{l:rows}: the row $\nu=0$ at a top has length at most $\aD{D}+e$ with $e=\bone{D}+\btwo{D}+r^{16}$, the rows
$0<\nu<\kD{D}$ have length at most $\wtD{D}{0}{\tau}+r^{16}$, and the other rows have length at most
$r^{16}$ (and $\bone{D}+\btwo{D}+r^{16}$ for $\nu=0$ off the tops). By Lemma~\ref{l:sizes}(b) each
$\jsz{D}{l(\tau,\nu)}$ exceeds $\jsz{D}{\aD{D}}$, $\jsz{D}{\wtD{D}{0}{\tau}}$ or $\jsz{D}{0}=0$ by at most the
excess times $\log(2\aD{D}+1)-\log(\OmP{P}{\Dmu{D}}/\gD{D})$, which is at most~$m$ by size. Summing over
$\nu<\nD{D}{i}$, over $\tau$ and over the $m$ stages gives the main term of the display above. The remaining
terms, namely $c_i$, $l/L\le2$, the excesses, $m\,x^{30}$, $x^{60}$, $x^{70}$ and $2b\fD{D}\log\ND{D}$, are
bounded by $\bud{D}$; the formal text does this by a chain of size comparisons (through $b$ and $\sqrt b$)
that was read and not replayed.
\end{proof}
\leanline{\lean{weylNormSqBmule}}

\begin{proof}[Proof of Theorem D]
Let $p\in U$. If $p\in\nsp{D}$ then $p\le\lfloor x_1\rfloor$, $\log p\le\log x_1$ and
$\err{D}{p}\log p\le3\fD{D}(\log b)^2\aD{D}\log x_1\cdot m^\ast/p$, where $m^\ast=m$ if $p$ is short for $D$ and
$1$ otherwise. If $p$ is short, $p=mc\pm1$ for a natural number $c$ with $1\le c\le x_1$, and
$m/p\le2/c$; if not, $m^\ast/p=1/p$, and put $c=p$. If $p=\ND{D}\notin\nsp{D}$, the term is
$4m\fD{D}b\log\ND{D}$. So every term is bounded by a term attached to a pair $(c,n)$, $1\le c\le\lfloor
x_1\rfloor$, $n\in\{0,1,2,3\}$, namely $3\fD{D}\aD{D}(\log b)^2\log x_1\cdot2/c$, the prime $p$ being the
$n$-th entry of $(\ND{D},mc-1,c,mc+1)$; different primes of $U$ have different pairs. Summing over all pairs,
with $\sum_{c\le x_1}1/c\le1+\log x_1$,
\[\sum_{p\in U}\err{D}{p}\log p\ \le\ 4\cdot6\,\fD{D}\aD{D}(\log b)^2\log x_1\,(1+\log x_1)\ \le\ \fD{D}\aD{D}(\log b)^3,\]
the last step by the condition $(\log x_1)^6\le\log b$ of the data and by size; the term of $\ND{D}$ is bounded
by the term of the pair $(1,0)$, by $m\log\ND{D}\le30m^2\log b$, $m^2\le4x$ and $x\,b\le\aD{D}$.
\end{proof}
\leanline{\lean{sumallowedLossmullogle}}

\sect{A by-product: the determinant period}{s:det}
The first part of the file on the chain serves another unit (the value families of the unit on Selmer groups).
It is independent of the data~$D$.

\begin{definition}[determinant kernel and determinant period]\label{d:det}
Let $E$ be a CM field with a CM type and Weil data, $I$ a finite set, $\alpha$ a continuous character of the
norm-one adeles of $E$ that is trivial on the norm-one elements of~$E$, $c$ a function
$\operatorname{Pl}(E)\to\NN$, $\mathfrak{m}$ a modulus of~$E$ (a non-zero ideal divisible by no prime together
with its conjugate), $\Psi'$ a complex function on the matrices indexed by $I$ over $\OO_E/\mathfrak{m}$, and
$(\psi_w)_w$ a family of characters of the residue fields of the finite places.
\begin{enumerate}[label=(\alph*)]
\item $\detK{I}{\alpha}{c}{\mathfrak{m}}{\Psi'}$ is the function $h\mapsto\theta_{\alpha,\alpha}[\Psi',\prod_\tau\det\nolimits_\tau^{\,c(\tau)}](1,h)$
on $\Gf{I}$: the theta kernel of the pair $(E^{I},E^{I})$ with both characters equal to~$\alpha$, the finite datum
given by $\Psi'$ modulo~$\mathfrak{m}$ and the polynomial $\prod_\tau\det_\tau^{c(\tau)}$ ($\det_\tau$ the determinant
of the matrix of variables at~$\tau$), at the points $(1,h)$.
\item $\detP{I}{\alpha}{c}{\mathfrak{m}}{\psi}$ is the sum over the classes of $\Kf{I}$, weighted by the inverses
of the orders of their stabilisers, of the function
$\detK{I}{\alpha}{c}{\mathfrak{m}}{\chi_{\mathfrak{m},\psi}\circ\det}$, where
$\chi_{\mathfrak{m},\psi}(y)=\prod_{w\mid\mathfrak{m}}\psi_w(y\bmod w)$.
\end{enumerate}
\end{definition}
\leanline{\lean{detKernel}\sep\lean{detPeriod}\sep background library: \lnm{modKernelI}, \lnm{letterMinor},
\lnm{stdClassSum}, \lnm{Modulus.charOfTamePart}}

\begin{lemma}[invariance and expansion]\label{l:det}
In the situation of Definition~\ref{d:det}, assume that $\alpha$ has weight $-c$ at the infinite places and tame
part $(\psi_w)_w$.
\begin{enumerate}[label=(\alph*)]
\item For every $\Psi'$, every $\gamma\in\Gam{I}$ and every $g\in\Gf{I}$:
$\detK{I}{\alpha}{c}{\mathfrak{m}}{\Psi'}(\gamma g)=\detK{I}{\alpha}{c}{\mathfrak{m}}{\Psi'}(g)$.
\item If $\psi_w=1$ for every $w$ that does not divide~$\mathfrak{m}$, then
$\detP{I}{\alpha}{c}{\mathfrak{m}}{\psi}$ is the weighted sum over the classes of $\Kf{I}$ of the function
$g\mapsto\sum_{y}\detK{I}{\alpha}{c}{\mathfrak{m}}{\mathbf{1}_1}(g\tilde y)$, where $y$ runs over the invertible
matrices over $\OO_E/\mathfrak{m}$, $\tilde y\in\Kf{I}$ is a chosen lift of~$y$, and $\mathbf{1}_1$ is the
indicator function of the unit matrix.
\end{enumerate}
\end{lemma}
\begin{proof}
(a) The polynomial is a product of determinants of the stacked matrix of variables, and for such polynomials
the kernel is invariant under the rational points on the left (unit~143), the characters having the weight $-c$.
(b) Expand $\chi_{\mathfrak{m},\psi}\circ\det$, which vanishes on the singular matrices, as
$\sum_y\chi_{\mathfrak{m},\psi}(\det y)\mathbf{1}_y$; right translation by $\tilde y$ moves $\mathbf{1}_1$ to
$\mathbf{1}_y$ and multiplies by $\alpha(\det\tilde y)$, which is $\chi_{\mathfrak{m},\psi}(\det y)$ by the
hypothesis on the tame part (background library).
\end{proof}
\leanline{\lean{detKernelembFinmul}\sep\lean{detPeriodeqstdClassSumsum}\sep unit 143: \lnm{modKernel\_embFin\_mul}\sep
background library: \lnm{modKernel\_mul\_Kf\_right}}

\begin{proposition}[adding a place to the modulus]\label{p:detins}
In the situation of Definition~\ref{d:det}, assume that $\alpha$ has weight $-c$ at the infinite places and tame
part $(\psi_w)_w$, and that $\psi_w=1$ for every $w$ that does not divide~$\mathfrak{m}$. Let $v$ be a finite
place that does not divide~$\mathfrak{m}$ and $\mathfrak{m}'$ a modulus whose places are those of $\mathfrak{m}$
together with~$v$. Then
\[\detP{I}{\alpha}{c}{\mathfrak{m}'}{\psi}=\prod_{l<|I|}\bigl(1-\mathrm{q}_v^{\,l}\,\alpha_v^{-1}\bigr)\cdot\detP{I}{\alpha}{c}{\mathfrak{m}}{\psi},\]
where $\alpha_v$ is the value at a uniformiser of $v$ of the Hecke character of~$\alpha$, normalised by
$\mathrm{q}_v^{\,|I|/2}$.
\end{proposition}
\begin{proof}
$\alpha$ is unramified at~$v$, and $v$ is not conjugate to a place of~$\mathfrak{m}'$. The local computation
behind Input~I9(iv), which the catalogue marks as substantial (the action of the Hecke coset sums on the
unramified vacuum, by counting subspaces of the residue space), shows that the datum for $\mathfrak{m}'$, which
differs from that for $\mathfrak{m}$ by the indicator of the invertible matrices at~$v$, is obtained from it by
the Hecke polynomial at $X=\alpha_v^{-1}$ acting through the Weil representation. Both class sums are class
integrals of functions that are invariant on the left by Lemma~\ref{l:det}(a) and on the right under $\Kf{I}$;
the Hecke polynomial then becomes a sum over coset representatives, each contributing the same integral, and
the count of the cosets is the $q$-binomial theorem
$\sum_{V}(-X)^{\dim V}\mathrm{q}_v^{\binom{\dim V}{2}}=\prod_{l<|I|}(1-\mathrm{q}_v^{\,l}X)$, the sum over the
subspaces $V$ of the residue space.
\end{proof}
\leanline{\lean{detPeriodinsert}\sep background library: \lnm{sum\_heckePolyCoeff\_smul\_pairWeil},
\lnm{q\_binomial\_submodule} (the $q$-binomial theorem; the catalogue: the theorem itself, place not confirmed)}

\sect{Where the results are used}{s:use}
By a search of the sources for the names of the declarations: Theorems A, B, C and D, with the definitions of
$\hrow{D}$, $\nsp{D}$ and $\err{D}{p}$, are used in the file of the endgame that proves Theorem~A of the handout
(short forms). Proposition~\ref{p:ram} is used by the unit on the filling, for the dimension of the ramified
columns in Theorem~B of the handout. Proposition~\ref{p:span} is used by the unit on the exchange, for its crude
rows; Lemma~\ref{l:prime} by the unit on the rows; the determinant period of §\ref{s:det} by the unit on Selmer
groups (its value families). The definitions of the columns, of the levels, of the restriction and of good data
(Definitions \ref{d:pi}, \ref{d:branch}, \ref{d:central}, \ref{d:good}) are used by the units on the seesaw and
on the exchange; the seesaw unit proves that good data exist and that the rows of the handout are of the form
$\Bch{D}{\vartheta}{\Psi}\psi$. I read the statements of these consumers only where this note quotes them.

\sect{Notes on the formal text}{s:notes}
\begin{enumerate}[label=\textup{\arabic*.},leftmargin=1.6em,itemsep=5pt]
\item \emph{Arguments: the text reading and Lean.} Nothing was compiled, so the argument lists of the
declarations are those that a program reads from the text, and in this folder they differ from Lean's in
several places, because section variables are declared again and included or omitted by commands. The
statements above follow the Lean text, as I read it; the points that only Lean can settle are these.
\emph{Entries of the text reading that are not arguments}: a place $w$ dividing $\piN{D}$ in
Definition~\ref{d:central}(d) (the operator takes the data and a set of places; the note writes
$\ramop{D}{\Sigma}$ through a macro of its own with these two arguments), in Lemma~\ref{l:level}(e), (g) and in
Lemma~\ref{l:intafter}(a); a number called $j$ with its primality in Definition~\ref{d:good}(b),
Lemma~\ref{l:chainsum} and Proposition~\ref{p:gchain} (in Lemma~\ref{l:level}(h) the same variable is the prime
that the note calls~$\ell$); the exponent $e$ in Definition~\ref{d:mark}(a) and Lemma~\ref{l:coll}(b); the set
$\mathcal{L}$ in Lemma~\ref{l:exmark}; the prime $p$ and its place in Theorem~C; $\vartheta$ and $\Psi$ in
Theorem~B, which the statement binds; a second record of data in Lemma~\ref{l:intafter}(a), where the formal
statement names only one; the Weil data and other instances of variables that were declared again, in
Definition~\ref{d:short}, Lemma~\ref{l:sub}, Proposition~\ref{p:span} and Definition~\ref{d:central}(a).
\emph{Arguments that the text reading does not show}: that $\Psi$ is good and $p$ is prime to
$\nrm{\vartheta}$ in the four parts of Lemma~\ref{l:four} (a command includes these hypotheses again); that
$\mathcal{L}$ is marked in Proposition~\ref{p:move} (used after the first local definition of the statement);
the prime $p$ and the homomorphism $j$ in Lemma~\ref{l:level}(i), Lemma~\ref{l:ind}(b),
Lemma~\ref{l:markprop}(a), (c) and Lemma~\ref{l:thick}. The note states all of them, through the Settings.
\item \emph{Statements that the line reading cuts.} The formal statements of Propositions \ref{p:hecke}
and~\ref{p:move} contain local definitions, and the program takes their first sign of definition for the end of
the statement; the two propositions were compared with the full Lean text. The same program counts the command
that follows Lemma~\ref{l:lastcol} and the one that follows Theorem~C into these declarations; they make some
definitions opaque and are no part of the proofs.
\item \emph{Hypotheses that follow from others, or are used in one case only.} In Proposition~\ref{p:hecke} the
hypothesis $\cj{\mathfrak{l}}\ne\mathfrak{l}$ follows from (M1). In Theorem~B the columns have to be ramified
only for $p=\ND{D}$. In the definition of $\err{D}{p}$ the case $p=\ND{D}$ is tested after $p\in\nsp{D}$; by
size $\ND{D}\ge x_1$, so the order does not matter. Lemma~\ref{l:chainsum} takes functions $\mathrm{N}_i$ and
numbers $r_i$ for all $i$ and uses those with $i\le m$. In Proposition~\ref{p:typez} the equalities ``$W$ is the
valuation ring at $\mathfrak{l}$'' and $p_T=p$ are hypotheses of the formal statement.
\item \emph{Forms of the statements.} Integrality of a column is the condition ``$j$-gauge at most $1$ at~$1$'';
Theorem~B(ii) is, in the formal text, membership of $p^{-\beta}\Bch{D}{\vartheta}{\Psi}\varphi$ in the module of
such forms on $\YD{D}{m}$. Theorem~B(i) uses the notion of index of unit~149 (Input~I2), which is stated without quotient
modules. In (G3) the numbers of pairs and of lifts
are written with $\lfloor m/2\rfloor-1$. The jet exponent contains the term $e=0$ for $p=d$. The definition of
the restriction contains the proof that it maps forms to forms. Lemma~\ref{l:level}(b) and
Lemma~\ref{l:served}(a) are instances in the formal text: the first is found by Lean whenever a level is needed,
the second feeds the tactic of size comparisons.
\item \emph{What was not replayed.} The dossier marks as of lower confidence the claims made for dense
automatic steps, whose intermediate goals are not visible without Lean. I read these lines; I could not replay
them. They are: the collection of the exponents into $R^{6M}$ in Proposition~\ref{p:span}; the extraction of
the coefficients and the exponent count in Proposition~\ref{p:hecke}; the count of the exponents in
Proposition~\ref{p:raise}; the identification of the letter matrices in Proposition~\ref{p:typez}; which of
the Hecke sizes each line proves in Proposition~\ref{p:type}; the rewriting of the paired chain in
Proposition~\ref{p:pairs}; the size comparisons at the end of the proofs of Theorems C and~D; and the two
computations in Proposition~\ref{p:detins}. For Proposition~\ref{p:jet} the formal text proves the termwise
inequality \eqref{e:term} by one automatic step after a case split; the proof given above is mine and was
checked against the formal definitions, not against the automatic step.
\item \emph{Size comparisons.} Many steps are comparisons of sizes of the data, proved in the formal text by a
tactic from the scale conditions; the dossier does not record the exponents, and the note says ``by size''.
They belong to the note on the admissible data.
\item \emph{What the final theorems use.} The dossier found no member of the unit that is evidently unused;
the count of the members reached from the final theorems by name alone is not meaningful here, and a kernel
export will decide.
\item \emph{References.} For the inputs from the background library the catalogue of the handout gives a
reference only for the Hecke transfer (Input~I9(iv), place not confirmed, a special case) and for the
$q$-binomial theorem (the theorem, place not confirmed); for the others: \emph{reference to be confirmed}. The
headers of the files of this folder name as related work, without places: for the scaled Hecke operators,
Rallis (1982), V.~Lafforgue (2011) and Z.~Liu (2025); for the divisibility of theta lifts, X.~Zhang (2022),
Berger (2009) and Prasanna (2006); for short columns, Gross (1999), Greenberg--Voight (2014) and the geometry of
numbers of Minkowski and Mahler; for the elementary divisors of twisted Hecke operators, Arsovski (2019),
Liu--Truong--Xiao--Zhao (2023) and Fakhruddin--Pilloni (2021). I did not check these.
\end{enumerate}

\subhead{How this note was written.} From the dossier of the unit: the readings of its \UnitMembers{}
declarations and its \UnitProofs{} structured proofs, all of which I read, with the reader's list of doubts. The
program check of the dossier covered \UnitSteps{} leaf steps of \UnitProofs{} proofs. I read the seven Lean
files of the folder completely, statements and proofs, and compared every statement of this note with its Lean
text, argument by argument, with the program's list of arguments beside it. Outside the folder I read the Lean
definitions of the objects of Inputs I2 to~I6 (index and type, the level prime and the twists, the lift kernel,
the data and the chain, induced primes and marked sets, the bookkeeping of unit~150 in its file), the
definitions of gauges, levels, classes, kernels, lifts and Hecke operators at split places in the background
library, and the statements, not the proofs, of the results quoted in Inputs I7 to~I12. The arrangement into
sections, the remark that the prime sums cancel in Theorem~C, the proof of \eqref{e:term} and the presentation
of the proof of Theorem~D are mine; elsewhere the proofs follow the structured proofs, with bookkeeping steps
merged.

\head{Appendix A. Concordance with the Lean artifact}
{\fontsize{8.6}{11.6}\selectfont\renewcommand{\leanfont}{\ttfamily\fontsize{7.7}{10}\selectfont\color{leancol}}
\setlength{\LTleft}{0pt}\setlength{\LTright}{\fill}\setlength{\tabcolsep}{0pt}
\begin{longtable}{@{}p{30mm}p{56mm}p{43mm}r@{}}
\textbf{In this note} & \textbf{Lean declaration} & \textbf{File} & \textbf{Lines}\\\midrule\endhead
\ConcordanceRows
\end{longtable}}

\subhead{Four statements as they are written in the artifact.}
\quoted{spanshortramifiedColumns}
\quoted{existsindexDvdBmudivisible}
\quoted{weylNormSqBmule}
\quoted{sumallowedLossmullogle}
In these texts {\leanfont D.ramifiedColumns} is $\Colram{D}$, {\leanfont IsShort} the notion of
Definition~\ref{d:short}, {\leanfont D.budget} is $\bud{D}$, {\leanfont integralColumns} is $\Colint{D}$,
{\leanfont IndexDvd} the index of Input~I2, {\leanfont Bμ ϑ Ψ} is $\Bch{D}{\vartheta}{\Psi}$,
{\leanfont IsGoodDatum ϑ Ψ} says that $\Psi$ is good for $(D,\vartheta)$, {\leanfont ϑ.levelNorm} is
$\nrm{\vartheta}$, {\leanfont integralFormsAt ⊤ 1} is the module of the forms of gauge at most $1$ at~$1$,
{\leanfont weylNormSq} is the Gram norm, {\leanfont D.heightPerRow} is $\hrow{D}$, {\leanfont D.nonsplitPrimes} is
$\nsp{D}$ and {\leanfont D.allowedLoss} is $\err{D}{p}$. Elsewhere the letters of the formal text were changed
where the write-up reserves them: the stage {\leanfont j} is~$i$ and the homomorphism {\leanfont ιₚ} is~$j$; the element {\leanfont θ} of the
null cone is $\varpi_0$, as in the note on the CM curve; the
exponent {\leanfont s} of the data is $\Dmu{D}$; a place {\leanfont σ} is~$\tau$; a marked set {\leanfont Λ} is
$\mathcal{L}$ and its elements {\leanfont l} are~$\mathfrak{l}$; a prime {\leanfont w} dividing {\leanfont D.π} is
$\mathfrak{p}$; a tuple of translations {\leanfont x} is~$\mathsf{x}$; the set {\leanfont ev} is~$\Xi$;
{\leanfont κ} is $\varkappa$; the number of boxes {\leanfont N} is~$L$; {\leanfont D.ϱ + 1} is $m/2$.

\head{Notation}
\begingroup\fontsize{8.6}{10.9}\selectfont\raggedright
\setlength{\columnsep}{9mm}\setlength{\parskip}{0pt}
\begin{multicols}{2}
\nt{$P=(N_0,d)$, $D$}{parameters; admissible data for $P$ (Setting~\ref{set:D})}
\nt{$x$, $x_1$, $b$, $m$, $r$}{components of $D$: the three scales, the number of lifts, the rank of the rows}
\nt{$\Dmu{D}$, $\ED{D}$, $\NnD{D}$, $\fD{D}$}{the exponent, the CM field, the order of its roots of unity, the
number of its infinite places}
\nt{$\KP{P}$, $\chiK{P}$}{$\Q(\sqrt{-d})$ and its quadratic character}
\nt{$\PlD{D}$, $\tau$, $\tau'$, $\tau_0$, $\tau_e$}{the group of infinite places of $\ED{D}$; places;
$\tau_e=\pin{D}{-p^{e}}\cdot\tau$}
\nt{$\emb{\tau}$, $\pin{D}{\nu}$}{the embedding at $\tau$; the place named by the integer~$\nu$}
\nt{$\kD{D}$, $\tD{D}$, $\aD{D}$, $\nD{D}{i}$}{$r+1-m$; $5r$; the order of the jets; the ranks
$\kD{D}+i(\tD{D}+1)$}
\nt{$\tops{D}$, $\thick{D}$, $\nomW{D}{\tau}$, $\wtD{D}{i}{\tau}$, $\gap{D}$}{tops, thick places, nominal weight,
weights, the gap $2r^8$}
\nt{$\bud{D}$}{$\fD{D}\aD{D}(\log b)^3$}
\nt{$E$, $\operatorname{Pl}(E)$, $\OO_E$, $\widehat{\OO}_E$, $\Af{E}$}{a CM field; its infinite places, integers,
integral adeles, finite adeles}
\nt{$w$, $\cj{w}$, $E_w$, $\mathrm{q}_w$, $v$}{a finite place, its conjugate, the completion, the size of the
residue field; $v$ also a finite place, or a column, or a family of data, where said}
\nt{$p$, $j$, $\CC_p$, $|\cdot|_p$, $\op$, $\Oj$}{a prime (a natural number where said); a ring homomorphism
$\CC\to\CC_p$; the absolute value; the valuation rings (Setting~\ref{set:p})}
\nt{$\ell$, $\iota$}{a prime number; a ring homomorphism $\CC\to\CC_\ell$}
\nt{$\sigma$, $\Aut$}{an automorphism of $\CC$; the group of all of them}
\nt{$I$, $I'$}{finite index sets}
\nt{$f$}{the number of infinite places of a CM field $E$}
\nt{$\varkappa$, $\mathfrak{e}$, $h$, $c$, $c'$}{column counts; determinant exponents; a height; Weyl indices
(tableaux); $c$ also a real or complex constant, a class, a list index, a function $\operatorname{Pl}(E)\to\NN$,
where said}
\nt{$\Gf{I}$, $\Gam{I}$, $\Kf{I}$, $\Kprin{I}{a}$, $\Kmax{I}{w}$}{finite adelic unitary group, its rational
points, the standard level, principal levels, the group at $w$ (Input~I1)}
\nt{$K$, $J$, $\Cl{K}$, $\chi$}{levels; classes; a character of a level, or of a lift}
\nt{$g$, $h$, $k$, $a$, $a'$, $d$, $d'$, $d_i$, $y$, $y'$, $y''$}{points of the groups; $y$ also a matrix over the
adeles, a vector of residues, a real number, or an integer, where said}
\nt{$\Wy{I}{\varkappa}$, $\rhow{\varkappa}{\mathfrak{e}}$, $\xi$, $\xi'$}{Weyl indices, the weight representation,
vectors of the weight}
\nt{$\Alg{I}{\varkappa}{\mathfrak{e}}$, $\varphi$, $\varphi'$, $\varphi_0$, $\psi$, $\mathrm{R}(g)$}{algebraic forms;
forms and columns; right translation}
\nt{$\gauge{\varkappa}{\mathfrak{e}}{g}{\xi}$, $\wn{\xi}$, $G$}{the gauge at $j$; the Gram norm and its matrix}
\nt{$\boxes{\varkappa}{n'}$, $L$, $l(\tau,\nu)$}{numbers of boxes; row lengths}
\nt{$O$, $V$, $M$, $S$, $u$, $z$, $s$, $l$}{Input~I2: a ring, a module, submodules, a matrix, a scalar, an
exponent, a natural number; $S$ also a set of places, $M$ a natural number, $z$ a complex number, $s$ a number
of lifts or an element of $\Z/n$, $V$ a subspace, where said}
\nt{$\XD{D}{i}$, $\YD{D}{m}$, $i$, $i'$}{index sets of the definite spaces; stages}
\nt{$\ND{D}$, $N$, $\piN{D}$, $\gD{D}$, $g_0$, $g_1$, $\pig{D}$}{the level prime (written $N$ where said), an integer
of norm $N$, the product of the two primes of the data, an integer of $\KP{P}$ of norm $\gD{D}$}
\nt{$N'$}{a natural number in the different of a CM field; for $\ED{D}$ it is $\NnD{D}\,d$}
\nt{$\Kch{D}{I}$}{the chosen level}
\nt{$\vartheta$, $Q(\vartheta)$, $\psi_w$, $\nrm{\vartheta}$}{a twist, its set of places, its tame part, the
product of the norms}
\nt{$\Wfin{D}{i}$, $\mathrm{U}(y)$, $\intvac{D,i}$, $\mathrm{P}^{\Sigma}$}{finite data, Heisenberg operators,
the integral vacuum, a projection (Input~I4)}
\nt{$\lk{D}{\vartheta}{i}$, $\lkom{D}{\vartheta}{i}$, $\lkjm{D}{\vartheta}{i}$, $\lklift{D}{\vartheta}{i}$}{the lift
kernel, its representation, its jet matrix, the lift}
\nt{$\raises{k}{d}{d'}{c}$}{data that multiply gauges by at most $c$ from $d$ to $d'$}
\nt{$\mathsf{x}$, $\td{D}{\mathsf{x}}{i}$, $\Phi$, $\Phi'$, $\Phi^{(2)}$, $\Psi$, $\Psi'$}{tuples of translations,
their data; data; a function on tuples; a function on matrices}
\nt{$\chn{D}{\vartheta}{\Phi}{s}$, $\chnn{D}{\vartheta}{\Phi}{s}$, $C$, $C_i$, $\Bch{D}{\vartheta}{\Psi}$}{chains
of lifts and of pairs; the chain followed by the restriction}
\nt{$\Col{D}$, $\Colram{D}$, $\Colint{D}$, $\Stg{D}{i}$}{columns, ramified columns, $j$-integral columns, stages}
\nt{$\piPl{D}$, $\mathfrak{p}$, $\Sigma$, $\colL{D}{i}$}{the primes dividing $\piN{D}$, one of them, a set of
them; the column levels}
\nt{$\lstEq{D}$, $\bdU{D}$, $\brexp{D}{\tau}$, $\Brm{D}$, $\resY{D}$}{Definition~\ref{d:branch}}
\nt{$\short{B}{\varphi}$, $B$, $B'$}{short forms; real bounds}
\nt{$\cgen{D}{\mathfrak{p}}$, $\cenrep{D}$, $\ramop{D}{\Sigma}$, $\nu$}{Definition~\ref{d:central}; $\nu$ an
adele, an integer, or the index of a row or of a pair, where said}
\nt{$\admT{q}$, $q$, $\pairT{v}{\nu}$, $\good{\vartheta}{\Psi}$}{Definition~\ref{d:good}}
\nt{$\ind{\tau}$, $\mathcal{L}$, $\mathfrak{l}$, $\even{\mathcal{L}}$, $\mexp{D}{\mathfrak{l}}$,
$\lc{D}{i}{\tau}$, $\lb{D}{i}{\tau}$}{Input~I5}
\nt{$\ask{D}{p}{y}{\tau}$, $\forced{D}{p}$, $\gain{D}{p}$, $\loss{D}{p}$, $\rr{D}{p}$, $\DeltaP{D}{p}$,
$\jc{D}{p}$, $\oz{D}{p}$, $\coll{D}{p}{\Xi}$, $\Xi$, $e$,
$\cc{p}{\varkappa}{n'}{\Xi}{\tau}{\tau'}{e}$, $\kap{D}{\tau}{h}$, $\last{D}{i}{\tau}$, $\kept{D}{h}$,
$\kapY{D}{\tau}{h}$, $\nsw{d}{p}$}{Input~I6}
\nt{$\jetexp{p}{i}{\Xi}$, $\liftloss{\mathcal{L}}{i}$, $\markop{\mathcal{L}}{\mathfrak{l}}$, $\Tl{w}{s}$, $Y$}{Definition~\ref{d:mark};
Hecke operators; the product of the mark operators}
\nt{$\hrow{D}$, $\nsp{D}$, $\err{D}{p}$, $\jsz{D}{n}$, $\Ask{D}{y}{\tau}$, $\dirc{D}$, $\lamd{d}{1}{\nu}$,
$\OmP{P}{\Dmu{D}}$, $\Lambda$}{Definition~\ref{d:size}; $\Lambda$ the von Mangoldt function}
\nt{$\beta$, $U$, $\eta$, $\delta$, $\delta_S$, $\omega$}{an integer; a finite set of numbers, a natural number or a
complex variable; $\aD{D}/(pm)$; an integer of $\KP{P}$; a diagonal matrix; in $\lkom{D}{\vartheta}{i}$}
\nt{$\intafter{\delta}{S}$, $\Std{\delta}{D,i}$, $\ThProd{P}{\mu}{\delta}{r}{I}$, $s_1$, $\per{P}{\mu}$,
$\Theta$, $\varpi$, $\varpi_0$, $T$}{objects of unit~150 (Inputs I8, I9); $\varpi$ also a uniformiser
(Input~I8); $T$ also twisting data}
\nt{$\mathrm{N}_i$, $r_i$, $t_i$, $t$, $t'$, $\beta_c$, $\mathcal{C}_i$}{sizes, factors, operators on the data, linear
maps, properties (Inputs I8, Lemma~\ref{l:chainsum}); $t$ also a bijection, $r$ a real number}
\nt{$n$, $n'$, $n''$, $o$, $o'$, $k'$, $a'$, $r'$, $l'$, $q'$, $c_0$, $c_1$, $c_2$, $c_3$, $c_n$, $e'$}{natural numbers
and constants defined where they occur}
\nt{$\mathrm{one}_\tau$, $\mathrm{low}_\tau$, $\mathrm{bx}_\tau$; $\mathrm{lv}_e$, $\mathrm{cl}_{\tau,e}$}{the three counts in the coupling count (Input~I6); the indicators of live pairs and of collisions (proof of Proposition~\ref{p:jet})}
\nt{$\mathrm{col}(z)$, $\mathrm{row}(z)$, $\mathrm{at}(g)$, $\mathrm{ev}(y)$}{matrices with one non-zero column or row; the point of a family of local matrices; evaluation vectors}
\nt{$\mathrm{ch}$, $\epsilon_\tau$, $J_{\mathfrak{l}}$, $u_\tau$}{charges; exponents of a diagonal point; sets
of indices; diagonal entries}
\nt{$p_T$, $\phi_s$, $\phi$, $W$, $R$, $X$, $\Pi$, $\alpha$, $\alpha'$, $\alpha_v$, $\alpha_{e}$, $H$}{twisting
data and Hecke polynomials: a prime, the letters, ring homomorphisms, a valuation ring (also a subspace), a ring
or a bound, a variable, a polynomial, values of characters; $\alpha$ also a character}
\nt{$\mathcal{N}$, $\mathbf{1}_{\mathrm{lc}}$, $F$, $\mathrm{F}$, $\mathfrak{m}$, $\mathfrak{m}'$, $g'$, $y_0$}{Setting~\ref{set:N} and Input~I11:
the split level, an indicator, a function, the Fourier operator, moduli, integers}
\nt{$\bone{D}$, $\btwo{D}$, $m^\ast$, $\gamma$, $A$, $A_\phi$, $Y(g)$, $c_{\tau,\nu}$}{multiples of $\NnD{D}$; a factor of the
loss; a rational point; matrices of the proof of Proposition~\ref{p:span}}
\nt{$\detK{I}{\alpha}{c}{\mathfrak{m}}{\Psi'}$, $\detP{I}{\alpha}{c}{\mathfrak{m}}{\psi}$,
$\chi_{\mathfrak{m},\psi}$, $\tilde y$}{Definition~\ref{d:det}}
\nt{$\ee$, $\pi$}{the base of the natural logarithm; in $\piN{D}$, $\pig{D}$}
\end{multicols}
\endgroup

\end{document}
